\subsubsection{Normalized points on the fitted tower}
\label{par:rank-four-promodular}
\hypertarget{source-ordinary-promodularity}{}
The orbit count gives Galois representations, but a Hecke trace count requires ordinary eigenpoints at the prescribed levels. Start with a normalized point and retain its four ordered ordinary characters and local flags throughout the comparison.
The same argument applies when the matrix algebra is full and
when the point has a full two-dimensional block in the positions
\(\{0,2\}\) or \(\{1,3\}\). We describe the passage from a
residual chain to membership in the modular locus. This part
will be used in both host cases.

\begin{paramdefinition}[Coefficient rings and flags on the fitted tower]\label{hyp:fitted-ordinary-tower}
Restrict to a totally real abelian field \(F\) of the tower and
let \(K\) be its fixed imaginary quadratic field. Put
\[
  L=FK,\qquad d=[F:\Q],\qquad
  g=\#\{v\mid\ell\text{ in }F\},\qquad d_v=d/g .
\]
The local degrees are equal since \(F/\Q\) is abelian. The labels of the
four residual characters give the diagonal order at the
distinguished element. Choose the integral basis and the
lower unipotent local flags before constructing the residual
chain. The normalization retains both the conjugacy class
over the fraction field and the marked inertia characters.

The small coefficient model is an injective map from a
complete local integral algebra to the valuation ring of the
point. Writing \(A_{\rm sm}\) for this algebra and \(O_y\) for the target valuation ring, the model contains the finite lists
\[
  A_{\rm sm}\hookrightarrow O_y,\qquad
  \{r(h_a)_{ij}\}_{a,i,j},\quad\{(b_v)_{ij}:i>j,\ v\mid\ell\},
\]
where \(h_a\) are the chosen group generators and \(b_v\) the lower
unipotent matrices of the ordinary flags. The seed uses an ordinary
integral coefficient ring. Use a common finite coefficient DVR \(\Lambda\) for the
coefficient constants: take the compositum of the seed and ordinary
coefficient fields and the unramified extension realizing the
finite chain residue fields. The small integral target model is
a coefficient algebra, with its given map to the target valuation
ring. Its deformation variables are retained in that algebra.
The integral marking at the target and its transported flags are
constructed in Lemma~\ref{lem:rational-rank-four-marking} and
Proposition~\ref{prop:rational-rank-four-torus}.
\end{paramdefinition}

Let \(B\) be the complete local base for the coefficients of
characteristic polynomials and the ordinary characters, and
let \(A\) be the finite type \(B\)-algebra carrying the
normalized matrix representation and flags. Denote the four
inertial characters at \(v\) by \(\psi_{v,i}\).
The ordinary construction gives an ideal \(J\subset B\).
The two properties proved below are
\[
 \mathfrak m_B^n\subset(J,\varpi_\Lambda,\psi_{v,i}(s)-1),\qquad
 JA\subseteq\mathfrak p
\]
for some \(n\) in the first statement and for the primes
\(\mathfrak p\) at the specified points in the second:

\begin{enumerate}
\item A power of \(\mathfrak m_B\) lies in the ideal generated
by \(J\), the coefficient uniformizer and all the
\(\psi_{v,i}(s)-1\) for \(s\) in inertia.
\item At the characteristic-zero automorphic seed, the
extension of \(J\) to \(A\) lies in every prime contained in
the kernel of evaluation. The analogous local vanishing holds
at a characteristic-\(\ell\) formal arc centered over
\(\mathfrak m_B\) which kills \(J\), has absolutely irreducible
generic representation and separates the four inertial
characters on the prescribed open inertia groups.
\end{enumerate}
The local vanishing clause concerns the components through the
point. Vanishing of the evaluations of \(J\) alone would only
put the point in its closed locus. The stronger clause is what
allows propagation to adjacent components.
\label{cor:ordinary-promodularity-1}

\subsubsection{Ordinary ideals and the closed specialization}
\label{par:count-two-ideals}

Membership in ordinary support is expressed by ideals from the two levels. At the closed point one first needs a bound on what remains after imposing these ideals, the marked characters and the coefficient uniformizer.

\begin{proposition}[An ideal primary to the closed tower point]\label{prop:tower-closed-specialization}
Take the two ordinary determinant links on \(G_L\) with the common
scalar twist, and let \(J\) be the sum of the extended ordinary
ideals in \(B\). For some \(d\ge1\), with \(\varpi\) the coefficient uniformizer, 
\[
 \mathfrak m_B^d\subseteq J+(\varpi)+(\psi_{v,i}(s)-1:s\in I_v,\ (v,i)\text{ marked}).
\]

\end{proposition}

\begin{proof}
\proofstep{Use the same tower and twisting character.}

The two auxiliary levels on \(L\) use distinct split primes
and a common scalar twist \(\kappa\). Each adds only one
possible exceptional prime to the split-level set. At each
of these levels take the universal ordinary determinant ring
for the twisted residual representation, and its ordinary
modular quotient with kernel \(J_i\). Pull the corresponding
closed loci back to \(\operatorname{Spec} B\) by the
determinant and ordered ordinary characters of the universal
tower representation. Writing the link maps as \(f_i:R_i\to B\) gives
\[
                   J=f_1(J_1)B+f_2(J_2)B .
\]
The extension is by these two maps for the same tower point. The change of
similitude on \(L\) supplied by \(\kappa\) gives the
polarization for these determinant rings. Both comparisons
use the same twist on the same tower.

Write \(R_i\) for one universal ordinary determinant ring
and \(f_i:R_i\to B\) for its map. If the characteristic polynomial of \(h\in G_L\) is \(P_h(X)\), twisting all four roots gives
\[
 P_h^{\rm tw}(X)=\kappa(h)^4 P_h(X/\kappa(h)).
\]
At the fixed marked place this same twist sends each ordered character \(\psi_a\) to \(\kappa\psi_a\). The determinant law descends to \(B\): its values on
sums of group elements are integral polynomials in the
characteristic coefficients of words in those elements.
The polarization and flag relations hold over \(A\), hence
in \(B\) by injectivity. Use the empty-auxiliary ordinary
module \(M_i\) at the level of \(R_i\), with the tame type,
weight character and residual localization of
Section~\ref{sec:count-two-level-data}.

In the paragraphs
below, \(J_i=\ker(R_i\to\operatorname{End}_{\Lambda_i}M_i)\)
throughout, where \(\Lambda_i\) is the ordinary weight
algebra. The determinant and ordered characters of the
constituent attachments in
Section~\ref{sec:count-two-level-attachments} give that action.
The quotient of the link ring embeds into the endomorphisms
of the empty-auxiliary module:
\[
    R_i/J_i\ \lhook\joinrel\longrightarrow
                         \operatorname{End}_{\Lambda_i}(M_i).
\]
The target is finite free over \(\Lambda_i\), so the image is
finite and torsion-free.
The enlarged levels in the patching argument have the
same base module when the auxiliary conditions are removed.
\proofstep{Make the determinant coefficients integral over the residue field.}

Write \(I\) for the ideal generated in \(B\) by
\(f_i(J_i)\), the coefficient uniformizer and
\(\psi_{v,a}(s)-1\) for all marked inertia elements \(s\).
Take a prime \(\mathfrak q\supset I\) and set
\(D'=B/\mathfrak q\), a local domain with the same
residue field \(k\) as the coefficient ring.
The weight map to \(D'\) sends the ordered characters on
inertia to the reductions of the fixed twisting character
\(\kappa\). The universal property of the weight algebra
therefore identifies this map with its residue specialization.
The map thus factors on the quotient through
\[
  (R_i/J_i)\otimes_{\Lambda_i}k\longrightarrow D' .
\]
The algebra on the left is finite-dimensional over \(k\), so the image of each element of \(R_i\) is integral over \(k\).

We must recover the generators on the tower side. For \(h\in G_L\), if \(a_j(r(h))\) is the coefficient of
\(T^{4-j}\), then
\[
                 a_j(\kappa(h)r(h))=\kappa(h)^j a_j(r(h)).
\]
The left side belongs to the image of \(f_i\).
As \(\kappa(h)\) is a coefficient unit, \(a_j(r(h))\) is
integral in \(D'\) over \(k\). For \(g\in G_F\) use
\(g^2\in G_L\).
Each coefficient of the characteristic polynomial at \(g\)
is integral over the ring generated by the coefficients at
\(g^2\).

To see the universal identity, use indeterminates
\(x_1,\ldots,x_4\). Each \(x_a\) satisfies the monic
equation \(\prod_b(T^2-x_b^2)=0\) at \(T=x_a\).
Consequently each elementary symmetric function lies in the
integral closure of
\[
        \mathbf Z[e_b(x_1^2,\ldots,x_4^2):1\le b\le4]
            \ \subset\ \mathbf Z[x_1,\ldots,x_4].
\]
The monic relations specialize to characteristic coefficients
of matrices, including the universal matrix over \(A\), by
the symmetric polynomial identities for the square of a
matrix. The inclusion \(B\hookrightarrow A\) descends those
relations to \(B\). Reduction modulo \(\mathfrak q\) and
transitivity of integrality give the assertion for all \(g\).
For the other generators, every decomposition group at a
marked \(v\mid\ell\) lifts to \(G_L\) since \(v\) splits in
\(L/F\). Its four ordered character values therefore come
from the ordinary determinant link, again up to a coefficient unit.

\proofstep{Recover the closed point from its generators.}

Integral dependence of the generators must now imply equality with their residues.


An element \(x\in D'\) integral over \(k\) equals its
residue \(a\in k\) inside \(D'\). The difference \(x-a\) lies in the maximal ideal, whereas \(k[x-a]\) is a finite-dimensional domain over \(k\), hence a field. If \(x-a\ne0\), then
\[
 1=(x-a)(x-a)^{-1}\in\mathfrak m_{D'},
\]
since the inverse lies in \(k[x-a]\subset D'\). Thus the
quotient map and the residue map into \(D'\) agree on the
characteristic coefficients and the marked local characters.
These determine maps out of the complete local base \(B\).
Hence \(\sqrt I=\mathfrak m_B\). By Noetherianity
\(\mathfrak m_B^n\subset I\) for some \(n\). Since
\(f_i(J_i)\subset J\), one link proves the first property.

\end{proof}

\subsubsection{Base and enlarged ordinary modules}
\label{par:ordinary-functions}

To control these ideals on nearby components, the two ordinary levels must admit compatible integral modules at auxiliary level. Reduction by the augmentation ideal should recover the specified base module.

\begin{proposition}[Ordinary modules for the auxiliary system]\label{prop:ordinary-level-modules}
For the stated tame coefficient, residual localization and weight
character, the enlarged ordinary modules are finite projective over
their auxiliary group algebras. Denote one such module by \(M_N\), its auxiliary group by \(\Delta_N\), the ordinary weight ring by \(\Lambda\) and the module at the auxiliary Iwahori level by \(M_0\). Base change along the augmentation gives 
\[
 M_N\otimes_{\Lambda[\Delta_N]}\Lambda\simeq M_0
\]
 for the associated Hecke quotient. The
classical evaluations considered below separate the characteristic-zero
Hecke action on these modules.
\end{proposition}

\begin{proof}
\proofstep{Integral coefficients and the fixed residual localization.}

Use the definite unitary group for the standard Hermitian
form over the CM field \(L\). Keep the tame coefficient
type, the split place \(v_0\) with principal congruence
level, and the ordinary \(\ell\)-adic data of the initial
level. The auxiliary product of groups \(\mathrm{GL}_n(L_w)\)
acts by right translation at the chosen places over
\(L^+\). The level at these places is initially
hyperspecial. Their later replacement by compact open
subgroups makes no change to the weight character or
to the tame coefficient type. Write \(\Lambda\) for the
weight algebra before adjoining the formal variable
\(s\).

At a selected place \(v\mid\ell\), local reciprocity identifies
the pro-\(\ell\) inertia units modulo torsion with
\[
                         U_v\simeq\Z_\ell^{d_v}.
\]
Here \(d_v\) is the local degree. Fixed residual and torsion
values leave the coordinate group \(\prod_{v\mid\ell}U_v^n\).
Using a basis in each \(U_v\), the completed group algebra is
\[
 \Lambda=O\left[\!\left[\prod_{v\mid\ell}U_v^n\right]\!\right]
       \simeq O[[w_1,\ldots,w_{n\sum_{v\mid\ell}d_v}]],
       \qquad \sum_{v\mid\ell}d_v=[L^+:\Q].
\]
For a basis element \(u_{v,a}\) of \(U_v\), the
corresponding coordinates are \([u_{v,a}]_i-1\), one for
each diagonal slot \(i\).
Use the induced character \(\theta\) of the product of
the compact diagonal tori. The inverse and cyclotomic
shifts converting Galois inertia to this torus character
are fixed twists, so the generators of the same free
lattice also give coordinates for \(\theta\).

The ordinary function space can be defined over any
complete Noetherian local coefficient ring \(A\) with
finite residue field carrying these data. Its functions
on the finite adeles are invariant on the left by the
rational unitary group, transform by the coefficient
type at the fixed tame level, and are invariant by the
upper unipotent subgroup \(N\) of the integral groups
at the \(\ell\)-adic places.

For \(f\) in the function space the local conditions on
the right are
\[
 f(xn)=f(x),\qquad f(xt)=\theta(t)^{-1}f(x),\qquad
 f(x\bar n)\equiv f(x)\pmod{\mathfrak m_A^j}
          \quad(\bar n\in\bar N_{d(j)}).
\]
Here \(n\in N\), \(t\in T\), and \(\bar N_d\) has off-diagonal
entries in the \(d\)-th power of the local uniformizer ideal. For
\(d\ge1\) put
\[
 J_d=N T\bar N_d.
\]
These conditions
are closed under uniform limits for the filtration
\(V_j\) by functions with values in \(\mathfrak m_A^j\).
Indeed a limit modulo \(\mathfrak m_A^j\) equals one
eventual member of its sequence, which has the required
lower unipotent invariance to that precision.

The stabilizers at the tame level and any of the
\(J_d\) are trivial. A rational stabilizer preserves
the definite Hermitian form at infinity and is in a
conjugate of a compact subgroup at every finite place.
This bounds its positive and negative powers locally,
so all characteristic roots are units at every finite
place. They have absolute value one at every
archimedean embedding. Kronecker's theorem makes the characteristic roots \(\alpha_i\) roots of unity. At an archimedean embedding the matrix \(a\) is diagonalizable, since it is unitary. For \(h=\operatorname{lcm}_i\operatorname{ord}(\alpha_i)\) this gives
\[
 a^h\sim\operatorname{diag}(\alpha_1^h,\ldots,\alpha_n^h)=1 .
\]
The embedding is faithful, so the rational element also satisfies \(a^h=1\).

At \(v_0\) the stabilizer is conjugate into
\(1+p_0M_n(\mathcal O_{L_w})\), with \(p_0\) odd and
unramified in \(L\).
The matrix logarithm is injective on
\(1+p_0 M_n(\mathcal O_{L_w})\), since the exponential
and logarithm converge there and are mutual inverses.
For a matrix \(a\) of finite order \(h\) in this congruence
subgroup,
\[
                    h\log a=\log(a^h)=0.
\]
The matrix logarithm is injective and the field has characteristic
zero, so \(a=1\). The argument
continues to apply when auxiliary places move and
when their compact levels change.

Let \(z\) at the \(\ell\)-adic places be the
strictly dominant diagonal translation, contracting
\(N\) by conjugation and contracting \(\bar N_d\) by
inverse conjugation. The ordinary operator \(U_z\)
is the sum of translations by \(b z\) over
\(b\in N/zNz^{-1}\) at these places. For each
\(j\), take \(c\ge1\) such that \(\theta\) on the diagonal
congruence subgroup \(T_c\) is trivial modulo
\(\mathfrak m_A^j\). The operator lowers the depth of
lower unipotent invariance modulo this ideal from
\(d+1\) to \(d\) whenever \(d\ge c\).

In fact set \(J'_d=N T_c\bar N_d\).
The Iwahori factorization gives
\[
 J'_d z J'_{d+1}
   =\coprod_{b\in N/zNz^{-1}} b z J'_{d+1}.
\]
Here \(z^{-1}\bar N_d z\subset\bar N_{d+1}\),
the torus commutes with \(z\), and the upper
unipotent intersection is \(zNz^{-1}\). Thus
the sum is invariant by \(J'_d\) on
\(J'_{d+1}\)-invariants. The compact torus normalizes
\(N\) and commutes with \(z\), preserving the
\(\theta^{-1}\)-equivariance throughout. This proves
the lowering assertion. Iterating the lowering assertion for a function of depth \(d\) modulo \(V_j\) gives
\[
 \operatorname{depth}(U_z^a f\bmod V_j)\le\max(c,d-a)
 \qquad (a\ge0).
\]
Thus the relevant threshold for the iterate index \(a\) is
\[
 a\ge\max(0,d-c)\quad\Longrightarrow\quad \text{depth of the iterate}\le c.
\]


That space has finitely many value coordinates, one
vector of coefficients for each double coset of the
tame level and \(J_c\). For any list of values there
is an \(A\)-linear lift to the original function space.
Choose representatives. In their right cosets use
the tame type and the diagonal part of the
\(\bar N_c T N\) factorization to extend the values.
Trivial stabilizers give well-defined values.
Right multiplication by \(N\) keeps the diagonal
part, and multiplication by \(T\) gives exactly
its character. Right multiplication by \(\bar N_c\)
changes the diagonal coordinates only modulo the
diagonal congruence subgroup \(T_c\), by triangular
elimination. This gives the specified reduction and
the required continuity at smaller lower unipotent
subgroups. Let \(b_j\) be evaluation on these cosets and \(a_j\)
the lifting just described. Evaluation recovers the lifted coordinates at this precision,
\[
 b_ja_j\equiv\mathrm{id}\pmod{\mathfrak m_A^j}.
\]
This gives the finite depth-\(c\) coordinate module.

Every orbit is eventually periodic modulo \(V_j\) because its
depth-\(c\) coordinate module is finite. For a periodic orbit
of period \(a_j\), each \(h!\) with \(h\ge a_j\) is divisible by
\(a_j\). Thus
\[
                    e_{\rm ord}=\lim_{h\to\infty}U_z^{h!}
\]
exists and projects onto the functions periodic at every precision. On those functions
\(U_z\) acts invertibly, by the inverse on each
periodic orbit and passage to the limit. The
projector preserves the filtration and commutes with
the away-from-\(\ell\) action. Put \(M=e_{\rm ord}S_A\),
where \(S_A\) denotes the original function space.
For \(f\in M\), periodicity and eventual membership
in the depth-\(c\) space give
\[
   f-e_{\rm ord}a_j b_j(f)\in V_j .
\]
Here \(b_j\) takes a function in \(V_i\) to values
in \(\mathfrak m_A^i\).

If \(f\in M\cap V_i\), the
approximation at \(j=i+1\) writes it modulo \(V_{i+1}\)
as a linear combination of the columns of
\(e_{\rm ord}a_{i+1}\), with coefficients in
\(\mathfrak m_A^i\). Replace these columns modulo
\(V_1\) by combinations of the columns of
\(e_{\rm ord}a_1\). The error is in \(V_{i+1}\).
Successive approximation and completeness now give
\[
   M\cap V_i=\mathfrak m_A^i M.
\]
Lift a basis of \(M/(M\cap V_1)\) and evaluate it in the
finite depth coordinates. The evaluation columns have full rank
modulo \(\mathfrak m_A\), so a maximal minor is a unit.
Successively using the value filtration equality now gives
\[
       A^{\dim_k M/(M\cap V_1)}\xrightarrow{\;\sim\;} M
\]
from the chosen lifts. This
proves finite freeness over \(A\).

For a normal auxiliary sublevel with quotient
\(\Delta\), the double cosets with full
\(\ell\)-Iwahori and the smaller auxiliary level
form a free \(\Delta\)-set. This follows from the
same stabilizer calculation and disjointness of
the auxiliary and \(\ell\)-adic factors. Choose one
point in each \(\Delta\)-orbit. The indicator function
of the resulting subset pulls back to the adeles,
where multiplication by it preserves the function
space. Project after multiplying, to obtain an
\(A\)-linear endomorphism \(B\) of \(M\). The translates
of the chosen subset partition the double cosets, so
\(\sum_{\delta\in\Delta}\delta B\delta^{-1}=1\).
The action map \(A[\Delta]\otimes_A M\to M\) has the
\(A[\Delta]\)-linear section
\[
 f\longmapsto\sum_{\delta\in\Delta}
                \delta\otimes B(\delta^{-1}f).
\]
Thus \(M\) is projective over \(A[\Delta]\).
Global localization at a residual Hecke eigensystem
takes a direct summand.

In fact a finite commutative
algebra of commuting Hecke operators over a complete
local ring decomposes by the maximal ideals of its
residue algebra. The residual localization idempotent \(e_\mathfrak m\)
lifts and commutes with \(\Delta\). It exhibits
\[
       M=e_\mathfrak m M\oplus(1-e_\mathfrak m)M
\]
as auxiliary group-algebra modules. The same
argument selects the further residual torus tuple
on the module at an auxiliary cover. Over the
local group ring used below, the resulting finite
projective modules are free.

\proofstep{Separation by classical weights and control at the selected type.}

The auxiliary module is an integral object varying with weight. Classical specializations must distinguish its elements if their tame types and operators are to determine the same ordinary module.


For the classical evaluations in this construction
it suffices to use algebraic weight zero. At a smooth
specialization of the compact torus weight character,
the depth \(c\) killing its finite part can be fixed
at every precision over the coefficient DVR
of the specialization. Write \(\mathcal O_x\) for the specialization DVR. For \(u\in\bar N_c\) lowering at each precision gives
\[
 f(gu)-f(g)\in\bigcap_{j\ge1}\mathfrak m_{\mathcal O_x}^{,j}=0 .
\]
Hence each specialized ordinary function is invariant by \(\bar N_c\). It is a classical automorphic
form of weight zero with the fixed tame type and
the specified compact torus character at \(\ell\).
The specialization of the integral ordinary module
identifies with this ordinary space over the new
DVR. Let \(\mathcal O_x\) be the specialization DVR with
residue field \(k_x\), and let \(V_c(k_x)\) be the finite space
of depth-\(c\) functions of its fixed type. Its ordinary
reduction is
\[
 V_c(k_x)^{\rm ord}=\bigcap_{a\ge0}\operatorname{im}
                       (U_z^a:V_c(k_x)\to V_c(k_x)).
\]

With \(V=V_c(k_x)\), the chain of images stabilizes as soon as two successive dimensions agree. For that index \(a_0\),
\[
 U_z^{a_0}V=U_z^{a_0+1}V=\bigcap_{a\ge0}U_z^a V .
\] Lift any fixed set of depth-\(c\) values to the
DVR and apply the integral projector to obtain every element
of that stable image. Conversely periodicity puts the reduction
of any ordinary function in the stable image. The equality of
the value and coefficient filtrations identifies both reduced
modules with \(V_c(k_x)^{\rm ord}\). The specialization map between the two finite free
\(\mathcal O_x\)-modules reduces to the identity of
\(V_c(k_x)^{\rm ord}\). In lifted bases let \(S\) be its matrix. The identity on the common reduction gives
\[
 \det S\equiv1\pmod{\mathfrak m_{\mathcal O_x}},
 \qquad S^{-1}=(\det S)^{-1}\operatorname{adj}(S)\in M_h(\mathcal O_x),
\]
where \(h\) is the common rank. The projections to residual Hecke summands
and the auxiliary translations commute with it.

We can impose that the \(n\) smooth characters of
the compact diagonal torus at each place above
\(\ell\) be pairwise distinct. Index the free weight
coordinates by a place, a diagonal index
\(i\in\{0,\ldots,n-1\}\), and a generator of the free
unit lattice at that place. In each coordinate the
specialization has the form \(c\zeta-1\), where the
fixed offset \(c\) accounts for the reference weight
and the reciprocity normalization. Fix an integer
\(M\ge1\) such that every offset satisfies
\(c^M\equiv1\) in the residue field. For each coordinate
choose an \(M\)-th root \(\omega\) with
\(c\omega\equiv1\). Take
\[
 \zeta=\omega\zeta_{\ell^{M(i+nb)}}
\]

for \(b\ge1\), using primitive roots for the indicated
\(\ell\)-power orders. Raising to \(M\) removes \(\omega\) and gives
\[
 \operatorname{ord}(\zeta^M)
   =\frac{\ell^{M(i+nb)}}{\gcd(M,\ell^{M(i+nb)})}
   =\ell^{M(i+nb)-v_\ell(M)} .
\]
This exponent is positive and strictly increases with \(i+nb\).
The integers \(i+nb\) for different \(i\) are disjoint modulo
\(n\), so roots assigned to different diagonal indices cannot
be equal. Indeed equality of the two selected roots would force equality of their orders after raising to \(M\). The strict dependence of that order on the exponent gives
\[
 i+nb=i'+nb'\quad\Longrightarrow\quad i\equiv i'\pmod n
 \quad\Longrightarrow\quad i=i'
\]
for the specified representatives \(0,\ldots,n-1\).

At each place the character values on the
first free unit generator are among these \(\zeta\)'s,
up to the common choice of inversion in reciprocity, so
they detect all the required differences. The finite
torsion character stays fixed. These choices are independent
and are smooth, since a character of the free unit lattice
with finite-order values on a basis has open kernel.

Each of the coordinate sets \(c\zeta-1\) is infinite
inside the open unit disc. A nonzero integral power
series in one variable has only finitely many zeros
there by Weierstrass preparation. For example, write
\[
 F=\sum_{j\ge0}F_j(X_2,\ldots,X_d)X_1^j.
\]

If \(F\) vanishes at all tuples, then at every tuple in the
last \(d-1\) sets the one-variable series vanishes on the
infinite first set. Since the resulting one-variable series is zero,
\[
      F_j(x_2,\ldots,x_d)=0
      \qquad\text{for all }j\text{ and all remaining tuples}.
\]
Induction on \(d\) gives \(F=0\). Each tuple uses a finite extension of
the coefficient field, so the partially evaluated
series in this argument still have integral coefficients.
For a free basis \(e_i\) and \(v=\sum_i F_i e_i\), detection of the module follows coefficientwise:
\[
 v(x)=0\text{ at every selected }x\quad\Longrightarrow\quad
 F_i(x)=0\text{ for all }i,x\quad\Longrightarrow\quad F_i=0\text{ for all }i .
\]

The global determinant action uses the classical
Galois attachments for the admitted packets. The
inputs here are the polarized attachment, the
ordinary local flag with its ordered characters
at \(\ell\), the prescribed inertia at the tame
type, and local--global compatibility at the
split auxiliary places. On the residual summand consider the finite algebra
\[
 \Lambda\big[T,U:T\text{ spherical away from the level},
       \ U\text{ a normalized ordinary translation at }\ell\big]
\]
in the endomorphism ring. Its action determines the determinant.

At a smooth weight just constructed, take the complex realization through the chosen coefficient isomorphism.
Write classical automorphic forms as sums of forms
in the admitted automorphic constituents. This
decomposition gives a spanning statement also after
the local conditions are imposed. Apply the tame
level and type projection, and average against the
smooth torus character on \(J_d\) at a common depth
\(d\) at which the character extends. If \(P\) is this composite projector and \(V_\pi\) is an admitted constituent, then
\[
 P(V_\pi)\subset V_\pi,\qquad P\Bigl(\sum_\pi V_\pi\Bigr)=\sum_\pi P(V_\pi).
\]
Thus imposing the local conditions preserves the spanning description. The two sets of generalized eigenvalues give the coprime factorization
\[
 m=m_{\rm ord}m_{\rm other},\qquad \gcd(m_{\rm ord},m_{\rm other})=1
\]
of the minimal polynomial on this finite-dimensional space.

Choose \(p\) by the Chinese remainder theorem.
\[
 p\equiv1\pmod{m_{\rm ord}},\qquad
 p\equiv0\pmod{m_{\rm other}},\qquad P_{\rm ord}=p(U_z).
\] Apply it to each summand. If \(P_{\rm type},P_\theta,P_{\rm ord}\) are these three
projections in the order applied, follow them by the
residual Hecke idempotent \(e_{\mathfrak m}\) acting on
the ordinary module. For the constituent decomposition
\(V_x=\bigoplus_\pi V_x(\pi)\), the projected space is
\[
 e_{\mathfrak m}P_{\rm ord}P_\theta P_{\rm type}V_x
    =\bigoplus_\pi
          e_{\mathfrak m}P_{\rm ord}P_\theta P_{\rm type}V_x(\pi).
\]
Each operator preserves each summand on its domain. Thus each
nonzero summand carries the corresponding constituent's
local parameters.

On each nonzero intersection the spherical Hecke
operators act through the spherical character of the
constituent. At a place above \(\ell\), use the Jacquet
map for the upper unipotent group over the local field.
If \(v\) is in its kernel, write
\(v=\sum_i(\pi(n_i)-1)v_i\) with \(n_i\) in the full
upper unipotent group. All \(n_i\) belong to \(z^{-a}Nz^a\)
once \(a\) is large. For \(H=z^{-a}Nz^a\) containing every \(n_i\), right invariance of probability Haar gives
\[
 \int_H \pi(h)(\pi(n_i)-1)v_i\,dh
 =\int_H \pi(hn_i)v_i\,dh-\int_H\pi(h)v_i\,dh=0 .
\]
For an \(N\)-fixed \(v\), double-coset decomposition writes
\(U_z^a v\) as this average followed by \(\pi(z^a)\) and the
index \([N:z^aNz^{-a}]\), with the indicated normalization.

If \(v\) is fixed by \(N\), write \(A_a v\) for this
probability average. The double coset sum gives
\[
       U_z^a v=[N:z^aNz^{-a}]\,\pi(z^a)A_av=0.
\]
Since \(U_z\) is invertible on its ordinary part, \(v=0\). In the quotient, every term in the sum for
an ordinary translation has the same image. The
normalized translation therefore acts as the torus
action times the fixed index and normalization factors.
A character quotient of the Jacquet module with the
specified torus type embeds the irreducible local
constituent into its principal series by Frobenius
reciprocity. The inducing characters
have pairwise distinct restrictions to the units,
hence no ratio is an unramified norm power. The
principal series is irreducible by the criterion for
unlinked singleton segments \cite[Theorem 4.2]{zelevinsky}.
The geometric lemma lists the permutations in its
Jacquet module with multiplicity one
\cite[Sections 2.9 and 2.12]{bernstein-zelevinsky}.

Let \(\eta_i\) be the restrictions of the inducing characters to units, numbered in the specified order. A Jacquet constituent indexed by \(w\) has this type precisely when
\[
 (\eta_{w(1)},\ldots,\eta_{w(n)})=(\eta_1,\ldots,\eta_n).
\]
The \(\eta_i\) are pairwise distinct, so only the identity permutation contributes. Its Jacquet space
is one-dimensional, on which every normalized
ordinary translation acts by the corresponding
torus scalar. Multiplicities of the same constituent
carry the same scalars. Thus each nonzero intersection
supplies a scalar evaluation of the core Hecke algebra.

Let \(\mathbb T\) be the core Hecke algebra and
\(\mathcal X\) the detecting evaluations just constructed,
with coefficient fields \(E_x\). The map
\[
                 \mathbb T\longrightarrow\prod_{x\in\mathcal X}E_x
\]
is injective. Indeed a kernel operator vanishes on every
contributing constituent at each chosen weight. Its action
on a basis vector therefore has coefficients vanishing at
all chosen weights and is zero by weight density.
The source is compact and the product of coefficient fields
is Hausdorff. Therefore the injective map is a homeomorphism
onto its closed image.
For a fixed Galois element \(g\), approximate its conjugacy
class by split good Frobenius classes. The image of their
\(i\)-th characteristic coefficient converges in the product
to the tuple \((a_i(r_x(g)))_{x\in\mathcal X}\), where \(a_i\)
is the \(i\)-th coefficient of \(\det(X-r_x(g))\)
under the common convention.

The tuple lies in the closed
image. Define the desired coefficient as the unique element
\[
 c_i(g)\in\mathbb T,\qquad
      \bigl(c_i(g)_x\bigr)_{x\in\mathcal X}
        =\bigl(a_i(r_x(g))\bigr)_{x\in\mathcal X}.
\] The tuple is continuous in \(g\), so its
preimage is continuous as well. The characteristic
coefficients on group elements determine the
whole determinant law. The integral determinant
formula for a sum expresses its coefficients in
the characteristic coefficients of cyclic products
of the summands \cite[Section 1]{chenevier}.
For the value \(R\) of any of the determinant, polarization or local relations, classical compatibility gives
\[
 R\in\bigcap_{x\in\mathcal X}\ker(\mathbb T\to E_x)=0 .
\]
Thus the identities descend to the finite algebra. This gives the
map from the universal ordinary determinant ring,
using the characters given by the normalized
ordinary torus operators. Its kernel on the base
module cuts out its support.

On replacing the new levels by hyperspecial
levels, the underlying ordinary functions on the
new and old systems agree, including their weight
characters. The enlarged Galois ramification set
removes some spherical operators from the list
used to select the residual block. The two blocks
are the same. At a classical eigensystem in the
common block, the characteristic polynomials modulo
the residue maximal ideal coincide with the
chosen residual polynomials at the remaining good
primes.

Write \(\overline P_g\) and \(\overline P'_g\) for the two residual characteristic polynomials. Their coefficients are continuous, so the dense set of good Frobenius classes gives
\[
 \overline P_g(X)=\overline P'_g(X)\qquad(\text{every }g).
\]
This includes the removed good primes, which are still hyperspecial at this level. The ordinary operators at \(\ell\)
are identical in the two lists. Separation by
classical eigensystems then identifies the integral
summands. The determinant actions agree after
pullback along the quotient of Galois groups,
by equality of the characteristic coefficients
at the common good primes and continuity.
The ordered local characters agree because the
weight character and normalized ordinary
operators agree. This is the hyperspecial module
identification needed in the transfer.

\end{proof}

\subsubsection{Integral transfer at the auxiliary places}
\label{par:ordinary-level}

The modules at the two levels have now been identified under augmentation. For patching one also needs the Hecke operators to correspond integrally, with the possible discrepancy in the trace relations controlled uniformly.

\begin{proposition}[Transfer and nilpotence at the auxiliary level]\label{prop:ordinary-integral-transfer}
At the split auxiliary places of this ordinary system the
augmentation comparison is integral and the double-coset
operators satisfy \eqref{eq:hecke-transfer-identities}.
For a trace relation \(r\) in the comparison and the product
\(f_N\) of the local factors below, the integral module \(M_N\) satisfies 
\[
 (f_Nr)^{n!}M_N=0.
\]

\end{proposition}

\begin{proof}
\proofstep{Finite group quotient and the augmentation trace.}

At a split auxiliary place let \(K_v=\mathrm{GL}_n(\mathcal O_v)\),
let \(I_v\) be the inverse image of the upper triangular
Borel in \(K_v\), and put \(I_{1,v}=\ker(I_v\to\Delta_v)\),
where
\[
 \Delta_v=((k_v^\times)_\ell)^n
\]
is the
\(\ell\)-primary part of the residue torus. If \(Q_v=|k_v|\)
and \(a_v=v_\ell(Q_v-1)\), each of its \(n\) factors
has order \(\ell^{a_v}\). Take products of the groups over the
auxiliary places, and write \(\Delta=\prod_v\Delta_v\).

For the ordinary module at product level \(I_1\), the
variables \(1+y_{vi}\) act by the generators of \(\Delta_v\).
The group algebra over the ordinary weight ring has the
presentation in Paragraph~\ref{par:ordinary-patching}. It is
local because it is finite over the local weight ring and its
residue algebra is the group algebra of a finite abelian
\(\ell\)-group in characteristic \(\ell\). For the residue field \(k\) the group algebra becomes
\[
 k[\Delta]\simeq k[y_{vi}:v,i]/(y_{vi}^{\ell^{a_v}}:v,i),
\]
since \((1+y)^{\ell^{a_v}}-1=y^{\ell^{a_v}}\) in characteristic \(\ell\). The augmentation ideal upstairs is generated by the \(y_{vi}\).

The trace from \(I_1\) to \(I\) gives
\[
            M_\Delta\longrightarrow M^\Delta,\qquad
              [m]\longmapsto\sum_{\delta\in\Delta}\delta m .
\]
On a free group-algebra module this sends the standard
coinvariant basis to the orbit-sum basis of the invariants. The same isomorphism holds on each
projective summand because both functors commute with direct
sums. Thus the trace identifies the augmentation fibre
at \(I_1\) with the fixed module at \(I\). It commutes
with the determinant action and the local torus operators.
We give the further comparison of this \(I\)-fixed module
with the hyperspecial module.

Use one auxiliary place temporarily, with residue size
\(Q\), uniformizer \(\varpi\), and write \(X_i\) for
\(Q^{i-1}\) times the torus operator at the diagonal tuple
with \(\varpi\) in slot \(i\). Put \(P=O[X_1,\ldots,X_n]\),
\(S=P^{\mathfrak S_n}\), and
\[
 \Delta_X=\prod_{i<j}(X_j-X_i),\qquad
 D_Q=\prod_{i<j}(X_j-QX_i),\qquad
 d(g)=\sum_{w\in\mathfrak S_n} w(g/\Delta_X).
\]
Then \(d\) takes values in \(S\), and the pairing
\((g,h)\mapsto d(gh)\) is perfect over \(S\).
Indeed \(P\) is the ordered splitting algebra of the
universal monic polynomial with coefficients the signed
elementary symmetric functions. Adjoin its roots in
order \(X_n,\ldots,X_1\), dividing out the linear factor
after each adjunction.

For a monic adjunction of degree \(r\) with root \(t\),
let \(d_r\) extract the coefficient of \(t^{r-1}\).
Its pairing matrix satisfies
\[
             d_r(t^it^j)=0\quad(i+j<r-1),\qquad
             d_r(t^it^j)=1\quad(i+j=r-1).
\] It is therefore
perfect. Lagrange interpolation expresses its coefficient
functional over the fraction field as the sum of the values
at the roots divided by the derivative of the polynomial.
Composing these functionals gives the displayed formula for
\(d\), with denominator \(\Delta_X\) at the ordered tuple.
\proofstep{Perfect pairings and double-coset identities.}

In particular this formula and perfectness hold integrally.

Write \(E_Q=\sum_w T_w\), with measure of the Iwahori
equal to one. This element acts as the unnormalized trace
\(I\to K_v\), since \(K_v\) is the union of the finite
Bruhat cells. For bases \((a_b)_b,(c_b)_b\) of \(P\) dual
over \(S\) under \(d\), the identities in the Iwahori
Hecke algebra are
\begin{equation}
\label{eq:hecke-transfer-identities}
 E_Q g E_Q=d(D_Qg)E_Q,\qquad
 \sum_b a_b E_Q c_b=D_Q .
\end{equation}
Here a polynomial means multiplication through the
Bernstein torus subalgebra. To see the identities, use
the Bernstein presentation
\cite[Sections 1.5 and 1.15]{haines-kottwitz-prasad}.
With the integral variables above, its action on the
spherical polynomial module is
\[
 T_{s_i}g=Q s_i(g)+(Q-1)X_i\frac{g-s_i(g)}{X_i-X_{i+1}}.
\]
The normalization follows by taking the ratios of the
Hecke translations at
\(\operatorname{diag}(\varpi\mathbf 1_i,\mathbf 1_{n-i})\)
for successive \(i\), and multiplying the resulting
\(i\)-th ratio by \(Q^{i-1}\). The coefficient of the reflection \(s_i\) in the fraction field is
\[
       Q-(Q-1)\frac{X_i}{X_i-X_{i+1}}
                  =\frac{X_i-QX_{i+1}}{X_i-X_{i+1}}.
\]

In the expansion of
\(T_w\) the coefficient of \(w\) comes from taking each
reflection term in a reduced expression and is nonzero.
The remaining terms are indexed by proper subwords.
This proves faithfulness on the algebra spanned by the
Bernstein monomials times \(T_w\), using independence of
the field automorphisms. The coefficient of the reversing
permutation \(w_0\) in \(E_Q\) is
\(w_0(D_Q/\Delta_X)\), since the reduced expression
encounters each positive root once. The equality
\(T_{s_i}E_Q=Q E_Q\) puts the image of \(E_Q\) in the
symmetric polynomials by the displayed action of \(T_{s_i}\).

Write \(E_Q=\sum_w c_w w\). Since its output is symmetric,
\[
 \sum_w c_w w=\sigma E_Q=\sum_w\sigma(c_w)\,\sigma w ,
 \qquad c_{\sigma w}=\sigma(c_w)
\]
by independence of the distinct automorphisms. The known coefficient at \(w_0\) therefore gives \(c_w=w(D_Q/\Delta_X)\) for every \(w\). Thus \(E_Q\) acts as \(g\mapsto d(D_Qg)\).
Both identities in \eqref{eq:hecke-transfer-identities}
follow from this action, using \(\sum_b a_b d(c_b h)=h\)
in the second. Faithfulness gives the identities before
passing to any global module.

Let \(V\) be the ordinary fixed module at \(I\) and \(H\)
the hyperspecial one. For all the auxiliary places at once,
take tensor products of the polynomial and symmetric
algebras, the coefficient functionals and the dual bases.
Let \(e\) be the composite of the traces \(E_{Q_v}\)
to hyperspecial level, and let \(D=\prod_v D_{Q_v}\).
The local identities commute across places. The action of
symmetric polynomials on \(H\) is specified by
\[
 e_k(X_{v,1},\ldots,X_{v,n})\big|_H
   =(-1)^k[Z^{n-k}]\det(Z-\rho(\varphi_v)).
\]
For this equality \(\det(Z-\rho(\varphi_v))\) denotes
the characteristic polynomial of the universal
determinant on the hyperspecial module.

At a separating
characteristic-zero evaluation, the usual Satake
isomorphism identifies the symmetric action with the
elementary symmetric functions of the algebraically
normalized Satake parameters. These are the roots of the
Galois Frobenius polynomial by local--global compatibility.
If the specialization of the vector is zero, the equality
on it holds as well. Separation thus proves the displayed
equality on the integral module. This shows that
\(P\otimes_S H\), with all places included in \(P,S\),
is precisely the module obtained by adjoining the ordered
Frobenius roots to \(H\).

Define \(F:P\otimes_S H\to V\) by \(F(g\otimes h)=gh\).
Define \(G:V\to P\otimes_S H\) by
\[
 G(v)=\sum_b a_b\otimes e(c_b v).
\]
The second map is \(P\)-linear. In fact the tensor
\(\sum_b a_b\otimes c_b\) corresponds to the identity
under the perfect pairing, so multiplication by \(g\)
on its first factor agrees with multiplication on the
second factor. Both maps commute with the global
determinant action, since the translations commute with
that action and hyperspecial restriction gives the same
determinant there. The second identity in
\eqref{eq:hecke-transfer-identities} gives \(FG=D\)
on \(V\).

For hyperspecial \(h\) the first identity gives
\(E_{Q_v}(g h)=d(D_{Q_v}g)h\) at that place. Indeed
\(E_{Q_v}h=[K_v:I_v]h\), and the nonzero index cancels
on the torsion-free integral module. Apply this at each place. The composite on a pure tensor is
\[
      GF(g\otimes h)=\sum_b a_b\otimes d(Dc_bg)\,h
                                      =Dg\otimes h .
\]
The last equality is duality of the bases over \(S\).
The polynomial \(D\) divides \(f_N\) from
Paragraph~\ref{par:ordinary-patching}. Multiply \(G\) by \(f_N/D\). The common scalar of the two composites becomes
\[
 (f_N/D)D=f_N.
\]
All of these maps
are integral. They respect projection to a marked
simultaneous residual tuple, since they commute with each
of the ordered variables. They also respect coefficient
extension and the subsequent finite truncations.

At the covering level \(I_1\), use the split torus action
on fixed vectors obtained by forming the dominant
translation operators and inverting them, then taking
ratios to obtain all translations. The Iwahori--Matsumoto
relations make each such translation invertible in the
pro-unipotent Iwahori Hecke algebra when \(Q\) is a unit
\cite[Theorem 1 and its corollary]{vigneras-iwahori}.
Pass to \(I_1\) by the idempotent averaging the
prime-to-\(\ell\) part of the residue torus. Writing \(e\) for this averaging idempotent and \(T\) for a translation, commutation gives the inverse explicitly in the corner algebra:
\[
 (eTe)(eT^{-1}e)=e=(eT^{-1}e)(eTe).
\]

The dominant translation coset sums have the same upper unipotent representatives at both levels. The torus operators for diagonal units
are the actions of the
residue torus. Under geometric Artin reciprocity the
marked Frobenius \(\varphi_v\) goes to \(\varpi\) and the
marked tame element \(\tau_v\) goes to a unit \(\delta_v\)
which generates the cyclic \(\ell\)-quotient of
\(k_v^\times\). Thus the action of \(1+y_{vi}\) is the
torus operator at the tuple supported in slot \(i\)
with value \(\delta_v\).

At a classical evaluation let \(\chi\) be a joint torus
eigencharacter of the relevant irreducible local
constituent at \(I_1\), and let \(\chi_i\) denote its
restriction to slot \(i\). The Jacquet comparison sends
this character to the inducing datum. In normalized
induction its characters are
\[
 \psi_i(a)=\chi_i(a)\,Q^{(i-1)\operatorname{ord}(a)}
                         |a|^{(n-1)/2}.
\]
The factor multiplying \(\chi_i\) is the square root of
the Borel modulus, whose exponent on the \(i\)-th
coordinate is \((n+1-2i)/2\). The common twist in the
algebraic normalization of local Langlands cancels
\(|a|^{(n-1)/2}\). After cancellation the Weil character on a reciprocity element is
\[
         a\longmapsto Q^{(i-1)\operatorname{ord}(a)}\chi_i(a).
\]
At \(a=\varpi\) this is the \(X_i\)-eigenvalue. At the unit
\(a=\delta_v\) it is the eigenvalue of \(1+y_{vi}\).

Suppose the local factor of \(f_N\) is nonzero at this
eigencharacter. The normalized inducing characters are
pairwise unlinked: a quotient by a single power
\(|\cdot|^{\pm1}\) would make two of the \(X_i\) values
differ by a factor of \(Q\). Hence the constituent is
the irreducible principal series of these characters
\cite[Theorem 4.2]{zelevinsky}. Its local parameter, in
algebraic normalization, is the sum of the listed
characters, with zero monodromy. By characteristic-zero
local--global compatibility, the attached Galois
representation has this Frobenius-semisimplified
Weil--Deligne parameter.

If only the Weil semisimplification is compared, the
distinct roots \(a_i\) still give simple Frobenius eigenlines.
The Frobenius relation for the monodromy operator \(N\)
gives, in one of the two reciprocal Frobenius conventions,
\[
                    (a_i/a_j-Q)N_{ij}=0 .
\]
The excluded pairs have neither ratio \(Q\) nor ratio
\(Q^{-1}\), so all entries are zero in either convention. The resulting local
Frobenius eigenspaces are one-dimensional and on each
one inertia acts by the corresponding Weil character.
\proofstep{Pass from the torus eigenvalues to the integral module.}

The calculation with torus eigenvalues gives relations after specialization. A nilpotence statement on the integral module is needed before these relations can be used in patching.


Consequently the two matrix relations for
\(\varphi_v,\tau_v\) in Paragraph~\ref{par:ordinary-patching}
hold, with the product in the second relation on the
left of \(\rho(\tau_v)-(1+y_{vi})\).

At a classical evaluation with local factor zero,
multiplying the trace relations by that factor gives
zero at the joint eigenvalues anyway. To obtain the
same relation on an integral module, for each irreducible constituent \(\pi\) use the Jacquet map
\[
 \pi^{I_1}\lhook\joinrel\longrightarrow J_B(\pi).
\]
For injectivity
one uses the Iwahori decomposition of \(I_1\) and a
strictly contracting diagonal element. A fixed vector
which maps to zero in the Jacquet module is killed by
averaging over a sufficiently expanded compact unipotent
subgroup. On the fixed space the corresponding
dominant translation, up to its nonzero index scalar,
is invertible. Let \(I(\chi)\) be the inducing principal series. The Bruhat filtration of its Borel Jacquet module \(J_BI(\chi)\) has one line for each permutation, so
\[
 \dim J_BI(\chi)=\sum_{w\in S_n}1=n!.
\]
The Jacquet functor is exact, so the dimension for the
constituent is at most \(n!\)
\cite[Sections 2.9 and 2.12]{bernstein-zelevinsky}.

In a simultaneous triangularization on this complex
space, a polynomial vanishing at each joint torus eigenvalue
has zero diagonal. In dimension \(h=\dim\pi^{I_1}\le n!\), take the flag \(V_j\) spanned by the first \(j\) basis vectors in the triangularization. Then
\[
 B V_j\subset V_{j-1}\quad\Longrightarrow\quad
 B^h V_h\subset V_0=0\quad\Longrightarrow\quad B^{n!}=0 .
\] Apply this to the trace relations
multiplied by the local factor of \(f_N\) and then
multiply by any remaining factors of \(f_N\). Separating
classical evaluations show that \((f_N r)^{n!}\)
annihilates the full integral module, where \(r\) is the
trace relation in question. The bound \(n!\) is independent
of the conductor of the tame character, the weight,
the depth \(N\) and the dimensions of the global modules.
\end{proof}

\subsubsection{Patching at an absolutely irreducible ordinary arc}
\label{par:ordinary-patching}

The finite-level transfer controls the trace relations before passing to a limit. Patching uses that control to deduce ordinary support along an arc when its generic representation satisfies the stated irreducibility and local separation conditions.

\begin{proposition}[Augmentation support below a marked arc]\label{prop:ordinary-one-point-patching}
In one rank-four determinant link, take a characteristic-\(\ell\)
ordinary arc through its ordinary support whose representation is
absolutely irreducible on the cyclotomic kernel of \(G_L\).
Assume the ordered characters separate on the prescribed inertia
groups, regular elements are dense in the cyclotomic kernel, and
the distinguished auxiliary place satisfies the local cyclotomic
intertwining condition of Proposition~\ref{prop:tower-below-arc}.
Assume \(\bar\epsilon(G_L)\ne\{1\}\), where \(\bar\epsilon\)
is the mod-\(\ell\) cyclotomic character.
With the finite-level modules and auxiliary comparisons of
Propositions~\ref{prop:ordinary-level-modules}
and \ref{prop:ordinary-integral-transfer}, if \(y\) denotes the arc and \(R_0\) the base ring of the Hecke ideal, then every generization satisfies 
\[
 P\subseteq\ker y\quad\Longrightarrow\quad J_0\subseteq P\cap R_0.
\]

\end{proposition}

\begin{proof}
\proofstep{One determinant link and its support condition.}

Fix one of the two determinant links and write \(R_0\) for its
universal ordinary determinant ring, \(J_0\) for its Hecke ideal,
and \(p\) for the evaluation kernel of the characteristic-\(\ell\)
arc. Write \(n=4\), \(d=[L^+:\Q]\), and
\[
 D=nd.
\]
We use the
one-point case in which \(p\) contains \(J_0\) and the arc
representation is absolutely irreducible on the cyclotomic kernel
of \(G_L\). Use the ordered local characters and auxiliary condition
of Proposition~\ref{prop:tower-below-arc}
and the residual tests of Lemma~\ref{hyp:ordinary-arc-residual-tests}. We describe
the passage from finite-level data to vanishing of \(J_0\) at the
generizations of \(p\).

Put
\[
 A=K[[t]]
\]
for the arc valuation ring after finite
extensions and \(E=K((t))\). The finite field \(K\) can be taken
to contain the coefficient residue field. Choose Taylor--Wiles
elements \(\sigma_j\), \(1\le j\le q\), in the full cyclotomic
kernel of \(G_L\) at which the characteristic polynomials are
separable. Here \(q\) counts the Taylor--Wiles places and is
fixed throughout the patching limit. The elements must detect
dual adjoint classes, as follows. On the cyclotomic kernel the
adjoint action twisted by the cyclotomic character becomes the
ordinary adjoint action. Write
\[
 V(g)=z(g)\rho(g)
\]
for a dual
cocycle \(z\) on that kernel. Its cocycle identity is
\[
 V(gh)=V(g)\rho(h)+\rho(g)V(h).
\]
Suppose \(\operatorname{tr} V(g)=0\) for every \(g\) in the
kernel.

If \(\sum_i a_i\rho(g_i)=0\), the lift identity
\(V(gh)=V(g)\rho(h)+\rho(g)V(h)\) gives
\[
 \operatorname{tr}\left(\Bigl(\sum_i a_i V(g_i)\Bigr)\rho(h)\right)
  = \sum_i a_i\operatorname{tr}V(g_i h)
        -\operatorname{tr}\left(\Bigl(\sum_i a_i\rho(g_i)\Bigr)V(h)\right)
  =0.
\]
Every \(g_i h\) here lies in the same cyclotomic kernel,
where the traces are assumed zero.
Absolute irreducibility gives the full span \(M_n(E)\).
Its trace pairing is perfect, so \(\sum_i a_i V(g_i)=0\).
Thus the proposed map on the spanning matrices is well-defined
and obeys the Leibniz rule.

Write \(D_\rho\) for this derivation and
\(e_{ij}\) for the matrix units. The matrix
\(Y=\sum_i D_\rho(e_{i1})e_{1i}\) implements the derivation.
Indeed the Leibniz rule applied to \(e_{ab}e_{i1}=\delta_{bi}e_{a1}\) gives
\[
 Ye_{ab}=D_\rho(e_{a1})e_{1b},\qquad
 e_{ab}Y=D_\rho(e_{a1})e_{1b}-D_\rho(e_{ab}).
\]
Hence \(V(g)=Y\rho(g)-\rho(g)Y\) on the cyclotomic kernel \(H\).
Subtracting the associated coboundary makes \(z|_H=0\).

For \(h\in H\) and \(g\in G_L\), apply the cocycle identity to
\(hg=g(g^{-1}hg)\). Normality and \(z|_H=0\) give
\[
 \rho(h)z(g)\rho(h)^{-1}=z(hg)=z(g).
\]
The span of \(\rho(H)\) is the full matrix algebra.
Its commutant is \(E\cdot I_n\), so \(z\) descends to a scalar
cocycle on the abelian quotient \(G_L/H\), with the residual
cyclotomic action \(\epsilon=\bar\epsilon\).
For the fields used here a place with nontrivial residual cyclotomic
value is provided by the marked-field constructions in
Propositions~\ref{prop:count-ii-auxiliary-primes}
and~\ref{prop:count-three-marked-elements}. Choose \(a\) in
the quotient with \(\epsilon(a)\ne1\) and put
\(u=z(a)/(\epsilon(a)-1)\). Since \(ab=ba\),
\[
 z(a)+\epsilon(a)z(b)=z(b)+\epsilon(b)z(a)
 \quad\Longrightarrow\quad z(b)=(\epsilon(b)-1)u .
\]
This exhibits the coboundary.

Passing from the real group to its index-two CM subgroup
loses no classes, since \(2\) is invertible. The regular
elements are dense in the kernel, so continuity suffices to
test traces on that locus. The space of Laurent series classes
is finite-dimensional by finite generation after passage to
the pro-\(\ell\) quotient. Choose a finite set of its trace
tests. After an extension of \(A\), each marked characteristic polynomial takes the form
\[
 P_{\sigma_j}(X)=\prod_{a=1}^n (X-\alpha_{ja}),\qquad
 \alpha_{ja}\in A,\qquad \prod_{a<b}(\alpha_{ja}-\alpha_{jb})\ne0\text{ in }E .
\]
The differences need only be nonzero in the fraction field here.

For every depth \(N\) choose distinct split places away from
the initial level with Frobenius approaching \(\sigma_j\) and
norm
\[
 Q_{j,N}\equiv1\pmod{\ell^N}.
\]

For precision \(m\), scale the fixed finite list of integral
cocycles once and record their reductions alongside the residual
matrix action modulo \(t^m\). For a cocycle \(z\), conjugation on its coefficient module gives the homomorphism
\[
 h\longmapsto (z(h),r(h))\bmod t^m,\qquad
 (u,A)(v,B)=(u+\operatorname{Ad}(A)v,AB).
\]

The coefficient and matrix rings modulo \(t^m\) are finite. Intersect their open kernels
with the kernel of the cyclotomic character modulo \(\ell^N\);
take the normal core in the group over the totally real field,
also recording its CM quotient. All these quotients are finite.
Use \(m\ge N\) tending to infinity in the resulting sequence.
Chebotarev gives infinitely many primes in each specified
conjugacy class, so we may avoid all previous choices and the
fixed bad primes.
In the resulting finite quotient \(G_{m,N}\) the target class is
\[
          [\sigma_j]_{G_{m,N}},\qquad
          \epsilon_\ell(\sigma_j)\equiv1\bmod\ell^N .
\]
The split condition comes from its membership in the CM
subgroup. Use the inverse conjugacy class for geometric Frobenius.
Choose the tame generators in the pro-\(\ell\) inertia factors. The open subgroup on which all
residual characters and the mod-\(\ell\) cyclotomic character
are trivial has pro-\(\ell\) image for the arc.

In fact its
residual matrices are unipotent, so the finite residual
image is an \(\ell\)-group by Cauchy's theorem. At depth \(m\) the kernels are filtered by the groups
\(1+t^aM_n(K[t]/(t^m))\), with successive quotients embedded in additive \(M_n(K)\) by
\[
 1+t^aX\mapsto X\bmod t.
\]
Indeed
\[
 (1+t^aX)(1+t^aY)\equiv1+t^a(X+Y)\pmod{t^{a+1}}
 \qquad (a\ge1).
\] Thus each is an \(\ell\)-group. This justifies the use of the
maximal pro-\(\ell\) quotient on this subgroup. The finite
levels are compared on a common topologically finitely
generated profinite group, with stable quotients in every
fixed finite approximation along an ultrafilter
\(\mathcal U\) containing the cofinite filter.
Frobenius and tame inertia converge in these quotients.

To obtain the common group,
bound the generator ranks of the open pro-\(\ell\)
subgroups uniformly. A homomorphism to \(\mathbf F_\ell\)
which kills all lifts of the new tame generators descends
to the base. If \(d_0\) is the base generator rank, \(q\) the number
of new places, and \(I_0\) the fixed open subgroup index,
the evaluations on at most \(qI_0\) such lifts give
\[
 \dim_{\mathbf F_\ell}\operatorname{Hom}_{\rm cont}
      (U_N,\mathbf F_\ell)\le d_0+qI_0 .
\]
The Frattini quotient identifies the left side with the
generator rank of the pro-\(\ell\) group \(U_N\).
Use a free profinite group on this many generators
together with lifts of the fixed finite quotient.

In each finite quotient of the common group there are
finitely many possibilities for the image of the level
kernel. The ultrafilter selects compatible stabilized
images. Their inverse limit is the kernel of the limiting
quotient. The decomposition maps are carried along in
this construction. The ring \(R\) below uses the
determinant on the common group, and \(I\) imposes
its descent to the limiting quotient.

Use a coefficient DVR \(O\) with residue field \(K\) after
finite unramified coefficient extension, and adjoin a variable
\(s\) on the base, evaluated at \(t\). These extensions are
flat. The completed ordinary weight algebra on the base is
\[
 \Lambda_0=O[[s,w_1,\ldots,w_D]].
\]
Its specialization to \(K[[t]]\) is onto because it sends
\(s\) to \(t\). The completion of \(\Lambda_0\) at the
kernel of this specialization is regular local of dimension
\(1+D\). Indeed the quotient has dimension one, and
\(\Lambda_0\) is regular local of dimension \(2+D\).

At the finite levels retain one tame character variable
\(y_{ji}\) for each Frobenius eigenline at each of the \(q\)
auxiliary places. Put \(\mathfrak y=(y_{ji})\). The required
finite-level ordinary module \(M_N\) is finite free over
\[
 \Lambda_0[[y_{ji}]]\Big/
 \bigl((1+y_{ji})^{\ell^{a_{j,N}}}-1\bigr)_{j,i},
 \qquad a_{j,N}\ge N .
\]

Adjoin the ordered roots of the characteristic polynomials at
the \(\sigma_j\) over the extended base ring, using one monic
polynomial and then its successive monic quotients. If \(P_0\) is the polynomial of degree \(n\) at one place, the successive extensions are
\[
 C_i=C_{i-1}[z_i]/P_{i-1}(z_i),\qquad
 P_i(X)=P_{i-1}(X)/(X-z_i),\qquad
 \operatorname{rank}_{C_{i-1}} C_i=n+1-i .
\]

Thus ordering the roots is finite free, including when the residual
polynomial has repeated roots. If the residual labeling singles out a
maximal ideal, take the corresponding direct factor \(R_0'\):
over a complete local base the idempotents of the finite special
fibre lift uniquely by Hensel's lemma. This summand is finite
projective, hence free over the local base. If \(C\) is the
finite splitting algebra and \(\bar e\) the residual projector, this
construction gives
\[
             C=eC\oplus(1-e)C,\qquad e^2=e,\quad e\bmod\mathfrak m=\bar e .
\]
Evaluation at the
chosen roots selects this summand and lifts the arc. Differences
which vanish at its closed point can still be nonzero at the
generic point of the arc. The base module \(X\) is the ordinary
Hecke module after this base change. It is finite free
over \(\Lambda_0\); its support is the inverse image of
\(V(J_0)\).

At a finite level use the roots \(U_{ji}\) of
the corresponding Frobenius. The maps between the
zero-\(\mathfrak y\) fibre of \(M_N\) and the base module
with those ordered roots have both compositions equal to
multiplication by
\[
 f_N=\prod_j\prod_{i}\prod_{i'\ne i}
 (U_{ji}-U_{ji'})(U_{ji}-Q_{j,N}U_{ji'}).
\]
These are the Hecke-module freeness and level comparison
inputs. The second factor is needed in the comparison
at finite depth. At the limiting arc both factors become
the same root difference, which is nonzero in \(E\).
The integral construction at these levels is described
in Paragraph~\ref{par:ordinary-level}. It gives the
trace at the tame cover, the two comparisons and the
uniform exponent for the local relations.

Here is the algebraic form of the remaining finite-level
compatibilities. One has one complete Noetherian local
ring \(R\) acting on the modules, with coefficient map
\[
 \Lambda_0[[y_{ji}]]\to R.
\]
The coefficients of the determinant of the limiting group supply the action of \(R\).
There is a surjection \(R\to R_0'\) killing \(\mathfrak y\),
and an ideal \(I\) in its kernel which imposes the limiting
ordinary, polarized, and prescribed inertia relations.
The Frobenius and tame trace relations use the following
elements in the group algebra or its Cayley--Hamilton quotient:
\[
       \prod_i(\sigma_j-U_{ji}),\qquad
       \Big(\prod_{i'\ne i}(\sigma_j-U_{ji'})\Big)
                       (\tau_j-1-y_{ji}).
\]
At the limiting point the marked Frobenius has simple spectrum. Its commutation with the tame element implies that each Frobenius line \(L_{ji}\) is tame-stable. The second relation then reads
\[
 \rho(\tau_j)|_{L_{ji}}=(1+y_{ji})\,\mathrm{id}_{L_{ji}}
\]
after the distinct-root differences have been inverted. Pairing these
relations with the traces of all group elements gives
coefficient relations.
\proofstep{Comparison element and patching depth.}

For the proof below the comparisons give an element
\(f\in R\) specializing to the product of the squared
root differences and the following two properties.
\begin{itemize}
\item For every fixed \(b\), for \(\mathcal U\)-almost
all \(N\) there are \(R\)-linear maps in both directions
between \(X/\mathfrak m_{\Lambda_0}^bX\) and
\[
 M_N/(\mathfrak y,\mathfrak m_{\Lambda_0}^b)M_N,
\]

with both composites multiplication by \(f\).
\item For a generator \(r\) of \(I\), some fixed power
\((fr)^e\) agrees modulo each fixed power of \(\mathfrak m_R\)
with an element annihilating \(M_N\) for
\(\mathcal U\)-almost all \(N\). The exponent is chosen
before the accuracy and depth.
\end{itemize}
The second property permits congruences of the action at
different finite levels. Here the generators of \(I\) are the trace relations for the group kernel,
polarization, ordinarity, the prescribed inertia, and
the two Frobenius and tame relations just described.

We prove the patching comparison for these data.
For \(k,b\ge1\) let \(V_{k,b}\) be the ultraproduct
along \(\mathcal U\) of
\[
 M_N/(\mathfrak y^k,\mathfrak m_{\Lambda_0}^b)M_N.
\]

For fixed \(k,b\) the group-ring relations are in
\((\mathfrak y^k,\mathfrak m_{\Lambda_0}^b)\) for all
large \(N\). For \(1\le c<k\le\ell^a\), the identity
\[
 c\binom{\ell^a}{c}=\ell^a\binom{\ell^a-1}{c-1}
 \quad\Longrightarrow\quad
 v_\ell\binom{\ell^a}{c}\ge a-v_\ell(c)
\]
is uniform in these finitely many \(c\) as \(a\) tends to infinity.
\proofstep{Compatible bases and the augmentation fibre.}

Choose a free basis of \(M_N\) over its coefficient group
ring at each \(N\). For fixed \(k,b\) the binomial valuation
bound kills all relations of positive \(y\)-degree less than
\(k\) modulo \(\mathfrak m_{\Lambda_0}^b\) for large \(N\).
Writing \(d_N\) for the size of the chosen basis, the
basis therefore gives the additive decomposition
\[
 M_N/(\mathfrak y^k,\mathfrak m_{\Lambda_0}^b)M_N
   \simeq\bigoplus_{|\alpha|<k}
       (\Lambda_0/\mathfrak m_{\Lambda_0}^b)^{d_N}y^\alpha .
\]
Reduction in \(b\) acts on the coefficients and reduction
in \(k\) discards the extra monomials. Thus both commute with
the identifications from the fixed bases. For fixed \(k\), take the inverse limit
over \(b\), giving \(V_k\). Multiplication by the monomials
identifies the successive graded pieces with copies of
\(V_1\). These identifications are \(R\)-linear, because
multiplication by each \(y_{ji}\) commutes with the
\(R\)-action.

At \(k=1\), the comparison maps give the ultraproduct and inverse
limit maps
\(X \xrightarrow{u} V_1\xrightarrow{v} X\).
To arrange compatibility in \(b\), choose maps at a depth
of truncation tending to infinity along the ultrafilter
and use their reductions. Indeed, if \(S_j\in\mathcal U\) is the good set for truncation \(j\), then
\[
 T_b=\bigcap_{j=1}^b S_j\in\mathcal U,\qquad T_{b+1}\subseteq T_b .
\]
On \(T_b\) maps at depth \(b\) can be reduced compatibly to every smaller truncation. At level \(N\) select the largest
available truncation up to \(N\), which gives the
required cofinal choices. The
ultraproduct of the constant finite module
\(X/\mathfrak m_{\Lambda_0}^bX\) is that same module.
Their composites satisfy
\[
      vu=f\,\mathrm{id}_X,\qquad uv=f\,\mathrm{id}_{V_1}.
\]
Inverting \(f\) gives \(V_1[1/f]\simeq X[1/f]\).

Now form
\[
 P=\varprojlim_k V_k[1/f]
\]
over the \(\mathfrak y\)-adic completion of \(R[1/f]\).
Flatness of localization preserves the transition kernels. The
monomial expansions identify
\[
   \operatorname{gr}_{\mathfrak y}P\simeq X[1/f][y_{ji}],
           \qquad P/\mathfrak yP\simeq X[1/f].
\]
The degree-\(a\) identification is multiplication by all monomials
of total degree \(a\).
For this assertion, a compatible sequence with zero
constant term is a sum of the \(y_{ji}\) times compatible
sequences. Solve this degree by degree with the
monomial identifications. The same argument applies in
each degree. Finite generators of \(X[1/f]\) lift to
generators of \(P\) by lifting the coefficient at each
successive degree and summing in the complete ring.

The relations in \(I\) act nilpotently on \(P\). To see
why congruences in the maximal ideal of \(R\) suffice,
fix \(b,k\). Some \(\mathfrak m_R^c\) kills
\(X/\mathfrak m_{\Lambda_0}^b X\), a finite module over
the local ring \(R\). The comparison at this truncation
therefore gives \(f\mathfrak m_R^cV_{1,b}=0\).
Iterating on successive \(\mathfrak y\)-quotients gives
\[
 f^k\mathfrak m_R^{ck}V_{k,b}=0.
\]
Use the annihilator approximation modulo
\(\mathfrak m_R^{ck}\). For every \(k,b\), the bound gives \(f^k(fr)^eV_{k,b}=0\).
After taking the inverse limit in \(b\) and inverting \(f\),
\[
             r^e V_k[1/f]=0 .
\]
Since \(e\) is independent of \(k\), for each compatible sequence \(v=(v_k)_k\in P\) we have
\[
 r^e v=(r^ev_k)_k=(0)_k.
\]
This gives \(r^e P=0\) on the inverse limit.

Let \(p_R\) be the inverse image in \(R\) of the lifted
arc prime in \(R_0'\), and put
\[
 C=\widehat{R_{p_R}}.
\]
Its residue field is \(E\). At each quotient by \(p_R^m\), the
inclusion \(\mathfrak y^m\subseteq p_R^m\) makes the preliminary
\(\mathfrak y\)-completion redundant, and \(f\notin p_R\).
Taking the localized inverse limit gives the same ring \(C\).
This base change on the intermediate Noetherian ring is
flat. Write \(P_C\) for the completed local base change
of \(P\), and \(\Lambda_p\) for the regular completed
localization of \(\Lambda_0\) above. The augmentation
\(P_C/\mathfrak y P_C\) is finite free over
\(\Lambda_p\).

To see this, the action on \(X\) factors
through a finite \(\Lambda_0\)-algebra inside
\(\operatorname{End}_{\Lambda_0}(X)\). After base change to the complete local ring \(\Lambda_p\), write the finite algebra as \(\prod_{\mathfrak q}T_{\mathfrak q}\) and its idempotents as \(e_{\mathfrak q}\). Then
\[
 X\otimes\Lambda_p=\bigoplus_{\mathfrak q} e_{\mathfrak q}(X\otimes\Lambda_p).
\]
The summand for \(p_R\) is projective over the local ring \(\Lambda_p\), hence finite free. This factor is
nonzero by \(J_0\subset p\). Flat base change preserves
the monomial identifications.

If \(P_C/\mathfrak yP_C\) has rank \(h\) over \(\Lambda_p\),
lift a basis. Successive monomial coefficients give the surjection
\(\Lambda_p[[y_{ji}]]^h\to P_C\). Indeed, put \(S=\Lambda_p[[y_{ji}]]\) and let \(e_i\) be the chosen lifts. A nonzero relation would have a least total degree \(a\). Thus its initial part would give
\[
 \sum_{i,\,|\alpha|=a} c_{i,\alpha}\,y^\alpha\bar e_i=0
 \quad\text{in }\mathfrak y^aP_C/\mathfrak y^{a+1}P_C .
\]
The monomial identification makes these vectors a \(\Lambda_p\)-basis, so all the initial coefficients vanish. Hence \(S^h\simeq P_C\). Completeness of \(C\) and finiteness
of \(P_C\) give completeness for the \(\mathfrak y\)-adic
lifting.

The support algebra \(C/\operatorname{Ann}_C P_C\)
embeds into the endomorphism algebra of this finite free
module over \(\Lambda_p[[y_{ji}]]\). It is finite and
torsion-free over that regular local ring. Every minimal
prime of the support algebra contracts to zero, since
each nonzero element of the regular ring is a
nonzerodivisor and minimal primes are associated.
Dimension for the finite integral extension now shows
that every minimal support prime \(Z\) has
\[
 \dim(C/Z)=1+D+nq .
\]
Each \(Z\) contains \(IC\), by the annihilator
calculation above. For \(Q\supseteq\mathfrak y\), Nakayama at \(Q\) gives
\[
    (P_C)_Q\ne0\quad\Longleftrightarrow\quad
                      (P_C/\mathfrak yP_C)_Q\ne0 .
\] In terms of \(R_0'\) it is
therefore equivalent to containment of \(J_0\) in
the contraction.

We explain the upper bound on the relative cotangent
dimension of \(C/IC\) needed to use the support.
The coefficient DVR at this point can be taken to be
the completed localization of \(O[[s]]\) at the
coefficient uniformizer. Its residue field is \(E\),
because \(s\) specializes to \(t\).
If \(T_{\rm rel}\) is the tangent space relative to this DVR, the first-order estimate to prove is
\[
 \dim_E T_{\rm rel}\le D+nq.
\]
Here the rank-four point has one
absolutely irreducible matrix block. The localized
Cayley--Hamilton algebra is the full matrix algebra
\cite[Theorem 2.22(i)]{chenevier}, since its determinant
over the residue field is that of the absolutely
irreducible representation over \(E\).
\proofstep{Recover matrices and flags from determinants.}

Thus a relative tangent gives a matrix cocycle modulo
infinitesimal conjugation. The polarization gives the transpose-inverse relation. For a fixed pairing matrix \(J\), an infinitesimal change \(1+\delta X\) with \(\delta^2=0\) preserves it exactly when
\[
 (1+\delta X)^{\mathsf T}J(1+\delta X)=J
 \quad\Longleftrightarrow\quad X^{\mathsf T}J+JX=0 .
\]
Normalization of the pairing uses division by two. The ordinary flags and their
character values are uniquely specified at first order
by the ordered eigenlines of a separating local
element. At each split \(\ell\)-place let \(\mu_v\) be the
preimage in the local Galois group of the \(\ell\)-primary
torsion in its topological abelianization. The ordered
characters \(\psi_{v,i}\) have prescribed values on
\(\mu_v\) from \(O^\times\). This is a condition in the
universal ordinary determinant ring, so an \(O\)-linear
tangent lift has the same values there, viewed as constants
in \(E[\delta]/(\delta^2)\).

For \(g\in\mu_v\) equality with the prescribed constant gives
\[
 \psi_{v,i}(g)(1+\delta c_{v,i}(g))=\psi_{v,i}(g)
 \quad\Longrightarrow\quad c_{v,i}(g)=0 .
\]
Here \(c_{v,i}(g)\in E\) and \(\psi_{v,i}(g)\) is a unit.
This condition survives reduction of an integral cocycle
modulo \(t^m\).
For the existence of the local flag, the trace relations and perfectness of the trace pairing give the ordered product identity
\[
 \prod_{i=1}^n\bigl(\rho(g_i)-\psi_i(g_i)I\bigr)=0\qquad(\text{all }g_1,\ldots,g_n).
\]
Take the same separating element at every
slot. Its distinct roots give the ordered spectral
projections, including over the dual numbers.
Replace only the \(k\)-th slot by \(g\). The factors
on the left kill the rows with index less than \(k\),
and those on the right kill the columns with index
greater than \(k\).

On the remaining rows and columns
they act by units. More explicitly, write \(L_k,R_k\) for the two products
of spectral factors at the separating element. For the
surviving rows and columns, the product equation reads
\[
 0=\bigl(L_k(\rho(g)-\psi_k(g))R_k\bigr)_{rc}
       =(L_k)_{rr}(\rho(g)-\psi_k(g))_{rc}(R_k)_{cc},
             \qquad r\ge k,\quad c\le k .
\]
Both outer factors are units by the distinct spectral values.
Taking \(k=c\) for each surviving column yields
\[
 \rho(g)_{rc}=0\ (r>c),\qquad \rho(g)_{cc}=\psi_c(g).
\]
Indeed the outer scalar factors are units, and the subtracted diagonal term occurs only when \(r=c\).
At the \(q\) new places simple Frobenius roots likewise
give the ordered lifts uniquely. The two tame and
Frobenius relations determine the new character
variables on those eigenlines. The determinant
coefficients and these parameters generate the chart.

Consequently the relative tangent injects into the
ordinary first cohomology of the limiting group with
its prescribed inertia conditions. For continuity of
the derivative, the kernel of evaluation onto \(K[[t]]\)
has a finite conormal module over \(K[[t]]\). The
derivative on it is bounded after multiplication by a
power of \(t\). It kills \(O[[s]]\), and if \(I\) denotes the evaluation kernel, the approximation by coefficient series reads
\[
 \mathfrak m_R^b\subseteq s^bO[[s]]+(I\cap\mathfrak m_R^b).
\]
Artin--Rees then gives a bound
tending to zero with \(b\). This applies to the
determinant coefficients and to the matrix coordinates
recovered by the inverse trace Gram matrix.

\proofstep{Local and global bound on relative tangents.}

For the dimension bound on this cohomology use the
first-order finite-level count. Over \(K[t]/(t^m)\)
it has the form
\[
 \bigl|\mathrm{Sel}_1(N,m)/\mathrm{Bd}_1(N,m)\bigr|
     \le |K|^{m(nd+nq)+c}.
\]
Here the bound holds for \(N\) sufficiently large for
each \(m\), with \(c\) independent of \(m,N\).
In the Euler-characteristic comparison let
\(\nu=n(n-1)/2\). At a split \(\ell\)-place, with
local degree \(d_v\) and principal-inertia rank
\(d'_v\), ordinary upper triangular cocycles contribute
at most
\[
 m\bigl(n(d'_v+1)+\nu(1+d_v)\bigr)+O(1)
\]
to the logarithmic count.

For the diagonal characters,
local class field theory identifies the free pro-\(\ell\)
part of inertia after quotienting by its torsion with
\(\mathbf Z_\ell^{d'_v}\). The vanishing on \(\mu_v\) leaves exactly \(d'_v\) inertia
parameters and one unramified parameter per character. Thus the
diagonals contribute logarithmic size at most \(mn(d'_v+1)\). For
upper entries filter the triangular cocycle space by
increasing distance from the diagonal. When all
preceding entries vanish, the \((i,j)\)-entry is a
cocycle with coefficients of character
\[
 \chi_{ij}=\psi_{v,i}/\psi_{v,j}.
\]

If \(\epsilon(g)/\chi_{ij}(g)-1\) has valuation \(b\)
at one local element, the invariants of the dual twist in
\(A/(t^m)\) lie in the kernel of multiplication by that scalar.
Over \(A/(t^m)\) the scalar generates the same ideal as \(t^b\). Consequently
\[
 \ker(t^b:A/(t^m)\to A/(t^m))=t^{\max(m-b,0)}A/(t^m),\qquad
 |\ker|=|K|^{\min(b,m)} .
\]
Writing \(D_{ij}\) for the rank-one coefficient module,
local duality, the Euler formula and the boundary size give
\[
 \begin{aligned}
 |Z^1(D_{ij})|
  &= |H^1(D_{ij})|\,|D_{ij}|/|H^0(D_{ij})|\\
  &= |D_{ij}|^{1+d_v}|H^0(D_{ij}^{\vee}(1))|
   \le |K|^{m(1+d_v)+b}.
 \end{aligned}
\]
Multiply these bounds at the successive strictly upper
positions of the filtration.

At the new places the unrestricted condition contributes
the invariants of the dual twist by local duality and
the tame Euler characteristic. For large \(N\) the
Frobenius matrix and cyclotomic value modulo \(t^m\)
are those of \(\sigma_j\). In the generic Frobenius eigenbasis the commutator has
multipliers \(a_i-a_j\) on off-diagonal entries. Hence
\[
     \dim_E\ker\operatorname{ad}(\rho(\sigma_j))=n.
\]
Only the diagonal entries contribute to the kernel rank. Smith normal form bounds the
size of the corresponding kernel modulo \(t^m\) by
\(|K|^{mn+O(1)}\).

The same argument bounds the term at
the old distinguished place by \(O(1)\) in logarithmic
size, using the vanishing local intertwining space.
The fixed-inertia primes give the unramified Euler
terms in the comparison. The global invariants
over \(E\) vanish. Schur's lemma on \(G_L\) makes an invariant \(X=\lambda I\). Invariance under complex conjugation then gives
\[
 \lambda I=-{}^t(\lambda I)=-\lambda I,\qquad 2\lambda=0 .
\]
The characteristic is odd, so \(X=0\). Finitely
many equations already have this kernel, so they
give a bounded invariant count at every truncation.

The dual auxiliary restrictions descend to the base
group, since they impose zero on the new tame generators,
and their Frobenius traces satisfy the chosen detection
tests. To bound this count modulo \(t^m\), work on the
topologically finitely generated group after the
pro-\(\ell\) quotient. The reduction of a cocycle also
factors through this quotient, by applying the
pro-\(\ell\) image argument to the cocycle and its
representation in their semidirect product.
A cocycle on the free profinite group with generators \(g_1,\ldots,g_m\) is determined by the tuple
\[
 z\longmapsto\bigl(z(g_1),\ldots,z(g_m)\bigr).
\]
Descent imposes
linear equations for relations. The rows of all these
equations and the finite trace tests span a finite
module of linear forms over \(A\), so finitely many
rows suffice even modulo \(t^m\).

Over \(E\) their
kernel consists of coboundaries by detection. Smith
normal form shows that the kernel modulo \(t^m\)
has size at most
\[
 |K|^{mr+O(1)}
\]
when its rank
over \(E\) is \(r\). For the coboundary map of the same generic rank \(r\), the
image has size at least \(|K|^{mr-O(1)}\). Thus division gives
\[
              |K|^{mr+O(1)}/|K|^{mr-O(1)}=|K|^{O(1)}.
\]
This bound is uniform in the truncation depth. These losses depend
on the fixed arc, the chosen local elements and the
finite trace tests. They are independent of \(N\) and
\(m\) in the indicated order of approximation.
Greenberg--Wiles duality subtracts
\[
 m\bigl((n+\nu)\#\{v\mid\ell\}+\nu d\bigr).
\]
With \(B_v\) the upper triangular matrix module at this
depth and \(A_v\) the full adjoint coefficient module, the
triangular coboundaries have order
\(|B_v|/|H^0(B_v)|\).

The ordinary local condition \(L_v\)
is an image of the triangular cocycle space, and dividing
by these coboundaries gives
\[
 \frac{|L_v|}{|H^0(A_v)|}
       \le \frac{|Z^1_{\rm ord}(B_v)|}{|B_v|}
                       \frac{|H^0(B_v)|}{|H^0(A_v)|}
       \le \frac{|Z^1_{\rm ord}(B_v)|}{|B_v|}.
\]
This is the subtraction of the triangular matrix size in
the Euler count.
At a real place the involution on the adjoint has
the form \(X\mapsto -S X^{\mathsf T} S^{-1}\) with \(S\)
symmetric and invertible. Since \(S^{\mathsf T}=S\), putting \(X=SY\) gives 
\[
 -SX^{\mathsf T}S^{-1}=X\quad\Longleftrightarrow\quad -Y^{\mathsf T}=Y.
\]
 This identifies the invariants with the skew matrices and accounts for \(\nu\) at each of the \(d\) real places.
Both sums of local degrees equal \(d\). The coefficient of
\(m\) after the global subtraction is
\[
 \sum_v\big(n(d'_v+1)+\nu(1+d_v)\big)
          +nq-(n+\nu)\#\{v\mid\ell\}-\nu d=nd+nq .
\]
This is the coefficient of the depth \(m\) in the first-order bound for one block. The local-condition comparison uses
the fixed torsion subgroup and the actual
decomposition maps at the finite levels.

For fixed \(m\), integral cocycles and the
relations at the limiting group are read at finite
levels on a set in \(\mathcal U\). Use finite generators and the
stable finite quotients in the common profinite group.
If \(a\) cocycles are independent over \(E\) modulo
boundaries, scale them by powers of \(t\) to make them
integral, and test their linear combinations against
coboundaries at finitely many elements. Such a finite
test suffices by Noetherianity of the module of linear
forms in the coefficients and the boundary matrix.

Smith normal form bounds the loss by a fixed power of
\(t\): if \(z_1,\ldots,z_a\) are these integral cocycles, then for \(m\ge c_0\)
\[
 \sum_j u_jz_j\equiv\text{a boundary}\pmod{t^m}\quad\Longrightarrow\quad(u_1,\ldots,u_a)\in t^{m-c_0}A^a.
\]
Thus a fibre of the map from the \(a\) coefficient tuples modulo \(t^m\) has size at most \(|K|^{a c_0}\), and the image has at least
\[
 |K|^{am}/|K|^{a c_0}=|K|^{a(m-c_0)}
\]
elements. The boundary matrices,
including their skew condition, lift from the finite
quotients. The cardinality estimate gives
\[
        a(m-c_0)\le m(nd+nq)+c \qquad(m\ge c_0).
\]
Dividing by \(m\) and taking the limit proves \(a\le nd+nq\). In this one-block case
the odd matrix module of the two-block comparison
is zero. Its second-order obstruction equations have
no odd entries. Only the first-order dimension bound
enters the cotangent argument.

By completeness and the relative cotangent bound,
\(C/IC\) is a quotient of a power-series ring in
\(D+nq\) variables over the coefficient DVR. Lift a
spanning set of the relative cotangent and use
successive approximation to obtain the surjection.
The power-series ring is a domain of dimension
\(1+D+nq\). Let \(S\) denote this power series ring and \(P\) the inverse image of \(Z/IC\). Its full-dimensional quotient forces
\[
 \operatorname{ht}P+\dim(S/P)\le\dim S=\dim(S/P),
 \qquad \operatorname{ht}P=0.
\]
Since \(S\) is a domain, \(P=0\), and the kernel of \(S\to C/IC\) lies in \(P\). Hence \(C/IC\) itself is a domain and
\(Z=IC\). Every prime containing \(IC\) now lies in
the support of \(P_C\).

\proofstep{Descend the vanishing to the original Hecke quotient.}

The patched calculation supplies support after the auxiliary choices and coefficient extensions. The point on the original Hecke quotient is the one needed in the propagation argument.


Finally, a prime below \(p\) lifts below the chosen
arc prime of \(R_0'\) by flat going down. Its inverse
image in \(R\) contains \(I\) and \(\mathfrak y\).
Lift this prime to \(C\) by faithful flatness of
completion of \(R_{p_R}\). The augmentation support
comparison puts \(J_0\) in its contraction to \(R_0\).
This proves the one-point implication from the
finite-level module and cohomological inputs above.
\label{cor:ordinary-ordinary-patching-1}
\end{proof}


\subsubsection{Vanishing below the formal arc}\label{par:count-two-arc}

The patching statement was made for one determinant link. Applied at each level, it gives an ideal containment below a formal arc on which the marked characters remain distinguishable.

\begin{auxlemma}[Regular elements in the cyclotomic kernel in rank four]
\label{lem:ordinary-four-regular-cosets}
Let \(A=k[[t]]\) with \(k\) finite of odd characteristic and
\(F=\operatorname{Frac}A\). Let \(G\) be a profinite group,
\(H\) a closed normal subgroup, and
\(r:G\to\mathrm{GL}_4(A)\) continuous with \(r|_H\) absolutely
irreducible over \(F\). Assume the residual semisimplification
on \(G\) is a sum of characters \(\chi_1,\ldots,\chi_4\), with
one ratio of order greater than four. Suppose \(s\in H\) has four
eigenvalues whose pairwise ratios have infinite order. For every
open normal subgroup \(U\) of \(H\) one has
\[
  \forall h\in H\quad\exists g\in hU,\qquad
               \operatorname{disc}_X\det(XI-r(g))\ne0 .
\]
\end{auxlemma}
\begin{proof}
\proofstep{Exclude a decomposition into four lines.}
Work over an algebraic closure of \(F\). Write \(C\) for the
Zariski closure of \(r(H)\), and \(C^\circ\) for its identity
component. Both are normalized by \(r(G)\). The restriction of
the four-dimensional space to \(C^\circ\) is semisimple: the sum
of its simple submodules is nonzero and \(r(H)\)-invariant, hence
the whole space. The isotypic components are permuted transitively.
Some positive power of \(r(s)\) lies in \(C^\circ\) and still has
simple spectrum, so no simple constituent occurs twice.
Consequently the possible dimensions of the summands are
\[
             (4),\qquad(2,2),\qquad(1,1,1,1).
\]

In the last case the four lines are permuted by \(G\). Let
\(N\) be the kernel of its permutation action. On these lines
\(r|_N\) acts by characters \(\psi_i\). Their values and inverses
are integral since they are roots of the characteristic
polynomials of elements of \(\mathrm{GL}_4(A)\) and their inverses.
Extend the \(t\)-adic valuation to the algebraic closure and
reduce the four characters there. For every \(h\in N\),
\[
 \prod_i(X-\overline\psi_i(h))=\prod_i(X-\chi_i(h)).
\]

In particular all the reduced characters take values in a common
finite group of roots of unity of order prime to the characteristic.
Let \(C\) be this finite group and \([\,\cdot\,]:C\to\mu_{|C|}(\mathbb C)\)
a fixed multiplicative lift. The equality of the multisets at \(h\) gives
\[
       \sum_i[\overline\psi_i(h)]=\sum_i[\chi_i(h)]\qquad(h\in N).
\] Independence of the
characters of the resulting finite abelian quotient of \(N\) gives
equality of the two multisets of characters on \(N\). The
\(\chi_i|_N\) are invariant under conjugation by \(G\). Since \(G\)
permutes the \(\psi_i\) transitively, all four reductions must
coincide. Thus every ratio \(\chi_i/\chi_j\) factors through the permutation image of \(G\). If \(\sigma\in S_4\) is the image
of a lift of a generator of this cyclic ratio group, then
\[
  |\operatorname{im}(\chi_i/\chi_j)|\mid\operatorname{ord}(\sigma),
          \qquad \operatorname{ord}(\sigma)\in\{1,2,3,4\}.
\] This contradicts the hypothesis on the residual ratio.

\proofstep{Find simple spectrum in every component.}
Let \(hU\) be a coset of an open normal subgroup of \(H\). The
Zariski closure of \(r(U)\) contains \(C^\circ\), since its finitely
many translates cover \(C\). It suffices to show the discriminant is
nonzero at some point of \(r(h)C^\circ\). The identity component of
the closure of the cyclic group generated by \(r(s)\) is a torus.
In an eigenbasis of \(r(s)\), its four coordinate characters are
distinct, since their ratios have infinite order on \(r(s)\). Choose
an integral cocharacter outside the finitely many hyperplanes where
two of these characters agree. It yields in \(C^\circ\)
\[
             D(z)=\operatorname{diag}(z^{d_1},\ldots,z^{d_4})
             \quad(d_i\text{ pairwise distinct}).
\]

First suppose \(r(h)\) preserves each simple summand. We need a point in the open set
\[
  \{X\in r(h)C^\circ:\prod_{i=1}^4 X_{ii}(X^{-1})_{ii}\ne0\}.
\] Each individual condition is a nonempty
open: in its summand the orbit of a coordinate vector spans by
irreducibility, and inversion carries this coset to one with
the same summands. The intersection is nonempty because
\(C^\circ\) is irreducible. Order the four exponents as
\(a<b<c<d\). For the characteristic polynomial of \(XD(z)\) let
\(v_j\) be the lowest \(z\)-degree of the coefficient of
\(Y^{4-j}\), with \(v_j=\infty\) for a zero coefficient. Expansion
in principal minors gives
\[
 v_0=0,\quad v_1=a,\quad v_3=a+b+c,\quad
 v_4=a+b+c+d .
\]

For \(v_3\) the relevant principal minor is a diagonal entry of the
adjugate and is nonzero. If the minor of the two smallest weights
is nonzero, \(v_2=a+b\), and the four Newton slopes are distinct.
Otherwise its point lies strictly above the middle Newton segment:
\[
             v_2\ \ge\ a+c\ >\ a+(b+c)/2=(v_1+v_3)/2 .
\]
That segment has length two and only its endpoints lie on it. Its residual polynomial
is a binomial of degree two with nonzero constant coefficient, hence
has two simple roots in odd characteristic. More explicitly, after
adjoining \(z^{1/2}\), rescale \(Y\) by the middle slope and
divide out the common order. The two nonzero residue roots lift
separately by Hensel's lemma. The outer segments give two further
roots of distinct valuations. In either case the discriminant is
nonzero over the Laurent series field.

In the remaining case \(r(h)\) exchanges the two planes. Each plane
is spanned by two of the chosen eigenvectors. Irreducibility on
each plane permits us to choose \(X\) in the coset with every
entry in both off-diagonal blocks nonzero, by the same open-set
argument. In block coordinates write
\[
 XD(z)=\begin{pmatrix}0&B D_2(z)\\ C' D_1(z)&0\end{pmatrix},
 \qquad
 \det(Y-XD(z))=Y^4-TY^2+P,
\]
where \(T=\operatorname{tr}(B D_2 C'D_1)\) and
\(P=\det(B D_2 C'D_1)\ne0\). If \(a<b\) and \(c<d\) are the
weights of the two planes, the term of degree \(a+c\) in \(T\)
is unique with nonzero coefficient. In particular
\[
 \operatorname{ord}_z T^2=2(a+c)<a+b+c+d=\operatorname{ord}_z P .
\]
The quadratic in \(Y^2\) has distinct nonzero roots, each with
two distinct square roots. Hence the discriminant is again
nonzero on the coset. Zariski density now supplies an actual
element of \(hU\) with nonzero discriminant. Such cosets form a
basis of the profinite topology.
\end{proof}

\begin{proposition}[Ordinary ideals below a separating formal arc]\label{prop:tower-below-arc}
Let \(y:A\to k'[[t]]\) have center at the closed point of \(B\),
kill \(J\), and be absolutely irreducible on the generic fiber
with all the prescribed inertia characters separated. Under the residual tests in
Lemma~\ref{hyp:ordinary-arc-residual-tests}, the containment at each generization is 
\[
 JA\subseteq P\qquad(P\in\operatorname{Spec} A,\ P\subseteq\ker y).
\]

\end{proposition}

\begin{proof}
Let \(y:A\to k'[[t]]\) be centered over the closed point of
\(B\) and kill \(J\). Assume its generic representation
is absolutely irreducible and its ordered characters
separate on the specified open inertia groups. The component
argument will construct such an arc. On \(G_L\) twist its matrices by
\(\kappa\). The resulting representation has the determinant
specified by \(y f_i\), its ordered characters are obtained
by the same twist, and its pairing and ramification
conditions come from the link. Continuity follows from
locality of the map on \(B\) and continuity of the
normalized matrices in their coefficient model.

\proofstep{Irreducibility on the CM field.}

The generic representation remains absolutely irreducible
on \(G_L\). Indeed, by index-two Clifford theory a reducible restriction here would split over an algebraic closure into two conjugate summands exchanged by the other coset. Relative to those summands every such element would have the form
\[
 \begin{pmatrix}0&A\\ B&0\end{pmatrix},\qquad
 \operatorname{tr}\begin{pmatrix}0&A\\ B&0\end{pmatrix}=0 .
\] For \(u\in G_F\setminus G_L\), write the four residual
characters as \(\bar\chi_a\). The reduction of the trace
at \(uh\), \(h\in G_L\), is
\[
           \sum_{a=0}^3\bar\chi_a(u)\bar\chi_a(h).
\]
Every coefficient is nonzero. The restrictions of the
\(\bar\chi_a\) to \(G_L\) are distinct, hence linearly
independent, so this function of \(h\) is nonzero. This contradicts vanishing. The argument can be
carried out after scalar extension to an algebraic closure
of \(k'((t))\).

\proofstep{Cyclotomic restriction and the residual tests.}

Let \(H\subset G_L\) be the full \(\ell\)-adic cyclotomic
kernel. Each pair of arc characters is separated on the
part of inertia in \(H\) that also kills the residual
characters. The tower's inertia subgroups at \(\ell\)
are contained in \(G_L\), and the scalar twist cancels
in a ratio. The ratios on that subgroup take values in \(1+tk'[[t]]\).

For \(u\in tk'[[t]]\setminus\{0\}\) and \(\ell\nmid b\) the
leading term of \((1+u)^b-1\) is \(bu\). Hence for
\(m=\ell^a b>0\) one has
\[
            v_t((1+u)^m-1)=\ell^a v_t(u)<\infty .
\]
Thus the group is torsion-free. A finite union of the kernels of the nontrivial
ratios cannot cover the compact group, since each has
infinite index and measure zero. Thus choose \(s\) outside the union of those kernels.
For any two local eigenvalues \(a_i,a_j\) at \(s\),
\[
 a_i/a_j\ne1 \quad\Longrightarrow\quad
          (a_i/a_j)^m\ne1 \qquad(m>0)
\]
by the valuation formula. Thus all powers of the matrix
at \(s\) have distinct eigenvalues.

\begin{established}[The residual character tests on both fitted links]\label{hyp:ordinary-arc-residual-tests}
On either fitted link let \(\mathcal C=\{\bar\chi_1,\ldots,\bar\chi_4\}\)
as a multiset on \(G_L\). Its quadratic twists and ratios satisfy
\[
 \{\delta:G_L/\ker\bar\epsilon\to\{\pm1\}:\delta\mathcal C=\mathcal C\}
    =\{1\},\qquad
 \max_{i,j}\operatorname{ord}(\bar\chi_i/\bar\chi_j)>4 .
\]
\end{established}
\begin{proof}
Divide all four labels by the first one. If a character \(\delta\)
permutes the labels after twisting, it permutes these normalized
labels, so in particular it is one of them. Indeed, putting \(\xi_i=\bar\chi_i/\bar\chi_1\) gives
\[
          \delta\{\xi_1,\ldots,\xi_4\}=\{\xi_1,\ldots,\xi_4\},
       \qquad \xi_1=1 \quad\Longrightarrow\quad
          \delta\in\{\xi_1,\ldots,\xi_4\}.
\]
In the even-order tower
the elements \(u_0,g_0\in\ker\bar\epsilon\) of
Proposition~\ref{prop:count-ii-starting-field} distinguish every
nontrivial normalized label from one on at least one of these
elements. The marked root at \(u_0\) has order greater than four.
Their restrictions are retained on the deep layer by
Proposition~\ref{prop:count-ii-auxiliary-primes}.

In the odd-order tower, at \(g_0\in\ker\bar\epsilon\) the four
values up to inversion are \(1,-1,z,-z\), where \(z\ne1\) has
odd order, by Proposition~\ref{prop:count-three-marked-elements}.
Only the first label is one there, and \(-z\) has order
\(2\operatorname{ord}(z)>4\). The powers retained in
Proposition~\ref{prop:count-three-fitted-layers} have exponent
prime to these orders. A common scalar twist of the labels affects
neither the ratios nor the permutation test.
\end{proof}

\proofstep{Verify each one-point patching condition.}

To apply patching to this arc one also needs the ordered local conditions.


The ordinary determinant criterion in characteristic \(\ell\) requires four conditions.
The ordinary characters must be distinct and must avoid a
cyclotomic ratio in the indicated order. The representation
must remain absolutely irreducible on \(H\). Elements with
separable characteristic polynomial must be dense in
\(H\). Finally the local cyclotomic intertwining space at
the auxiliary split prime must vanish. The separation on
inertia proves the first condition.

For the second, an
irreducible rank-four representation whose restriction to
\(H\) is reducible has a nontrivial quadratic self-twist:
the cyclotomic quotient is abelian and the number of
conjugate constituents divides four. The element of \(H\)
with distinct eigenvalues rules out a repeated isotypic
constituent. The pro-\(\ell\) kernel of cyclotomic reduction is
in the kernel of a quadratic self-twist, since \(\ell\) is odd.
Thus reduction gives a permutation
\[
               \{\bar\chi_a\}=\{\bar\delta\,\bar\chi_a\}
\]
with \(\bar\delta\) trivial on the mod-\(\ell\) cyclotomic
kernel. The first residual test excludes it.

For the third condition apply
Lemma~\ref{lem:ordinary-four-regular-cosets} with \(G=G_L\) and its
full cyclotomic kernel \(H\). The element \(s\) constructed on inertia
has the required simple spectrum on every positive power. The marked
residual ratio is tested on \(G_L\); twisting all four characters by
\(\kappa\) preserves it. Writing \(H_{\rm reg}\) for its separable-spectrum locus, the result gives
\[
     H_{\rm reg}\cap hU\ne\varnothing
       \qquad(h\in H,\ U\trianglelefteq H\ \text{open}).
\]

This is the density required for the trace tests. For the fourth condition choose the element
of the auxiliary decomposition group on which the four
residual eigenvalues are equal to the same nonzero scalar
and the mod-\(\ell\) cyclotomic value is different from
one. An intertwiner \(U\) at the decomposition element \(g\)
would satisfy \(\operatorname{Ad}(r(g))U=\epsilon_\ell(g)U\).
Multiply by a power of \(t\) so that its reduction is
nonzero. The eigenvalues of \(\operatorname{Ad}(\bar r(g))\) are all
one, since they are ratios of the common residual eigenvalue.
Therefore
\[
   \bigl(\operatorname{Ad}(\bar r(g))-\bar\epsilon_\ell(g)I\bigr)
        \bar U=0\quad\Longrightarrow\quad\bar U=0.
\] This verifies the auxiliary condition
at the actual prime of each link.

Apply Proposition~\ref{prop:ordinary-one-point-patching} to the point
\(y f_i\).  The hypothesis \(y(J)=0\)
puts \(J_i\) in its evaluation kernel. The marked auxiliary
decomposition element also supplies the nontrivial
mod-\(\ell\) cyclotomic value for detection. For a prime \(\mathfrak q\subset\ker y\) in \(A\), the
one-point conclusion on each link is
\[
    f_i^{-1}(\mathfrak q)\subseteq\ker(yf_i),\qquad
    J_i\subseteq f_i^{-1}(\mathfrak q).
\]
Extend the ideals to \(A\) and add the two inclusions to get
\(JA\subseteq\mathfrak q\).

\end{proof}

\subsubsection{Vanishing below the characteristic-zero seed}\label{par:count-two-seed-support}

Propagation also needs a starting component in characteristic zero. At its known automorphic point one must verify the same containment of ordinary ideals under generization as for the formal arc.

\begin{proposition}[Ordinary ideals below the seed]\label{prop:tower-below-seed}
For the characteristic-zero automorphic seed of the fitted tower,
with the shift of Lemma~\ref{hyp:ordinary-seed-shift}
and the residual tests of Lemma~\ref{hyp:ordinary-arc-residual-tests},
if \(x_{\rm seed}\) denotes the seed evaluation, then 
\[
 P\subseteq\ker x_{\rm seed}\quad\Longrightarrow\quad JA\subseteq P\qquad(P\in\operatorname{Spec}A).
\]

The shift preserves the original extended ideals.
\end{proposition}

\begin{proof}
\proofstep{The allowed shift of the ordinary weights.}

\begin{established}[An integral parallel shift at the seed]\label{hyp:ordinary-seed-shift}
Let \(w_0,\ldots,w_3\) be the doubled exponents at a chosen
ordinary seed, in their prescribed order. For the link's
anticyclotomic half-weight character \(\beta_i\), put \(n=|w_0|\)
and twist the scalar by \(\beta_i^n\). Up to finite-order inertia
characters the shifted exponents \(a_j\) satisfy
\[
        a_j=(w_j+n)/2\in\mathbb Z,\qquad
        a_j-a_k=(w_j-w_k)/2 \quad (0\le j,k\le3)
\]
at the chosen half, with opposite shift at its conjugate. The shift is
residually trivial, unramified outside \(\ell\), and trivial
on the prescribed local torsion.
\end{established}
\begin{proof}
Here is the character construction at the split coefficient prime.
Write \(K\) for the marked imaginary quadratic field and
\(\ell=\lambda\bar\lambda\) in \(K\). In local reciprocity units
at \(\lambda,\bar\lambda\), use the principal-unit parts of the
two cyclotomic characters with opposite exponents \(1/2,-1/2\)
in the order of the link. Writing \(\langle\epsilon_w\rangle\) for the principal-unit part, the
unit character in this orientation is
\[
 (x_\lambda,x_{\bar\lambda})\longmapsto
  \langle\epsilon_\lambda(x_\lambda)\rangle^{1/2}
  \langle\epsilon_{\bar\lambda}(x_{\bar\lambda})\rangle^{-1/2}.
\]
Squaring is invertible on principal units since \(\ell\) is odd. This defines an anticyclotomic
character of the local unit ideles, trivial on torsion and on
units at all other finite places. It kills the global units of
\(K\), which are roots of unity.

It therefore defines a character on the image of these units
in the finite idele class group modulo \(K^\times\), with the
complex infinite factor suppressed. For \(U_f=\prod_{v<\infty}O_{K_v}^\times\) the quotient is
\[
      \mathbb A_{K,f}^\times/(K^\times U_f)=\operatorname{Cl}(K).
\] To extend the character, lift generators
of its cyclic factors. If one such lift has order \(h\) modulo
the preceding subgroup, choose an \(h\)-th root of its prescribed
\(h\)-th power value, choosing residue one. These finitely many
choices take values in the principal units of a finite extension.
The extension \(\gamma\) is continuous because its original
domain is open. If \(\kappa_0\) denotes the prescribed character on
the local units \(U\), put
\[
 \kappa=(\gamma/\gamma^c)^{1/2},\qquad
 \kappa^c=\kappa^{-1},\qquad
 \kappa|_U=(\kappa_0^2)^{1/2}=\kappa_0 .
\] Restrict it to \(G_L\) by reciprocity.
This constructs the required companion, with either orientation
of the split pair. Its square has the designated integral
cyclotomic exponent on open inertia.

The two seeds have even differences \(w_j-w_0\), by their displayed
inertia exponents in Propositions~\ref{prop:count-two-integral-induced-seed}
and~\ref{prop:count-three-tensor-endpoint}. Hence
\[
          w_j+n\equiv w_j-w_0\equiv0\pmod2,\qquad
          (w_j+n)/2-(w_k+n)/2=(w_j-w_k)/2 .
\]
This proves integrality and preserves the required gaps (at least
two). At the conjugate embedding the sign of the shift reverses
and the same parity calculation applies.
\end{proof}

\proofstep{Invariance of the Hecke ideals.}

We check invariance of this particular \(J_i\).
Every finite \(\ell\)-power order approximation \(\chi\)
of the companion with these conditions acts by scalar
twist on the ordinary eigensystems. For an idele \(a\in\mathbb A_F^\times\) fixed by conjugation, anticyclotomicity gives
\[
 \chi(a)=\chi(a^c)=\chi(a)^{-1}.
\]
Since \(\ell\) is odd, the order comparison gives
\[
 \operatorname{ord}(\chi(a))\mid\gcd(2,\ell^b)=1
\]
for some \(b\). Hence \(\chi(a)=1\). The map
\[
 a\mapsto a/a^c
\]
and Hilbert 90 descend
\(\chi\) to a character of the norm-one torus over \(F\).
Twisting the unitary forms by its determinant character
keeps the tame coefficients and changes only the torus
character at \(\ell\) and the corresponding Hecke scalars.
The residual localization and the prescribed torsion
characters agree. The finite twist and its inverse therefore
preserve the detecting eigenpoints of this ordinary family.

For the adic assertion reduce \(\beta_i^n\) modulo
\(\mathfrak m_{R_i}^b\), let \(Q_b\) be its finite image,
and tensor the universal determinant by the resulting
character with values in the group basis of \(R_i[Q_b]\).
This is an allowed determinant over the local group ring,
with the same tame relations and the conjugate polarization.
At a detecting field evaluation \(\rho\), adjoin the
characters of \(Q_b\). If
\[
 \operatorname{Tw}_{Q_b}(x)=\sum_q c_q[q]
\]
for \(x\in J_i\),
the finite-twist argument gives
\[
 \sum_q \rho(c_q)\chi(q)=0\quad(\chi\in\widehat Q_b),
 \qquad
 \rho(c_q)=\frac1{|Q_b|}\sum_{\chi\in\widehat Q_b}\chi(q)^{-1}
                             \sum_h\rho(c_h)\chi(h)=0 .
\]
Thus the group-basis coefficients lie in the kernel at
every detecting evaluation.
Such evaluations cut out \(J_i\), by
Paragraph~\ref{par:ordinary-functions}.
Substitute the original scalar values modulo
\(\mathfrak m_{R_i}^b\). Separatedness lets \(b\) tend to
infinity, so the twist preserves \(J_i\). Its inverse
gives equality, hence
\[
                         f_{i,n}(J_i)B=f_i(J_i)B.
\]

The seed is cuspidal after the indicated scalar twist
on \(G_L\), and its polarization and flag after the
shift are those of the specialized determinant point.
The ratio of the shifted scalar character to the seed
automorphic twist is algebraic: the inertia weights have
matching parity, so their difference gives integral
cyclotomic exponents up to finite order. Thus algebraic
twisting preserves the automorphic attachment.

The seed
is unramified at both auxiliary places. Its absolute
irreducibility on the tower's CM cyclotomic subgroup
implies that of the link on the full CM cyclotomic kernel,
by restriction and cancellation of the scalar twist.
The distinct inertial exponents rule out a cyclotomic
ratio on any open subgroup of inertia.
Section~\ref{sec:count-two-seed-specialization} puts the seed
on \(M_i\) with the specialized Frobenius polynomials and
the ordered local characters. Writing \(y_i\) for the shifted evaluation, these support conditions give
\[
 J_i\subseteq\ker y_i.
\]
This uses one common
coefficient ring for the two links, the seed and the
residual chain.

We use the characteristic-zero adjoint comparison to
strengthen this last containment. A dual-number
lift of the specialized determinant that kills the
infinitesimal weight parameters is represented by a
matrix lift, by absolute irreducibility, with its pairing
and ordered ordinary flags. If \(\widetilde r\) denotes the lift, write
\[
 \widetilde r(g)=(1+e c(g))r(g),\qquad e^2=0.
\]
Multiplication gives
\[
 c(gh)=c(g)+\operatorname{Ad}(r(g))c(h).
\]
The ordinary condition at \(\ell\) forces the inertial diagonal variations to vanish, and at the remaining allowed places the cocycle is unramified.

At the
auxiliary place the base representation is unramified
and the residual cyclotomic separation gives the
vanishing of the infinitesimal inertia term. The
gaps between the \(m_{v,i}\) rule out cyclotomic
ratios and give the separated ordinary parameters.
Here one may apply the published adjoint vanishing theorem
directly to the polarized seed. We give the passage to its
local conditions. After a finite local extension killing
the finite inertia factors, the two copies of each graded
character in the dual-number flag form an unramified
self-extension up to twist. They are crystalline. Denote
these two-dimensional graded pieces in their flag order
by \(E_i\). To extend an already crystalline preceding
part \(V\) by \(E_i\), use
\[
 W=\operatorname{Hom}(E_i,V).
\]

Each Hodge jump of \(W\) is a difference of jumps of the
ordered graded characters. The ordered gap makes all jumps of \(W\) strictly negative
and its character exponents at least two. Therefore
\[
           \mathrm{Fil}^0D_{\rm dR}(W)=0,\qquad
                         H^0(W^\vee(1))=0 .
\]
With \(d_v=[L'_v:\Q_\ell]\), local duality and the Euler
characteristic give
\[
 \dim H^1(W)=\dim H^0(W)+d_v\dim W
 =\dim H^0(W)+\dim D_{\rm dR}(W)/\mathrm{Fil}^0
 =\dim H^1_f(W).
\]
The middle dimension is over the coefficient field, with
\(D_{\rm dR}(W)\) viewed by restriction of scalars. The
last equality uses the crystalline dimension formula.
The inclusion \(H^1_f(W)\subseteq H^1(W)\) and the dimension equality give 
\[
 H^1(W)/H^1_f(W)=0.
\]
 Hence every extension represented by one of these classes is crystalline. Induction on the flag now proves crystallinity
of the enlarged representation.

The extension of the trivial
character by the adjoint defined by the tangent class is
crystalline there too. Indeed over \(B_{\rm cris}\) the
enlarged representation splits above the original by
exactness of the crystalline functor on crystalline
representations, giving a splitting in their Hom space.
On \(H^1(B_{\rm cris}\otimes\operatorname{ad})\) for the
original local field the composite of restriction and
corestriction is
\[
                \operatorname{cor}\circ\operatorname{res}
                           =[L'_v:L_v]\operatorname{id}.
\]
The group is a vector space in characteristic zero. Hence the
splitting over \(L'_v\) gives zero image in this cohomology group
over \(L_v\), the required finite condition.
Away from \(\ell\) the inertia subgroup killed by the link
has finite index in inertia outside the auxiliary place.

At that place a possible tame infinitesimal inertia value
would intertwine the matrix of the chosen decomposition
element with its cyclotomic twist, by the action of this
element on tame inertia. Its eigenvalues reduce to a common
scalar and its mod-\(\ell\) cyclotomic value differs from one.
Writing \(U\) for that tame value at a generator and
\(F\) for the matrix of the distinguished element, the
local relation has the form
\[
                   F U=\epsilon_\ell(g)\,U F .
\]
The spectra of \(F\) and \(\epsilon_\ell(g)F\) reduce to
different scalars, so the two spectra are disjoint. The
associated Sylvester operator \(U\mapsto FU-U\epsilon_\ell(g)F\)
is invertible, forcing \(U=0\).

Wild inertia at residue characteristic
different from \(\ell\) has no additive \(\ell\)-adic
quotient, and finite-quotient averaging removes the
remaining inertia class. Polarization extends
the class to the adjoint over \(L^+\) (equivalently,
use invariants for the order-two extension in characteristic
zero). The seed on \(G_L\) comes from a polarizable, regular algebraic
cuspidal representation of \(\mathrm{GL}_n(\mathbb A_L)\). Its
restriction to \(G_{L(\mu_{\ell^\infty})}\) is absolutely irreducible,
as checked above before taking the tangent. Thorne's adjoint Selmer
vanishing theorem therefore applies to its polarized adjoint. The
finite local conditions just proved give, by Theorem~6.1 of \cite{thorne-adjoint},
\[
 [c]\in H^1_f(L^+,\operatorname{ad}r)=0
 .
\]

Thus the characteristic-zero tangent class is a coboundary.
Choose \(Y\) giving the infinitesimal conjugation. Then the lift has
matrix \((1+\delta Y)r(g)(1-\delta Y)\) in the dual numbers.
Conjugating by the inverse gives
\[
     (1-\delta Y)\,r_\delta(g)\,(1+\delta Y)=r(g).
\]

The ordered eigenlines give the same ordered characters.
Thus both the determinant lift and flag-character lifts are
constant.
\proofstep{Tangents and the dimension of the Hecke quotient.}

Invariance preserves the Hecke conditions under the allowed weight shift. The size of the Hecke quotient at the shifted point still has to be compared with the dimension of the deformation component.


This is the vanishing of the relative tangent space
over the weight algebra.
\hfill\mbox{\cite{thorne-adjoint}}\label{cor:ordinary-promodularity-2}

Here is the commutative-algebra implication. At the
characteristic-zero evaluation prime, vanishing of the
relative tangent space gives an injection of absolute
tangent spaces into that of the weight algebra. If
there are
\[
 N=4[L^+:\Q]
\]
weight variables, the embedding
dimension of the localized determinant ring is at
most \(N\). The residue fields are algebraic over the
coefficient fraction field and separable. The evaluation
prime of the weight algebra has height \(N\). In fact
its image is integral over the coefficient DVR. If an
element in the image had absolute value greater than
one, all its conjugates would have the same absolute
value. Its minimal polynomial divided by its constant
term would then have constant term \(1\) and other
coefficients in the maximal ideal of the DVR. Applying
that polynomial to a preimage would give a unit mapping
to zero, a contradiction.

Let \(\mathfrak p_\Lambda\) be the evaluation prime. Its image ring is an integral extension of the coefficient DVR and has dimension one. Catenarity of the regular local weight algebra now gives
\[
 \operatorname{ht}\mathfrak p_\Lambda
 =\dim\Lambda-\dim(\Lambda/\mathfrak p_\Lambda)
 =(N+1)-1=N .
\]
The localization
of the modular quotient has dimension at least \(N\).
Indeed its finiteness and torsion-freeness imply that
each minimal prime contracts to zero over the weight
algebra. Going down from a component through the
evaluation prime bounds its dimension below by the
height \(N\) of the characteristic-zero weight prime.

For the localized link ring \(R'\) and its modular quotient,
the dimensions just proved give
\[
       N\le\dim(R'/J_iR')\le\dim R'\le\operatorname{edim}R'\le N .
\]
The ambient local ring is regular of dimension \(N\).
A nonzero ideal in it would lower dimension, so \(J_iR'=0\).
In particular \(J_i\) belongs to every prime under the
shifted evaluation kernel. Contracting a prime of \(A\)
through the shifted determinant map and then using the
equality of the extended ideals proves the seed assertion
for each of the two original links.
\label{cor:ordinary-promodularity-3}

\end{proof}

\subsubsection{Propagation across components of a completion}
\label{par:count-two-components}

The preceding containments control primes below particular points. To cross an intersection of components, one needs curves through the closed point whose generic points avoid the excluded degeneracy loci.

\begin{auxlemma}[A curve avoiding prescribed closed loci]\label{lem:local-curve-avoidance}
Let \((A,\mathfrak m)\) be a Noetherian local domain and let \(0\ne u\in\mathfrak m\).
There exists \(P\in\operatorname{Spec}A\) satisfying 
\[
 u\notin P,\qquad \dim(A/P)=1.
\]

\end{auxlemma}
\begin{proof}
Choose \(P\) maximal among primes avoiding \(u\), and pass to \(B=A/P\).
Every nonzero prime of \(B\) contains \(u\). If \(0\ne x\in\mathfrak m_B\),
a minimal prime over \((x)\) has height one by the principal ideal theorem.
It contains \(u\), hence is also minimal over \((u)\). Therefore
\[
 \mathfrak m_B\subseteq\bigcup_{\mathfrak q\in\operatorname{Min}(B/(u))}
 \mathfrak q ,
\]
where the primes are viewed in \(B\). The union is finite. Prime avoidance makes \(\mathfrak m_B\) one of these height-one primes.
\end{proof}

\begin{proposition}[Propagation at a permitted vertex]\label{prop:tower-component-propagation}
At a permitted residual vertex with the local flags and size
conditions of the fitted tower, suppose one component of its
normalized completed local deformation algebra contains \(J\).
If \(P_1,\dots,P_h\) are the minimal primes of the normalized completed local algebra at the vertex, this propagates the inclusion to each branch: 
\[
 J\subseteq P_i\quad\text{for every }1\le i\le h.
\]
 The
conclusion descends through the faithfully flat completion.
\end{proposition}

\begin{proof}
\proofstep{The complete local presentation.}

Let \(x:A\to k_\Lambda\) be a permitted residual representation
with its flags, and let \(R\) be its complete local deformation
algebra, with the map \(A\to R\) of the normalized space.
Call a minimal prime of either algebra modular if it contains
the image of \(J\). The local completion is faithfully flat.
Thus each minimal prime of \(A\) under \(\ker x\) lifts to a
minimal prime of \(R\), and modularity of such a lift descends.

The residual presentation over \(\Lambda\) has \(a\) variables
and \(b\) equations with
\[
                      b+2+2d\leq a+2g .
\]
Indeed the adjoint module is the symplectic Lie algebra. The
ordinary Borel subalgebra at each place over \(\ell\) has
dimension six. At a real place use dual bases of the paired Lagrangians \(V_+,V_-\). A centralizing endomorphism is block diagonal. The linearized symplectic equation reads
\[
 \begin{pmatrix}X^t&0\\0&Y^t\end{pmatrix}
 \begin{pmatrix}0&I\\-I&0\end{pmatrix}
 +\begin{pmatrix}0&I\\-I&0\end{pmatrix}
  \begin{pmatrix}X&0\\0&Y\end{pmatrix}=0
 \quad\Longleftrightarrow\quad Y=-X^t .
\]
The free block \(X\) has dimension four. The dual \(H^0\) term vanishes by the exclusion
of the mod-\(\ell\) cyclotomic character from the adjoint root
characters.

The relative Poitou--Tate formula therefore gives
\(h^1_{\rm rel}-h^2_{\rm rel}\geq h^0_{\rm rel}+2d\).
The normalized framed presentation before imposing the local
generator equations has at least
\(h^1_{\rm rel}-h^0_{\rm rel}+2\) tangent variables.
The final two directions come from the symplectic torus
centralizing the diagonal normalizing matrix, with the
global infinitesimal stabilizer subtracted in degree zero.
At most \(h^2_{\rm rel}\) relations arise from the
relative obstruction space, and the two independent diagonal
conditions at each local mark add at most \(2g\).
Combining with the tangent count gives
\[
 a-b
 \ \ge\ (h^1_{\rm rel}-h^0_{\rm rel}+2)-(h^2_{\rm rel}+2g)
 \ \ge\ 2d+2-2g .
\]
This is the stated presentation bound.

\proofstep{Connectedness above the dimension threshold.}

Set \(N=2d+1-2g\). The regular power series ring over
\(\Lambda\) has dimension \(a+1\), so Grothendieck
connectedness connects the minimal primes of \(R\) by
intersections of dimension at least \(a-b\geq N+1\).
Partition the minimal primes into two nonempty sets \(\mathcal P,\mathcal Q\) and let \(U,V\) be the corresponding unions of components. Noetherianity makes the unions finite, so
\[
 U\cap V=\bigcup_{\substack{P\in\mathcal P\\Q\in\mathcal Q}}V(P+Q),
 \qquad
 \dim(U\cap V)=\max_{\substack{P\in\mathcal P\\Q\in\mathcal Q}}\dim V(P+Q).
\]
Connectedness in dimension gives
\[
 \dim(U\cap V)\ge a-b.
\]
A pair attaining the maximum supplies an edge across the partition. It suffices to
propagate modularity across one intersection.
\hfill\mbox{\cite{stacks-connectedness}}\label{cor:ordinary-promodularity-4}

\proofstep{Bound the reducible loci using the permitted supports.}

The constant \(C\) used in the dimension bounds is chosen
on the tower after the auxiliary ramification has been
killed. It bounds \(g\), the global cocycles in the even
root degrees and the global cocycles in all other degrees
which are locally coboundaries at \(\ell\). For an odd
root degree it bounds the excess dimension above \(d\)
of the space of global cocycles. These estimates follow
from Ferrero--Washington on the abelian residual kernel
field and the local Euler characteristic. For a fixed allowed prime \(v\), the open decomposition subgroup \(D_v\) in the tower group \(G\) gives the uniform bound
\[
 \#\{w\mid v\text{ in a layer}\}\le[G:D_v]<\infty.
\]
The even ranks use the trivial action of complex
conjugation, and the odd ranks include one dimension per
real place. The estimates persist under extension of the
residue field, so \(C\) is chosen before the deep layer
and before its seed and coefficient field.

Here is the reducible dimension bound. For the support list abbreviate the four root degrees by
\[
 (a_0,b_0,p_0,t_0)=(w_\alpha,w_\beta,w_\Psi,w_\vartheta).
\]

The diagonal projectors at the normalizing element lie
in the image algebra. Write \(E_{ij}\) for a matrix unit. For a matrix \(X\) in the image algebra the diagonal projectors isolate
\[
 E_{ii}XE_{jj}=X_{ij}E_{ij},\qquad E_{ij}E_{jk}=E_{ik}.
\]
At a field point every nonzero \(X_{ij}\) can be divided out. This proves transitivity of the allowed arrows. They are stable under
symplectic reversal, since taking the symplectic adjoint
of a group element is a scalar times its inverse.
The representation is absolutely irreducible when all
arrows occur, since they then span the full matrix algebra.
There are eight independent off-diagonal arrows under
reversal.

In the other cases their root degrees are
contained in one of the following nine supports:
\[
\begin{gathered}
\{\pm b_0,-a_0,-t_0,-p_0\},\quad
\{\pm b_0,a_0,t_0,p_0\},\\
\{\pm t_0,p_0,b_0,-a_0\},\quad
\{\pm t_0,-p_0,-b_0,a_0\},\\
\{\pm a_0,-t_0,-p_0,-b_0\},\quad
\{\pm a_0,t_0,p_0,b_0\},\\
\{\pm p_0,b_0,-a_0,-t_0\},\quad
\{\pm p_0,-b_0,a_0,t_0\},\quad
\{\pm b_0,\pm t_0\}.
\end{gathered}
\]
One checks the list by taking any missing arrow \(j\to i\)
and using transitivity to exclude each two-step path from
\(j\) to \(i\). Equivalently, enumerate the \(2^8\)
reversal-stable supports and retain those closed under
composition. Setting the entries outside each support to
zero gives the finite union of closed reducible loci.

At a residual vertex the diagonal matrix entries and the
ordinary characters are the marked reductions \(\bar\chi_i\).
The global support is either contained in
\[
 N_P=\{b_0,-a_0,-t_0,-p_0\},\qquad
 N_M=\{-a_0,-b_0,-p_0,-t_0\},
\]
with the additional frame support
\(Q=\{-a_0,-p_0,-t_0\}\) in the \(N_P\) case.
Write \(E(w)\) if a global matrix has a nonzero entry of
degree \(w\). The allowed vertices satisfy at least one of
these four tests.
\begin{enumerate}
\item \(E(-p_0)\) and either \(E(b_0)\) or \(E(-a_0)\).
\item \(E(-a_0)\) and \(E(-b_0)\).
\item The global support is in
\[
 V=\{b_0,-a_0,-t_0\},
\]

      each frame has support in \(Q\) with
      \(n_v(2,0)\ne0\), and \(E(-a_0)\) holds.
\item The same support and frame conditions as in the third
      test hold, together with \(E(b_0)\), \(E(-t_0)\) and
      \(\bar\chi_0/\bar\chi_1\ne1\) as a character on each
      decomposition group at \(\ell\).
\end{enumerate}
These are conditions on the actual vertices. In particular,
every support \(S\) for the tangent calculation already
contains every degree \(w\) required by the corresponding test \(E(w)\).

Let \(x\) be a point on this support and \(x_0\) its limit under the indicated diagonal torus. All points of the punctured torus orbit have equally large tangent spaces, so upper semicontinuity gives
\[
 \dim T_x=\dim T_{D(t)x}\le\dim T_{x_0}\quad(t\ne0).
\]
We bound the last dimension by \(d+20C\) after reducing the global matrices to the four characters. For the first two
tests use the torus exponents \((0,2,1,3)\) on \(N_P\)
and \((0,1,2,3)\) on \(N_M\).

If the exponents are \(b_i\), conjugation by the diagonal \(D(t)=\operatorname{diag}(t^{b_i})\) gives
\[
 (D(t)XD(t)^{-1})_{ij}=t^{b_i-b_j}X_{ij}.
\]
Under the diagonal one-parameter subgroup with exponents \(a_i\), an \((i,j)\) entry is multiplied by
\[
 t^{a_i-a_j}.
\]
Positive differences therefore disappear at \(t=0\), and zero differences preserve the entry. The two paired sums agree, and the
sign on every supported global off-diagonal entry is
positive, so the torus limit of the global matrix is diagonal.
Its frame limit is the identity. The lower odd root directions are then locally
coboundaries by the flag condition. In the nine displayed
supports the first two tests allow precisely rows \(1,5,7\).
Each has at most one upper odd entry. Thus at most one
direction uses the bound \(d+C\) for global cocycles.

For the last two tests use the exponents \((0,1,0,1)\).
The degrees in \(V\) have positive exponent, and precisely
the \(-p_0\) entries of the lower frame survive. The residual
point degenerates to the same diagonal representation with
flags
\[
 1+u_v(e_{20}+e_{31}),
\]
\(u_v\ne0\). The character
of degree \(p_0\) is trivial on the local decomposition
groups. Write \(Z\) for the local infinitesimal matrix
including the changes of frame. The lower-triangular
entries of the ordinary equation give
\[
 Z_{10}=-u_v Z_{12},\qquad Z_{30}=-u_v^2 Z_{12},
 \qquad Z_{21}=2u_v Z_{01}+u_v^2 Z_{03}.
\]
Each equation only couples entries with the same local
character. Hence if one of the first three entries
represents a local coboundary then all three do.

For the
last triple, the same conclusion follows when two are
local coboundaries. The supports containing \(-a_0\)
or both \(b_0,-t_0\) are rows \(1,3,5,6,7,9\)
when the nine sets above are read left to right. Fix the global cocycle at the entry indicated for each row.
\[
 \begin{array}{c|rrrrrr}
 \text{row}&1&3&5&6&7&9\\ \hline
 \text{entry}&(1,2)&(1,2)&(2,1)&(0,1)&(1,2)&(2,1)
 \end{array}
\]
In its kernel at least one entry of the first triple and at least two of the second are zero as global cocycles. The other
odd entries are then locally coboundaries. This leaves
at most one direction with bound \(d+C\).

After fixing this one global cocycle, at most ten
independent global entries remain, each of dimension
at most \(C\). The infinitesimal frame at each of the
\(g\) places is determined by at most six lower entries.
Let \(T_R\) be the tangent space of this fixed support
and \(e\) evaluation at the chosen global cocycle. The
two steps of the calculation give
\[
 \dim \operatorname{im} e\le d+C,\qquad
 \dim\ker e\le10C+6g,\qquad
 \dim T_R\le d+C+10C+6g\le d+20C .
\]
Thus the cotangent generators and the coefficient uniformizer
generate the maximal ideal up to its square. Adding the coefficient uniformizer gives
\[
            \dim R_{\rm supp}\le\operatorname{edim}R_{\rm supp}
                        \le d+20C+1,
\]
where \(R_{\rm supp}\) is the local ring of the fixed-support locus.
\hfill\mbox{\cite{ferrero-washington}}\label{cor:ordinary-promodularity-5}

On the irreducible locus,
passing from the characteristic-polynomial and character
base to the normalized representation adds dimension at
most \(2+4g\). The latter number comes from two torus
coordinates and the four entries of each lower unipotent
symplectic local frame. To see the first contribution, the
trace pairing on the full matrix algebra recovers the
relations among the matrices of group elements from their
traces of products. After making the four eigenvalues at the
normalizing element rational over the characteristic-polynomial
field by an algebraic extension, the matrix algebra splits
over an algebraic closure of that field. Two absolutely
irreducible realizations with these traces are conjugate
there, and the ordered diagonal normalization makes the
conjugation diagonal.

For an invertible diagonal conjugator \(D=\operatorname{diag}(d_0,\ldots,d_3)\), the similitude condition is
\[
 d_0d_3=d_1d_2,\qquad
 \dim\bigl(\{D:d_0d_3=d_1d_2\}/\mathbf G_m\bigr)=4-1-1=2 .
\]
The four entries \((1,0),(2,0),(3,0),(2,1)\) determine a
lower unipotent symplectic frame. Summing the parameters gives
\[
    \operatorname{trdeg}_{\operatorname{Frac}(B/\mathfrak q)}
        \operatorname{Frac}(A/\mathfrak p)\le 2+4g ,
\]
where \(\mathfrak p\) is the point's prime and
\(\mathfrak q=\mathfrak p\cap B\). The two nonflag coordinates
are the torus coordinates modulo scalars. The complete local dimension
formula for the finite type representation chart gives the
same bound on the difference of quotient dimensions at such
a prime.

For a characteristic-\(\ell\) prime of \(B\) containing
\(J\), coincidence of two characters on the prescribed
open inertia group at a single \(v\) gives a dimension
bound \(2d+1-d_v\). Two independent symplectic characters
are determined on all of inertia at a place by their values
at \(d_v\) local units for each character. By local class
field theory the local principal units modulo torsion have
pro-\(\ell\) rank \(d_v\). For \(u=1+x\) with \(x\) in the maximal ideal of a characteristic-\(\ell\) local domain,
\[
 u^{\ell^a}=1+x^{\ell^a},\qquad
 (u^m-1)/(u-1)=1+u+\cdots+u^{m-1}\equiv m\pmod{\mathfrak m}
 \quad (\ell\nmid m).
\]
The second expression denotes the polynomial geometric sum. The two identities exclude both \(\ell\)-power torsion and prime-to-\(\ell\) torsion. Thus these parameters suffice on the prime quotients, including in the ramified layers.
\proofstep{Exclude the inertia coincidence loci.}

Connectedness supplies an intersection of components of sufficiently large dimension. The dimension bounds must also leave points where two marked inertia characters do not coincide.


Use the prescribed coincidence on the subgroup in the
cyclotomic kernel and impose one more equality on an inertia
generator for the cyclotomic direction at \(v\). Coincidence
on a subgroup of finite index in the kernel suffices,
since any remaining value of the ratio would then be a
torsion unit congruent to one.
The two characters then agree on inertia. For symplectic partners the common value \(u\) on an inertia element satisfies
\[
 u^2=1,\qquad u\equiv1\pmod {\mathfrak m},\qquad
 (u-1)(u+1)=0 .
\]
As \(\ell\ne2\), \(u+1\) is a unit, so both characters are one on inertia. Otherwise, killing one character
also kills the other there.

In either case all inertia
characters are killed by the following parameters: \(d_v\)
per place for a first independent character, \(d_v\) at
the other \(g-1\) places for a second, and the one equality
at the marked place. The first property of \(J\) applies after
\(d+(d-d_v)+1=2d-d_v+1\) cuts. The height theorem bounds the
coincidence-locus dimension in \(B\) by this count. For the
irreducible locus in \(R\) add the chart parameters:
\[
                         \dim\le 2d+1-d_v+2+4g .
\]

In the deep layer one has \(d\geq30(C+1)^2\) and
\(d_v\geq30(C+1)\). In particular
\[
                   d+20C+1<N,\qquad
                   2d+1-d_v+2+4g<N .
\]
Suppose \(P,Q\) are adjacent minimal primes of \(R\) and
\(P\) is modular. In their intersection, cut by the
uniformizer and take a closed irreducible subset of
dimension at least \(N\). It exists by the principal ideal
dimension bound. It avoids containment in the reducible
loci and, at an irreducible representation, in any of the
coincidence loci. In the resulting local domain choose a
nonzero element in each of their defining ideals, and
choose also a nonzero element of its maximal ideal.

Write \(t\) for the chosen maximal-ideal element and \(u_i\) for the chosen ideal elements. In this domain
\[
 0\ne t\prod_i u_i\in\mathfrak m.
\]
Lemma~\ref{lem:local-curve-avoidance} gives a one-dimensional quotient in which this product survives. Its generic point therefore avoids all the prescribed closed loci.
Normalization of this complete local one-dimensional
domain gives a formal arc to \(k'[[t]]\) for a finite field
\(k'\). The generic representation of the arc is
irreducible and separates all the required inertia
characters. The arc factors through \(P\), so it kills \(J\), and both
\[
 P,Q\subseteq\ker y.
\]
The second property of \(J\), namely
Proposition~\ref{prop:tower-below-arc}, gives \(J\subseteq Q\). Thus \(Q\) is modular. This proves that
if one component through a permitted residual point is
modular then every component through that point is modular.
\label{cor:ordinary-promodularity-6}

\end{proof}

\subsubsection{Propagation along the residual path}
\label{par:count-two-propagation}

The component argument is local at one residual point. When two successive points admit the required comparison, their common component carries ordinary support to the other endpoint. Iteration reaches the target.

\begin{proposition}[Propagation from the seed to the target]\label{prop:tower-residual-propagation}
Assume the normalized seed and target lie at opposite ends of
a permitted residual path at a common coefficient level. Suppose
every vertex admits the component propagation and the seed
satisfies the containment of Proposition~\ref{prop:tower-below-seed}.
If \(x_{\rm tar}\) is the resulting target evaluation with its original flags, 
\[
 J\subseteq\ker x_{\rm tar}.
\]

The argument allows a reducible characteristic-zero target.
\end{proposition}

\begin{proof}
The residual chains consist of points joined by normalized
matrix and flag families over \(k_\Lambda[t]\) with
endpoints at \(0\) and \(1\). The lifts to \(A\) at the
endpoints are determined by their matrix and flag entries.
Let \(I_{\rm line}\) denote the prime kernel of the family
over \(k_\Lambda[t]\). A minimal prime \(P\) below it
satisfies
\[
        P\subseteq I_{\rm line}\subseteq I_0\cap I_1
\]
for the two endpoint kernels.

If every component through one
endpoint is modular, \(P\) is a modular component through the
other endpoint. Proposition~\ref{prop:tower-component-propagation} applies at the
other endpoint. The seed gives the first modular component by
taking a minimal prime under its characteristic-zero point
kernel and using local vanishing there. Inducting backwards
along the chain shows that all components at the reduction
of the target are modular. The kernel of its integral point
in the small coefficient model lies under that of its
reduction, so any minimal prime below the former is modular.
This proves that the integral point kills \(J\), and the
assertion survives the injection into the target valuation
ring. The argument applies to the same residual point from
which the matrix and its flags were constructed. It uses
no characteristic-zero irreducibility of the target.
\label{cor:ordinary-promodularity-7}

\end{proof}

\subsubsection{The fitted starting fields for the even-order case}
\label{par:count-ii-tower}

To apply propagation in even order, the starting field must retain enough marked Galois elements to separate the ordinary labels. This requirement governs which character kernels one may impose.

\begin{proposition}[Starting field and marked groups]\label{prop:count-ii-starting-field}
Assume the even-order case of
Lemma~\ref{lem:odd-prime-positive-ranges} and use its auxiliary
character from Lemma~\ref{lem:marked-dyadic-character}.
Write \(H_0\) for the subgroup cut out by the finite character conditions in \eqref{eq:count-ii-starting-groups}. The starting fields are specified by 
\[
 G_{F_0}=H_0,\qquad G_K=\ker\lambda,\qquad G_{k_0}=\ker\phi_0.
\]
 The marked normalizing and inertia elements
are retained, and the chosen quadratic extension splits at the
coefficient places used for the ordinary flag.
\end{proposition}

\begin{proof}
\proofstep{The rational characters and the starting field.}

Use the rational character data of
Paragraph~\ref{par:dyadic-residual} and the \(\eta\) of
conductor \(2^n\) there. The order of the prime-to-\(\ell\)
part of \(\xi\) away from \(\ell\) is even in this case.
Write \(\kappa_0\) for the \(\ell\)-primary part of \(\xi\)
away from \(\ell\), and \(\ell^b\) for an exponent killing
it. These factors are taken first on the character group
and then on the two coprime conductor factors. In
particular \(\kappa_0\) is unramified at \(\ell\) and away
from \(S\).

Let \(m_0\) be the order of the prime-to-\(\ell\)
part of \(\xi\). Write \(\xi=\xi_{\ell'}\kappa_0\xi_{\ell,\ell}\), where the last factor is the \(\ell\)-primary component at \(\ell\). On the simultaneous kernel we have
\[
 \kappa_0(g)=1,\qquad \epsilon_\ell(g)=1\Longrightarrow\xi_{\ell,\ell}(g)=1,
 \qquad \xi(g)^{m_0}=\xi_{\ell'}(g)^{m_0}=1 .
\] On inertia at a place away from
\(\ell\), the same implication holds using just
\(\kappa_0=1\). On that part of inertia the residual cyclotomic factor is
trivial, so \(\bar\Psi=1\) implies \(\bar\xi=1\). Reduction is
injective on prime-to-\(\ell\) roots, giving
\[
              \xi^{m_0}=1,\ \bar\xi=1\ \Longrightarrow\ \xi=1 .
\]

Choose three quadratic signs \(\phi_0,\lambda,\lambda'\)
of distinct odd prime conductors \(p_3,p,p'\) away
from \(S\), with parities \(+,-,-\), respectively.
To do so, given a parity \(s\in\{1,-1\}\) and
the finite set of primes \(S'\) already used or
required to split, take a prime
\(P\equiv s\pmod{8\prod_{q\in S'}q}\) outside \(S'\)
by Dirichlet's theorem. For odd \(q\in S'\), quadratic reciprocity gives
\[
 \left(\frac qP\right)
 =\left(\frac{s}{q}\right)(-1)^{((q-1)/2)((P-1)/2)}=1 .
\]

For \(s=-1\) the two signs are each \((-1)^{(q-1)/2}\), and for \(s=1\) each is positive. Also \((2/P)=1\) since \(P\equiv\pm1\pmod8\). The quadratic character
of conductor \(P\) has parity \(s\), is ramified at
\(P\), and is unramified elsewhere at finite places.
Start with \(S'=S\) and add each chosen prime before
the next choice. An inertia element at any of the
three conductors can change the corresponding sign
independently. It fixes the values of \(\xi,\Psi,\omega,\eta,\kappa_0\),
since these are unramified at the three primes.

In \(G_{\Q}\) take the open subgroup
\begin{equation}\label{eq:count-ii-starting-groups}
 H_0=\{g:\bar\Psi(g)=\phi_0(g),\quad
             \lambda(g)=\lambda'(g),\quad\kappa_0(g)=1\}.
\end{equation}
The signs in the first equality are taken in the
residue field. This subgroup is open because the
three conditions are finite character equalities,
and it contains the commutator subgroup. It contains
complex conjugation, since \(\Psi\) is even,
\(\lambda,\lambda'\) are odd and \(\ell\) is odd.
Its index is even since \(\lambda/\lambda'\) has
order two.

Let \(F_0\), \(K\) and \(k_0\) be the fixed
fields of \(H_0\), \(\ker\lambda\) and \(\ker\phi_0\),
respectively. The fixed fields have
\[
 \begin{gathered}
 F_0/\Q\text{ abelian and totally real},\qquad 2\mid[F_0:\Q],\\
 K/\Q\text{ imaginary quadratic},\qquad
 k_0/\Q\text{ real quadratic}.
 \end{gathered}
\]
The chosen congruences make \(K\) split at \(S\) and \(k_0\)
split at \(2,\ell\). To check the ramification of \(L_0=F_0K\) over \(F_0\),
for a rational inertia group \(I_v\) use \(I_v\cap H_0\).
At \(p\), \(\lambda'|_{I_p}=1\) and
\(\lambda=\lambda'\) on \(H_0\); away from \(p\), \(\lambda\)
is already unramified. Thus
\[
                   \lambda|_{I_v\cap H_0}=1
                         \qquad(v\text{ finite}).
\]

At \(p_3\) use \(\bar\Psi|_{I_{p_3}}=1\) because \(p_3\) lies
outside the host conductor. The other places are already unramified
for \(\phi_0\). Therefore
\[
 \phi_0|_{I_v\cap H_0}=1\quad(v\text{ finite}),
\]
so both quadratic composita are unramified over \(F_0\) there. The field \(F_0\) is tamely
ramified at \(\ell\), because on inertia there all
three quadratic signs and \(\kappa_0\) are trivial,
and the residual \(\bar\Psi\) has order prime to
\(\ell\). Inertia of \(F_0\) at \(\ell\) kills the
entire diagonal residual representation since
\(\eta\) is unramified there. Above each prime of
\(F_0\) dividing \(\ell\), choose its extension to
\(L_0\) inducing the original marked place of \(K\).
Exactly one of the two extensions does so, since
\(K\) splits at \(\ell\). Choose decomposition
embeddings representing these places in the fixed
algebraic closure.

The subgroup contains local elements with the
required large values of \(\bar\eta\). At \(\ell\) take
the Frobenius lift \(f_\ell\) with \(\epsilon_\ell(f_\ell)=1\).
Choose \(2^a u\), \(u\) odd, killing \(\xi\), as in
the dyadic calculation. With
\[
 e=2^a u\ell^b,
\]
the chosen lift satisfies
\[
 \epsilon_\ell(f_\ell^e)=1,\qquad
 \xi(f_\ell^e)=1,\qquad \kappa_0(f_\ell^e)=1.
\]
It therefore kills \(\Psi\), and all three auxiliary signs are trivial on this decomposition group. If the \(\bar\eta\)-value had order at most four, injectivity
of reduction would give \(\eta(\ell)^{2^{a+2}u\ell^b}=1\).
The dyadic test then gives
\[
                 2^n<\ell^2 2^{a+2}\le 2^{n_{\rm II}},
\]
contrary to \(n=n_{\rm II}\).

At two use inertia with cyclotomic
residue \(3\bmod 2^n\) for the \(2^n\)-th roots
of unity. It exists by total ramification there.
Raise it to \(2^{c_2(\xi)}\ell^b\) to kill the
host and \(\kappa_0\). If its \(\bar\eta\)-value had order at most four, the same
test with residue \(3\) would give
\[
       2^n<9\cdot2^{c_2(\xi)+2}<2^{c_2(\xi)+6}\le 2^n .
\]
Thus its order is greater than four.
Its \(\omega\)-value is one. Denote this element by
\(u_0\). Both local elements belong to \(G_{L_0}\).
They are available at every place in question by
conjugation, since the groups defining the fields
are normal and the characters are abelian.

We need two further elements of \(G_{L_0}\) with prescribed character values.
Complex conjugation has \(\xi,\Psi,\eta,\omega\)-values
\((\xi(-1),1,-1,-1)\). Set the three auxiliary signs
to \((1,1,1)\) by the inertia adjustments above.
The resulting element \(a_0\) has
\[
 (\bar\eta,\bar\Psi,\bar\omega)=(-1,1,-1).
\]

To construct the second element, recall that the even-order prime-to-\(\ell\)
factor of \(\xi\) off \(\ell\) takes the value \(-1\).
Choose an element with this value and
\(\epsilon_\ell=1\), using the independent cyclotomic
quotients at coprime moduli. Raise it to \(\ell^b\)
to kill \(\kappa_0\) and set the auxiliary signs to
\((-1,1,1)\). The adjustments preserve
\(\epsilon_\ell=1\). If the resulting element \(g\) has \(\bar\eta(g)^2=1\), replace it by \(gu_0\). Indeed
\[
 \bar\eta(gu_0)^2=\bar\eta(u_0)^2\ne1,\qquad
 \bar\Psi(gu_0)=\bar\Psi(g),\qquad \epsilon_\ell(gu_0)=1 ,
\]
using the marked values of \(u_0\) and triviality of the cyclotomic character on inertia at two.

We obtain the marked normalizing element \(g_0\)
with \(\bar\Psi(g_0)=-1\), \(\epsilon_\ell(g_0)=1\) and
\(\bar\eta(g_0)^2\ne1\). The labels on it are the
four distinct values \(-\bar\eta(g_0),-1,\bar\eta(g_0),1\),
in the order allowed by the rational chart. In
particular the quadratic sign stays nontrivial on
\(G_{L_0}\cap\ker\epsilon_\ell\).

The values needed for separation are
\[
 \begin{array}{c|ccc}
       &\bar\eta &\bar\Psi &\bar\omega\\ \hline
 a_0   &-1 &1 &-1\\
 u_0   &\zeta &1 &1\\
 g_0   &z &-1 &1
 \end{array}
 \qquad \operatorname{ord}(\zeta)>4,\quad z^2\ne1.
\]
The order of \(\zeta\) is a power of two. For even \(i\)
the torus character \(\bar\eta^i\bar\Psi^j\) is \(1\) at
\(a_0\), whereas for odd \(i\) it is \(\zeta^i\ne1\) at
\(u_0\). Both cases distinguish the character from \(\bar\omega\). Every root character
is nontrivial, using \(u_0\) for odd first exponent
and \(g_0\) for the two roots \(\pm w_\Psi\).
The first root also has order greater than four.
Suppose a residual quadratic character trivial on
\(\ker\omega\cap G_{L_0}\) permutes the four residual
labels by multiplication. Applied to the trivial
label it must equal one of the labels.

Since \(u_0,g_0\in\ker\omega\), a putative twisting label \(\lambda\) must satisfy
\[
 (\lambda(u_0),\lambda(g_0))=(1,1).
\]
Both labels involving \(\bar\eta\) fail the first equality and the other nontrivial label fails the second. Thus \(\lambda=1\). At each place above \(\ell\) the local
test of nontrivial \(\bar\eta\) with trivial sign and
Teichm\"uller value is given by the Frobenius power
constructed above. The quadratic restriction also
makes \(\eta|_{G_{k_0}}\) totally odd and unramified
at the two degree-one places above \(\ell\). Its
value on their Frobenius has order greater than
four, as does its restriction on inertia at either
norm-two place.

\end{proof}

\subsubsection{Auxiliary inertia and deeper layers}

The starting field supplies the marked elements. Enlarging it to kill auxiliary inertia must still leave enough primes with the same residual Frobenius behavior for the two level comparisons.

\begin{proposition}[Auxiliary primes on the deep layer]\label{prop:count-ii-auxiliary-primes}
After enlarging the even-order starting field to kill the added
inertia, one can choose arbitrarily many distinct rational primes \(p\), outside the fixed bad set, splitting in \(L\). For a Frobenius \(f_p\) the two auxiliary levels use the simultaneous conditions 
\[
 \bar r(f_p)\text{ is scalar},\qquad \bar\epsilon_\ell(f_p)\ne1.
\]
 Both levels use the same scalar twist of \(\bar r\); twisting multiplies the Frobenius matrix by a scalar.
\end{proposition}

\begin{proof}
The marked cup construction of
Paragraph~\ref{par:marked-cup} enlarges \(F_0\) by
the degree-\(\ell\) cyclotomic fields of the added
rational primes. Write \(F_1\) for the resulting
compositum. This is a finite abelian extension of
odd degree over \(F_0\), unramified at \(\ell\).
It kills the new inertia in the residual chain
and preserves the tame property of the base
field at \(\ell\). In every subsequent abelian
\(\ell\)-power extension of \(F_1\), a common
\(\ell\)-power of the marked elements belongs to
the deeper absolute Galois group. Every marked finite value \(z\) used for these comparisons
has order prime to \(\ell\), so
\[
             \operatorname{ord}(z^{\ell^a})=\operatorname{ord}(z)
\]
for all \(a\). The powers retain the required distinctions.

Taking these powers also preserves the specified decomposition and inertia subgroups. The exact
conditions \(\epsilon_\ell=1\) and \(\omega=1\)
on the indicated elements persist too. The local
frames are transported using the selected places
on the same side of \(K\). The normalizing element
at a new layer is the corresponding power of
\(g_0\), so the residual labels stay distinct there.

Use the cohomological bound \(C\) chosen on \(F_1\)
as in the component argument, before extending
the field further. To make the host restriction
have order prime to \(\ell\), take the compositum
of \(F_1\) with the degree-\(\ell^m\) layer of the
cyclotomic \(\mathbf Z_\ell\)-extension of \(\Q\).
Its absolute Galois subgroup is
\[
 G_{F_1}\cap\{g:\epsilon_\ell(g)^{\ell-1}
                        \equiv1\pmod{\ell^{m+1}}\}.
\]
The kernel of \(\epsilon_\ell^{\ell-1}\) on
\(G_{F_1}\) lies in \(\ker\xi^{m_0(\ell-1)}\).
Indeed \(g^{\ell-1}\) has cyclotomic value one
there and \(G_{F_1}\subseteq\ker\kappa_0\). This
latter kernel is open. Writing \(H_m\) for the displayed nested subgroups, their
intersection is \(\ker(\epsilon_\ell^{\ell-1}|_{G_{F_1}})\).
They are compact, so the finite character kernel gives
\[
         H_m\subseteq\ker\xi^{m_0(\ell-1)}
                  \qquad\text{for all sufficiently large }m .
\]
Otherwise the nonempty closed sets of counterexamples would have
nonempty intersection, contrary to the assertion at the limit.

The exponent \(m_0(\ell-1)\) is prime to \(\ell\).
For degree growth choose a positive prime-to-\(\ell\)
power \(s\) of the marked inertia element of
infinite cyclotomic order such that \(s\in G_{F_1}\).
This is possible by tameness. For a prescribed \(B\), none of the finitely many \(s^k\), \(1\le k<B\), lies in the limiting subgroup, since their cyclotomic images have infinite order. The subgroups decrease, so after taking one common depth,
\[
 s^iG_F=s^jG_F\quad\Longrightarrow\quad s^{j-i}\in G_F
 \quad\Longrightarrow\quad i=j\qquad(0\le i,j<B).
\]
Choose \(B=30(C+1)^2\). Since the \(B\) powers have distinct images at all sufficiently
deep levels,
\[
 [F:\Q]\ge[G_{F_1}:G_F]\ge B.
\]
Increase the depth
once for this bound and the finite character condition; the fitted
data are retained. The seed
and the three residual lines can now be formed
at this layer with the transported labels, the
same \(K\), and the real quadratic character
\[
 \phi_0|_{G_F}.
\]


The transported \(a_0\) still has \(\bar\omega=-1\),
so \(L=FK\) does not contain the \(\ell\)-th roots
of unity. Take a power of the transported \(u_0\)
whose \(\bar\eta\)-value is \(-1\) and multiply
it by the transported \(a_0\). This gives an element
of \(G_L\) with trivial residual diagonal characters
and nontrivial mod-\(\ell\) cyclotomic value. Apply
Chebotarev to the finite quotient given by \(L\),
the residual characters and the mod-\(\ell\) cyclotomic
character. This provides arbitrarily many rational
primes split in \(L\) with these Frobenius values.
The auxiliary-prime separation used in the two
ordinary comparisons is therefore available at
distinct primes outside the fixed bad set. A scalar
twist still gives a single residual eigenvalue at
each of these primes.
\label{cor:ordinary-count-ii-tower-1}
\end{proof}


\subsubsection{The even-order induced seed}

The propagation argument needs one known automorphic representation at the chosen levels. In the even-order case it will be induced from a rank-two Hilbert system with two different local weights.

\begin{proposition}[Inducing characters of the even-order seed]\label{prop:count-two-seed-characters}
For the fitted even-order fields of Paragraph~\ref{par:count-ii-tower}, the ordered residual labels
at the deep normalizer prescribe the inducing characters.
The construction selects positive multiples \(k,t\) of \(\ell-1\)
for which the mixed ordinary Hilbert source \(r_2\) satisfies 
\[
 \det r_2=\epsilon_\ell^k\eta,\qquad \bar r_2^{\,\rm ss}\simeq1\oplus\bar\eta,
\]
 with ordered exponents \((k,0)\) and \((k+t,-t)\) at the two places of \(k_0\) over \(\ell\).
\end{proposition}

\begin{proof}
\proofstep{Order the four residual labels.}

Use the fitted fields and characters of
Paragraph~\ref{par:count-ii-tower}. Write \(\tau\) for the
diagonal normalizing element on the deep layer. On
\(G_F\cap G_{k_0}\) the four reduced labels are
\(1,\bar\eta\), each twice. The two
possible orders of the labels exchange the roles of the
two places of \(k_0\) above \(\ell\).

The rank-two input is a cuspidal mixed ordinary Hilbert
representation \(r_2\) over \(k_0\). For \(k,t>0\) divisible by \(\ell-1\), the required data are
\[
   \det r_2=\epsilon_\ell^k\eta,\quad
   \bar r_2^{\,\rm ss}=1\oplus\bar\eta,\quad
   \text{ordered exponents }(k,0)\text{ and }(k+t,-t).
\]
The determinant exponents sum to \(k\) at both places.
Its first ordinary line has the residual character \(1\)
at the first place and \(\bar\eta\) at the second. Away
from \(\ell\) the two inertial eigencharacters are \(1,\eta\)
with the trivial extensions over their common inertia kernel
at the allowed primes. The construction needs irreducibility
on every open subgroup.

Here is the cuspidality test in the production of this
mixed rank-two input. The character \(\eta\) is finite order
and totally odd, has nontrivial finite conductor, and is
unramified at both places above \(\ell\). The weight-one
Eisenstein series with characters \(1,\eta\) at this
conductor has nonconstant adelic coefficients
\[
 a(\mathfrak a)=\sum_{\mathfrak b\mid\mathfrak a} \eta(\mathfrak b),
\]
with \(\eta\) extended by zero on
ideals sharing a prime with the conductor. Its constant
coefficient is independent of the narrow class
\cite[Section 2.2, Proposition 2.1]{dasgupta-darmon-pollack}.

For \(w\mid\ell\) write \(\mathfrak a=w^e\mathfrak c\) with \(w\nmid\mathfrak c\). The divisor sum gives
\[
 a(w^e\mathfrak c)=a(\mathfrak c)\sum_{i=0}^e\eta(w)^i .
\]
The operator \(V_w\) lowers \(e\) by one in this coefficient, giving zero at \(e=0\). At the raised level,
\(U_w V_w E=E\) and
\(U_w E=(1+\eta(w))E-\eta(w)V_w E\) for the starting
form \(E\), by the coefficient formulas. Let \(\alpha_w,\beta_w\) be the two roots \(1,\eta(w)\)
in the desired order. Their sum and product give
\[
 U_w(E-\beta_w V_wE)
   =(1+\eta(w)-\beta_w)E-\eta(w)V_wE
   =\alpha_w(E-\beta_w V_wE).
\]
Thus the stated stabilization has eigenvalue \(\alpha_w\).

Choose \((\alpha_{w_1},\beta_{w_1})=(\eta(w_1),1)\) and
\((\alpha_{w_2},\beta_{w_2})=(1,\eta(w_2))\).
Since the constant term of \(E\) is the same on every
component, \((1-V_{w_1})E\) kills it identically:
\[
     a_0\bigl((1-V_{w_1})E\bigr)=a_0(E)-a_0(E)=0.
\]
The \(V_w\) at different primes commute and permute the
component indices of the constant terms. The resulting
form has integral nonconstant coefficients, first
coefficient one, and the required two ordinary roots.

For weight raising one can use
the total Hilbert Hasse invariant of weight \(\ell-1\).
The quadratic field is unramified at \(\ell\) and both
places are split over \(\Q_\ell\). At a sufficiently small prime-to-\(\ell\) level the invariant
has expansion \(1\) at each standard Tate cusp. Large parallel
powers admit a characteristic-zero lift by the base change
theorem for parallel weight
\cite[Lemmas 2.2 and 2.3]{deo-dimitrov-wiese}.
One may take a power whose weight is at least three and
divisible by \(2(\ell-1)\).

To obtain full level, pull the
lift to a normal principal tame level contained in this
level. For a lift \(H\) at the normal tame level, form the norm over
the finite tame quotient \(C\):
\[
               H_{\rm full}=\prod_{g\in C}gH.
\]
Each \(gH\) reduces to the same Hasse power by naturality
under prime-to-\(\ell\) level change. The product is invariant
under \(C\) and its expansions are integral and congruent to \(1\). This gives the parallel weight-\(k\)
form with \(\ell-1\mid k\) used below. Higher powers allow
any fixed lower bound for \(k\).

\proofstep{Lift the induced seed and its marked ordinary eigenvalue.}

For the simultaneous lift let
\(F\) be the product of the weight-one form and this
weight-\(k\) form, and use a finite coefficient DVR
\(O\) for the two expansions. The normalized first
nonconstant coefficient of \(F\) is one modulo the
maximal ideal. For an ideal \(\mathfrak a\) and a prime
\(w\) away from the level, the coefficient of \(T_w F\)
at \(\mathfrak a\) is
\[
 a_F(w\mathfrak a)+\eta(w)(Nw)^k a_F(\mathfrak a/w),
\]
where the second coefficient is zero if \(w\nmid\mathfrak a\).

At a prime dividing the level the \(U_w\)-operator has
only the first term. The weight-one formula has exponent
zero. For \(w\) off the level, \(w\nmid\ell\) and
\(\ell-1\mid k\) give \((Nw)^k\equiv1\pmod\ell\).
At the \(\ell\)-adic places use the one-term \(U_w\) rule.
Since \(\bar F\) has the expansions of the stabilized
weight-one form, for every \(w\) these rules give
\[
          \overline{T_wF}=\bar\lambda^{(1)}_w\bar F
\]
with \(T_w\) replaced by \(U_w\) at the level.

\proofstep{Use the integral lattice before restriction.}

The eigenvalues determine the characteristic-zero representation. Its reduction along the tower depends on an integral model, which must be fixed before restriction.


For the integral lifting take the coefficient-integral
lattice in the characteristic-zero space generated by
the Hecke translates of \(F\). The nonconstant coefficient
functionals separate forms of this positive weight. A
form with all of them zero would be constant on each
component, and the weight transformation under a lower
unipotent of the level would force that constant to be
zero. Finitely many coefficient functionals suffice by
finite dimension. They embed the lattice into a finite free
\(O\)-module, so the lattice is finite free. It is stable
under all the displayed Hecke operators by their
coefficient formulas.

If \(L\) denotes the coefficient-integral lattice and \(\pi\) a uniformizer, its definition gives
\[
 \{f\in L:a(\mathfrak a,f)\in\pi O\text{ at every coefficient}\}
 =\pi L .
\]
Indeed \(f/\pi\) then has integral coefficients and belongs to the same characteristic-zero Hecke span. Thus coefficient congruences detect reduction. The finite commutative Hecke algebra splits into complete local
factors indexed by the maximal ideals of its reduction. Write \(e\)
for the idempotent of the stabilized system. Its residual summand
contains the vector with nonzero coefficient at the unit ideal, so
\[
       eM/\varpi eM\ne0\quad\Longrightarrow\quad
                    eM\otimes_O\operatorname{Frac}(O)\ne0 .
\]
Here \(M\) is the free Hecke lattice, and \(eM\) is a direct summand.
Any simultaneous eigenvector in that fibre after scalar
extension has the required reductions for all the operators.
In particular its two \(\ell\)-adic eigenvalues are units.

If \(\bar\lambda_1,\bar\lambda_2\) are the two reduced
characters of a possible noncuspidal lift, at good Frobenius
their trace and product give
\[
       (X-\bar\lambda_1)(X-\bar\lambda_2)
                    =(X-1)(X-\bar\eta).
\]
At each good Frobenius one reduced character therefore equals one.
Reciprocity and Chebotarev give \(G=\ker\bar\lambda_1\cup
\ker\bar\lambda_2\) on the full group \(G\). Write \(H_i=\ker\bar\lambda_i\). If \(x\in H_1\setminus H_2\) and \(y\in H_2\setminus H_1\), then
\[
 \bar\lambda_1(xy)=\bar\lambda_1(y)\ne1,\qquad
 \bar\lambda_2(xy)=\bar\lambda_2(x)\ne1 .
\]
This contradicts the union, so the kernels are nested. Thus one kernel is all of \(G\). Therefore the reductions
are globally \(1,\bar\eta\), up to permutation.

At a
place above \(\ell\), the positive power of the residue
norm makes the second term of the ordinary Eisenstein
eigenvalue nonunit when that term occurs. The first
Eisenstein character must consequently give the chosen
unit eigenvalue at both places. It cannot have residual
value \(\bar\eta\) at one and \(1\) at the other, since
\(\bar\eta\) is nontrivial at both. This excludes every
noncuspidal lift. Take the starting parallel weight \(k+1\ge4\), with
\(\ell-1\mid k\), and denote its normalized Hecke eigenvalues
by \(\lambda_w\). At a prime dividing the tame conductor
the stabilized Eisenstein eigenvalue is \(1\), so the higher
weight eigenvalue there is also a unit.
\label{cor:ordinary-promodularity-8}

\end{proof}

\subsubsection{Changing one Hilbert weight}
\label{par:mixed-hilbert}

The two places over the residual prime have different roles in the induction. Changing the Hilbert weight at one while retaining both ordinary lines produces the required local orders.

\begin{proposition}[Simultaneous cuspidal lift at a mixed weight]\label{prop:mixed-hilbert-lift}
For the marked Eisenstein data over the real quadratic field
\(k_0\), the ordinary specialization below lifts the residual
system with one weight changed, retaining the specified unit-root
ordinary lines. Let \(t_w=t\) at the selected place and \(t_w=0\) at the other. The holomorphic weights are \(k+1+2t_w\). The original and lifted eigenvalues obey 
\[
 \ell^{t_w}\lambda^{\prime}_w\equiv\lambda_w\pmod{\mathfrak m}\quad(w\mid\ell),\qquad \lambda^{\prime}_w\equiv\lambda_w\pmod{\mathfrak m}\quad(w\nmid\ell),
\]
 after finite coefficient extension.
\end{proposition}

\begin{proof}
Let \(P\) consist of the place of \(k_0\) above \(\ell\)
at which the Eisenstein unit root is the Hecke value of
\(\eta\). At the other place it is \(1\). Write \(c(w)\)
for the conductor exponent of the Hecke character associated
to \(\eta\) away from \(\ell\), and put \(c(w)=1\) at both
places above \(\ell\). We need a cuspidal system of the
same central character and level with holomorphic weights
\(k+1+2t_w\), where \(t_w=t\) on \(P\) and \(t_w=0\)
at the other embedding. The comparison required of its
eigenvalues is
\[
 \ell^{t_w}\lambda'_w\equiv\lambda_w\pmod{\mathfrak m}
       \quad(w\mid\ell),\qquad
 \lambda'_w\equiv\lambda_w\pmod{\mathfrak m}
       \quad(w\nmid\ell).
\]
Here and below reduction of a normalized eigenvalue uses a
finite extension of \(\Q_\ell\) containing it.

We give a proof using the definite quaternion algebra available
over this real quadratic field. Its two ramified places are
the two real places, so all finite Hecke operators keep their
local double cosets. Global Jacquet--Langlands transfers a
cuspidal representation that is discrete series at both real
places to this algebra. In the opposite direction, an
automorphic representation of the quaternion algebra of
dimension greater than one transfers back with the same finite
components. At a real place the compact type for weight
\(h\ge2\) restricts to \(\operatorname{Sym}^{h-2}\) on the
norm-one group. The central character fixes its extension.
\hfill\mbox{\cite{jacquet-langlands}}

Put \(d_0=k-1\). At the two split \(\ell\)-adic places use
the integral coefficient module
\[
 V_{\mathbf t}=\bigotimes_{w\mid\ell}
       \bigl(\operatorname{Sym}^{d_0+2t_w}(O^2)\otimes\det^{-t_w}\bigr).
\]
For a scalar matrix \(zI\), the two actions on one tensor factor multiply to
\[
 z^{d_0+2t_w}\det(zI)^{-t_w}
 =z^{d_0+2t_w}(z^2)^{-t_w}=z^{d_0}.
\]
Thus the central exponent is independent of the shift at \(w\). Use a sufficiently large finite
coefficient DVR \(O\). We describe the
Hecke normalization. If \(D\) denotes the quaternion algebra,
write \(x\in D^\times(\mathbb A_f)\) and use finite level
\(K_c=\prod_w K_1(w^{c(w)})\) under the finite splittings.
Quaternionic forms are functions invariant on the left by
\(D^\times(k_0)\), of the given central character with its
algebraic normalization, and satisfying
\[
 f(xu)=\rho_{\mathbf t}(u_\ell)^{-1}f(x)
\]
for \(u\in K_c\).
Here \(\rho_{\mathbf t}\) is the displayed coefficient action.
For a double coset with right cosets \(a_iK_c\) the operator
is \(\sum_i\rho_{\mathbf t}(a_{i,\ell})f(xa_i)\).

To see its scalar normalization, start in the unitary
normalization with the compact type at infinity, write
the form as a vector-valued function on the finite ideles,
and transfer the algebraic action on its values to \(\ell\)
using the fixed embeddings. Multiply its value at \(x\) by
\[
 |\operatorname{nrd}x|_f^{-d_0/2}\,
                \rho_{\mathbf t}(x_\ell)^{-1}.
\]
For a rational quaternion the real normalization of its
algebraic coefficient action contains the factor
\(|N_{k_0/\Q}\operatorname{nrd}x|^{-d_0/2}\). The product
formula thus gives the asserted left invariance.

For a finite central idele \(z\), apart from the finite-order
character, the factor is
\[
             |z|_f^{-d_0}\prod_{w\mid\ell}z_w^{-d_0}.
\]
Because \(\ell\) splits in \(k_0\), the valuation of the
central multiplier is
\[
  d_0\sum_{w\mid\ell}v_\ell(z_w)
                       -d_0\sum_{w\mid\ell}v_\ell(z_w)=0 .
\]
Both coefficient modules have this same central factor. The scalar multiplying a
degree-one Hecke double coset relative to its unitary action
is \((Nw)^{d_0/2}\), for both weights. At \(w\mid\ell\) we
further multiply by \(\ell^{t_w}\) to obtain an integral
operator. Indeed its representatives there have the form
\(\left(\begin{smallmatrix}\ell&b\\0&1\end{smallmatrix}\right)\),
\(b\bmod\ell\). For such a representative the normalized matrix on the coefficient module is
\[
 \ell^{t_w}\det(a_i)^{-t_w}\operatorname{Sym}^{d_0+2t_w}(a_i)
 =\operatorname{Sym}^{d_0+2t_w}(a_i) .
\]
It preserves the integral symmetric power because \(a_i\) has integral entries. Away from \(\ell\) the representatives may
be taken to have identity \(\ell\)-components.

The double quotient by the center and \(K_c\) is finite by
finiteness of the ideal classes of a definite quaternion
order and finite index of the level. Evaluation at its
representatives gives a direct sum of coefficient modules
fixed by their projective stabilizers, with the indicated
central factors. The stabilizers are finite: the rational
projective group is discrete adelically, the archimedean
group is compact, and the finite level is compact modulo
the center.

At the ramified place of \(\eta\) of norm two
the positive level forces the local projective stabilizer
into a conjugate of the image of \(K_1(w^{c(w)})\).
Its reduction is contained in the order-two unipotent subgroup
\[
 \left\{\begin{pmatrix}1&a\\0&1\end{pmatrix}:a\in\mathbf F_2\right\},
\]
and its congruence kernel is pro-\(2\). The rational
projective group maps injectively to the local projective
group, since a quaternion which becomes scalar there was
already scalar. For the reduction map on a finite stabilizer \(H_x\), both the kernel and the image have 2-power order. Hence
\[
 |H_x|=|\ker(\mathrm{red}|_{H_x})|\,|\mathrm{red}(H_x)|
\]
is a power of two.
\[
 |H_x|=|\ker(\mathrm{red}|_{H_x})|\,
       |\mathrm{red}(H_x)| .
\]
The kernel is contained in a pro-2 group, and the image lies in the upper unipotent group over \(\mathbf F_2\).

For a projective stabilizer \(H_x\) at a coset \(x\),
the \(O\)-linear projector on its fixed values is
\(e_x=|H_x|^{-1}\sum_{h\in H_x}\rho_x(h)\), including
the fixed central multipliers in the action \(\rho_x\).
The order of \(H_x\) is a power of two and \(\ell\) is odd,
so \(e_x\) is integral. The splitting \(V=e_xV\oplus(1-e_x)V\) reduces to the analogous splitting over the residue field. For an equivariant surjection \(u:V\to W\), if \(w=e_xw=u(v)\), then
\[
 w=e_xu(v)=u(e_x v).
\]
This proves surjectivity on the two fixed submodules. Summing over the finite coset set
proves both assertions for forms.

In each symmetric power let the last-coordinate quotient
be given by the coefficient of the pure power of the
second basis vector. Modulo \(\mathfrak m\) the matrices
of \(K_1(w)\) are upper triangular, so this is a quotient
of their action. For \(t\) a positive multiple of \(\ell-1\)
the determinant factor on the units is trivial after
reduction. The two weights induce the same character on
the last-coordinate quotient, since all diagonal reductions
lie in \(\mathbf F_\ell^\times\). The normalized operators
at \(w\mid\ell\) also act by the same local quotient map:
their representatives have lower right entry one and the determinant factor is
\[
 \ell^{t_w}\det(a_i)^{-t_w}=1.
\]

Tensor the two quotient maps. Write \(Q\) for the product of last-coordinate maps
on the residue modules of forms. In either weight,
\[
                   \ker Q\subseteq\ker(U_{w_1}U_{w_2}) .
\]
This inclusion is at coefficient reduction modulo \(\mathfrak m\)
and uses the two normalized operators. Indeed for a local representative \(a_i\), modulo \(\ell\)
its matrix on the standard space is
\(\begin{pmatrix}0&\bar b\\0&1\end{pmatrix}\).
Write \(D_w=d_0+2t_w\), the exponent in the coefficient
module \(V_{\mathbf t}\). On its symmetric power the
reduced matrix acts by
\[
 e_1^a e_2^{D_w-a}\longmapsto
 \begin{cases}
 0,&a>0,\\
 (\bar b e_1+e_2)^{D_w},&a=0 .
 \end{cases}
\]
The determinant normalization contributes one. Thus after
taking both factors, every tensor monomial in the kernel
of the two last-coordinate quotients is killed. The
level preserves the quotient and its kernel, so the
calculation holds for all double-coset representatives.

Normalize an eigenvector in the starting weight to be
primitive in the integral lattice of forms. Its reduction \(\bar f\) is nonzero and satisfies
\[
 U_{w_1}U_{w_2}\bar f
          =(\bar a_{w_1}\bar a_{w_2})\bar f\ne0
\]
because both eigenvalues are units. The inclusion just proved
therefore gives \(Q(\bar f)\ne0\). The quotient of the mixed-weight lattice surjects onto the
same residual Hecke module. If \(\mathfrak m\) is the
character's maximal ideal, this gives
\[
                    (M_{\mathrm{mixed}})_\mathfrak m\ne0 .
\]
Thus the residual character is in its support.

Let \(\mathbb T\) be the finite torsion-free algebra generated by the commuting operators on the lattice, and choose a minimal prime \(P\) below the residual maximal ideal. A uniformizer \(\pi\) acts injectively on \(\mathbb T\), whereas each element of a minimal prime is a zero divisor. Hence
\[
 \pi\notin P,\qquad (\mathbb T/P)[1/\pi]\ne0 .
\]
In the normalization of \(\mathbb T/P\), choose a prime above
\(\mathfrak m/P\) by lying over. Completing at this prime yields
a characteristic-zero character over a finite coefficient extension
whose reduction is the prescribed residual character. 

This character occurs on the lattice after inverting \(\ell\) and extending
scalars, by faithfulness and simultaneous triangularization.
The degree \(d_0>0\) of the compact type at the unchanged
embedding rules out all one-dimensional quaternionic
representations. Transfer back gives the claimed cuspidal
Hilbert eigensystem. The two normalized eigenvalues above \(\ell\) reduce to the
prescribed nonzero elements \(\bar a_{w_1},\bar a_{w_2}\), so both
ordinary lines survive. The weight can be chosen with
\[
                    t\in(\ell-1)\mathbb Z_{>0},\qquad t\ge2
\]
as large as required for the subsequent local construction.
\label{cor:ordinary-mixed-hilbert-1}

\end{proof}

\subsubsection{Local and global images of the rank-two seed}

An ordinary eigenform with the desired weights supplies a Galois representation. Using that representation after extending the base field requires control of its image on open subgroups.

\begin{proposition}[Absolute irreducibility of the mixed-weight seed on open restrictions]\label{prop:mixed-hilbert-restriction}
The resulting cuspidal rank-two system has the local images used
for the ordinary seed. In the chosen residual situation, with
\(\bar\eta^2\ne1\), let \(V\) be the two-dimensional space of \(r_2\) over \(\overline{\mathbb Q}_\ell\). Absolute irreducibility on open restriction means 
\[
 \forall U\le_{\rm open}G_{k_0},\qquad \{W\subseteq V:\dim W=1,\ r_2(U)W=W\}=\varnothing.
\]

\end{proposition}

\begin{proof}
Apply local--global compatibility for a cuspidal Hilbert
eigenform of discrete-series weights at least two, all of
the same parity and with algebraic central normalization.
The mixed weights just constructed have these properties.
At \(w\nmid\ell\), newform theory and the fixed vector at
\(K_1(w^{c(w)})\) give the upper bound for the local parameter
\(W_w\):
\[
               c(w)=a(\det W_w)\le a(W_w)\le c(w).
\]
The middle inequality follows by comparing invariant
codimensions over each ramification subgroup.

For the finite inertial part of the
Weil--Deligne parameter, the codimension of invariants on
each ramification subgroup is at least that for the
determinant, and monodromy can only increase the
conductor. At positive conductor equality in the inertia
term forces a one-dimensional fixed space and distinct
inertial characters \(1,\eta\). In the two distinct inertia lines the commuting operator \(N\) is diagonal. Consequently
\[
 N=\operatorname{diag}(n_1,n_2),\qquad N^2=0
 \quad\Longrightarrow\quad n_1=n_2=0 .
\] At zero conductor
the parameter is unramified. This gives the full inertia
action away from \(\ell\), in particular at the norm-two
place.
\hfill\mbox{\cite{skinner-hilbert}}\label{cor:ordinary-mixed-hilbert-2}

At the two split places the normalized ordinary eigenvalues
are units. In the normalization before the final cyclotomic
twist the determinant is \(\epsilon_\ell^{-k}\eta\).
The two exponents are
\[
 m_0=t_w,\qquad m_1=-k-t_w,\qquad m_0+m_1=-k.
\]
A nonzero \(U_w\)-eigenvalue
at \(K_1(w)\) level supplies a Borel Jacquet exponent.
With arithmetic Weil degree its normalized Weil character is
\(s\mapsto(\lambda'_w)^{-\deg s}\). Indeed the translation
by the uniformizer on the Jacquet module gives the \(U_w\)
eigenvalue after the Borel modulus factor.

Passing to the
algebraic Weil parameter cancels this factor. The second
Weil character is unramified too since the determinant is
unramified on the Weil inertia group. The Hodge filtration
jumps are \(-m_0,-m_1\). The normalized unit reading and the determinant give the
Newton slopes
\[
        v_\ell(\lambda'_w)=-m_0,\qquad
        k-v_\ell(\lambda'_w)=-m_1 .
\]
Here crystalline Frobenius is inverse to arithmetic Weil
Frobenius in the stated convention. On the
potentially semistable module the smaller-slope line is
stable under monodromy, which lowers slope by one. It is
stable under descent, since each slope occurs once. Its
induced Hodge number is at least the smaller jump. Weak
admissibility gives the reverse bound by the Newton number,
so the line is weakly admissible. The equivalence for
semistable representations and descent give a Galois line.
Its character \(\kappa_0\) has inertia \(\epsilon_\ell^{m_0}\)
by the rank-one comparison.

On an arithmetic Frobenius
\(\sigma\) the same comparison gives
\[
     (\kappa_0\epsilon_\ell^{-m_0})(\sigma)^{-1}
                 = \ell^{m_0}\lambda'_w .
\]
The other character follows from the determinant. The two reductions of
the first character are \(\bar\eta\) on \(P\) and \(1\)
on its complement, by the weight congruence and local
reciprocity. Twisting by \(\epsilon_\ell^k\) leaves the reductions fixed
since \(\ell-1\mid k\). The determinant becomes
\(\epsilon_\ell^k\eta\) and the ordered exponents become
\[
       (m_0+k,m_1+k)=(k+t_w,-t_w).
\]
These are \((k,0)\) and \((k+t,-t)\) at the respective places.
Away from the level and \(\ell\), comparison of the Hecke
polynomials with the Eisenstein system gives residual trace
\(1+\bar\eta\) and determinant \(\bar\eta\). In particular at good Frobenius elements \(g\) the residual polynomial is
\[
 X^2-(1+\bar\eta(g))X+\bar\eta(g)=(X-1)(X-\bar\eta(g)).
\]
Chebotarev and Brauer--Nesbitt identify the residual semisimplification with \(1\oplus\bar\eta\).
\hfill\mbox{\cite{skinner-hilbert,colmez-fontaine}}\label{cor:ordinary-mixed-hilbert-3}

For irreducibility one can also use the unequal gaps at the
two embeddings. If the representation were reducible, the
two constituents would be global \(\ell\)-adic characters,
locally algebraic at both places by the flags and of finite
inertial image away from \(\ell\). In a totally real field
the cyclotomic exponents of such a character are parallel.
Indeed, restrict by reciprocity to totally positive units in
the kernels of all finite inertial factors. A finite-index
subgroup of the totally positive units remains, and the
product of its embedded values to the indicated integer
powers is one.

For the finite-index group of totally positive units,
transfer their algebraic identities to \(\R\). The local
exponents \(a_\sigma\) then satisfy
\[
          \sum_\sigma a_\sigma\log|\sigma(\epsilon)|=0
                        \qquad(\epsilon\text{ in that group}).
\]
Their logarithms span the trace-zero hyperplane by the unit
theorem, so all \(a_\sigma\) are equal. The possible exponents \(k,0\) at one place and
\(k+t,-t\) at the other are disjoint.

Finally let \(H\) be the normal core of any open subgroup.
An inertia element at the unchanged place has infinite
cyclotomic image and a positive power belongs to \(H\).
The two eigenvalues of this power are distinct. If \(H\)
had an invariant line, all such lines would be among
the two eigenlines of that element, and the full group
would permute them. Global irreducibility forces a
transitive action on two lines. The character \(\nu\)
of this permutation is a nontrivial quadratic self-twist:
in the basis of the two lines, conjugation by
\(\operatorname{diag}(1,-1)\) twists the representation by
\(\nu\). Reducing characteristic polynomials gives
\[
 \{1,\bar\eta\}=\{\bar\nu,\bar\nu\bar\eta\}
\]
as
multisets of characters. The nontrivial quadratic \(\nu\) has nontrivial reduction since \(\ell\ne2\). The two entries must therefore be exchanged:
\[
 \bar\nu=\bar\eta,\qquad \bar\nu\bar\eta=1
 \quad\Longrightarrow\quad\bar\eta^2=1,
\]
contrary to the order of the chosen \(\bar\eta\). Thus every open restriction
is absolutely irreducible.
\end{proof}


\subsubsection{The induced representation on the fitted CM layer}

Induction combines the mixed-weight representation and its conjugate into dimension four. Choosing an integral frame is essential because the residual extension class and the order of the diagonal labels enter propagation.

\begin{proposition}[The integral four-dimensional seed after induction]\label{prop:count-two-integral-induced-seed}
Take the mixed-weight Hilbert system \(r_2\) just constructed
and the normalizer \(\tau\) on the fitted layer. There is an
integral symplectic induced frame with the specified nonsplit
rank-two reduction and the four ordered residual labels.
The scalar correction \(\delta\equiv1\) gives the required
similitude. The induced system is absolutely irreducible on
\(G_L\cap\ker\epsilon_\ell\); after removing \(\delta\) its restriction to \(G_L\) has the cuspidal attachment with ordered exponents 
\[
 (k+t,k,0,-t)
\]
 in the untwisted convention.
\end{proposition}

\begin{proof}
\proofstep{Choose the nonsplit lattice in the marked eigenbasis.}

Choose a stable lattice for \(r_2\) over \(G_{Fk_0}\)
in which \(r_2(\tau^2)\) is diagonal and the reduction
is a nonsplit extension of its two distinct residual
characters. One can arrange the upper triangular
orientation by Ribet's lattice argument. The eigenprojectors at \(\tau^2\) preserve every stable
lattice, since its two eigenvalues differ by a unit. If \(b_g,c_h\)
are the upper and lower entries, the residual trace comparison gives
\[
                        b_g c_h\equiv0\pmod\varpi
                        \qquad(g,h\in G_{Fk_0}).
\]
Rescale the eigenvectors relatively so the smallest valuation
of the \(b_g\) is zero. The products and the unchanged diagonal
then give strictly positive valuation for the \(c_h\).
This gives the nonsplit reduction and retains the
diagonal residual characters.

\proofstep{Induce in an integral symplectic basis.}

Use \(\mu=\epsilon_\ell^k\eta\phi_0\) on \(G_\Q\).
If the two diagonal entries at \(\tau^2\) are \(a^2,b^2\),
choose the signs of \(a,b\) to agree with the residual
labels at \(\tau\), and choose \(ab=\mu(\tau)\). This is
possible since the determinant at \(\tau^2\) is
\(\mu(\tau)^2\) and \(\ell\) is odd. On the two cosets of
\(G_{k_0}\) the induced matrix of \(\tau\) is
\(\left(\begin{smallmatrix}0&\operatorname{diag}(a^2,b^2)\\
I_2&0\end{smallmatrix}\right)\). The change of basis
\[
 P=\begin{pmatrix}
 0&a&0&a\\
 b&0&-b&0\\
 0&-1&0&1\\
 1&0&1&0
 \end{pmatrix}
\]
is integral with inverse integral, and diagonalizes the
operator. If \(R_\tau\) is the displayed two-coset matrix, then
\[
                       R_\tau P=P\operatorname{diag}(b,-a,-b,a).
\]

For \(J_2=\left(\begin{smallmatrix}0&1\\-1&0\end{smallmatrix}\right)\),
the pairing in the induced basis is \(\operatorname{diag}(J_2,abJ_2)\)
and has similitude \(\mu\). Indeed it is
preserved up to \(\mu(h)\) on each block for
\(h\in G_{k_0}\), and the matrix of \(\tau\) exchanges
the blocks with the factor \(ab\). In the new basis the pairing satisfies
\[
       P^{T}\operatorname{diag}(J_2,abJ_2)P=-2abJ,
\]
where \(J\) is the standard alternating form pairing \(0,3\)
and \(1,2\) in this order.
On \(G_F\) the ratio \(\vartheta/\mu\) reduces to one,
so it has a continuous square root \(\delta\) reducing
to one, obtained from the inverse of squaring on the
group of units congruent to one. Twisting the induced
representation by \(\delta\) gives the required
similitude without changing its reduction.

\proofstep{Compare the two constituents on the cyclotomic kernel.}

To apply cuspidal induction the two conjugate constituents must remain distinguishable after restriction. Their behavior on the cyclotomic kernel provides this distinction.


For irreducibility put
\[
 U=G_L\cap\ker\epsilon_\ell.
\]

The restriction of \(r_2\) to \(U\cap G_{k_0}\) is
irreducible. Indeed this is a normal subgroup with
abelian quotient in an open subgroup of \(G_{k_0}\).
If it has one or two invariant lines, an open subgroup
of the larger group fixes a line. If it preserves at least three lines, its image is scalar. In a basis for the first two lines a diagonal matrix preserving a third has equal diagonal entries. Since the quotient is abelian, for elements \(g,h\) of the larger group
\[
 r_2([g,h])=cI,\qquad c^2=\det r_2([g,h])=1 .
\] On a
sufficiently small open subgroup they are the identity,
and a commuting set of \(2\)-by-\(2\) matrices has a
common eigenline over the algebraic closure. Either
possibility contradicts irreducibility on open subgroups.

If the constituents on \(U\cap G_{k_0}\) were equivalent,
an intertwiner would be equivariant on \(G_L\cap G_{k_0}\)
up to a scalar character. Schur's lemma gives this scalar and
equality of determinants makes its square one. At a power of the marked inertia element
with infinite cyclotomic image, square the element to
remove the sign. If the cyclotomic value of the squared
element is \(u\), the characteristic polynomials would
have the two root multisets
\(\{u^k,1\}\) and \(\{u^{k+t},u^{-t}\}\).
Evaluation of the two characteristic polynomials at \(1\) would now give
\[
 0=(1-u^{k+t})(1-u^{-t}).
\]
Both factors are nonzero since \(k+t,t>0\) and \(u\) has infinite order. The constituents are inequivalent, and \(U\) meets
the other coset, so induction is irreducible on \(U\).

\proofstep{Attach the cuspidal induction at the separating weights.}

Over \(L\) the untwisted induction is the automorphic
induction from \(Lk_0\) of the base change of the
rank-two Hilbert system. Solvable base change and
quadratic automorphic induction give the Frobenius
polynomials of the induced representation. Cuspidality
uses the irreducibility of the rank-two restriction
and inequivalence of the conjugate systems. Their weights at the two exchanged places also distinguish them.
The increasing gaps may be chosen with \(k,t\geq2\),
giving the four ordered exponents \(k+t,k,0,-t\).
The character used to recognize the automorphic seed
is \(\delta^{-1}\) on \(G_L\), which cancels the scalar
twist. Indeed
\[
 \delta^2=\epsilon_\ell^{j+1-k}\omega^{-1}\xi\phi_0^{-1}
\]

on \(G_F\), where \(\xi\) is the finite host character. Put
\(b=j+1-k\). Thus the doubled modular weights are
\[
         2(k+t,k,0,-t)+(b,b,b,b).
\]
Their successive gaps are \(2t,2k,2t\), each at least four, and the parity and finite-order
condition after removal of the cyclotomic factor are
the ones used for ordinary classicality.
\hfill\mbox{\cite{arthur-clozel,clozel-rajan}}\label{cor:ordinary-promodularity-9}

\end{proof}

\subsubsection{The integral flag of the induced seed}
\label{par:induced-flag}

The induced representation is integral, but its ordinary flags must reduce in the marked order. We verify this in the actual basis on each of the two local blocks.

\begin{lemma}[Integral ordinary flags on the two induced local blocks]\label{lem:count-two-local-induced-flags}
At each marked split decomposition group of
Proposition~\ref{prop:count-two-integral-induced-seed}, use its
integral basis \(P\) and the chosen ordered local lines. The
matrix \(N\) below is integral symplectic and its flag is
ordinary in the specified order. It obeys the residual grading condition and 
\[
 N_{20}=\pm1,\qquad N_{21}\equiv0\pmod{\mathfrak m},
\]
 where \(\mathfrak m\) is the maximal ideal of the coefficient ring.
\end{lemma}

\begin{proof}
\proofstep{Local lines in the integral basis.}

Use the integral frame of Proposition~\ref{prop:count-two-integral-induced-seed}.
Write \(\gamma=b/a\) in its basis \(P\). On each
marked decomposition group the split quadratic extension
gives the two matrices \(A=r_2(h)\) and
\(B'=r_2(\tau^{-1}h\tau)\). Both reduce to upper triangular
matrices with diagonal \((1,\bar\eta(h))\). There is a
decomposition element \(s\) at which \(\bar\eta(s)\ne1\):
the marked nontrivial power of the \(\alpha\)-root in the
quadratic tower supplies it. This remains necessary after
restricting the two ordinary lines from \(k_0\).

For an ordinary line reducing to character \(\bar\eta\),
write a nonzero eigenvector as \((x,y)\). At \(s\) its
eigenvalue minus the upper left entry is a unit. For a matrix \(A=(A_{ij})\) at \(s\), the first eigenvector
equation is
\[
 (\lambda-A_{11})x=A_{12}y.
\]
A nonzero eigenvector therefore has \(y\ne0\) and
\[
 u=x/y=A_{12}/(\lambda-A_{11})\in O .
\]
Rescaling gives the integral generator \((u,1)\).
For a line reducing to character \(1\), use
\[
        (\lambda-A_{22})y=A_{21}x .
\]
Here the left coefficient is a unit and \(A_{21}\) is in the
maximal ideal, giving \((1,w)\) with \(|w|_\ell<1\). Each generator works
on its entire decomposition group.

The following two frames record the integral pivots. Put
\[
 N(p,q,r,s,z)=
 \begin{pmatrix}
 1&0&0&0\\ p&1&0&0\\ q&s&1&0\\ r&z&-p&1
 \end{pmatrix}.
\]
It is symplectic if \(z=q-ps\). If the first block has
line \((u,1)\) and the second has line \((1,w)\), use
\[
 (p,q,r,s)=\left(\gamma u,-1,\gamma u,
                \frac{-2w}{1+\gamma uw}\right).
\]
If the first has line \((1,w)\) and the second \((u,1)\), use
\[
 (p,q,r,s)=\left(-u,1,u,\frac{-2w}{\gamma+uw}\right).
\]
Both denominators are units. To check the flag, denote the
two lines by \(U_\eta,U_1\) according to their reductions.
In the original two-block coordinates the first two columns of
\(PN\) are
\[
 \begin{pmatrix}a(p+r)\\ b(1-q)\\ -p+r\\ 1+q\end{pmatrix},
 \qquad
 \begin{pmatrix}a(1+z)\\ -bs\\ z-1\\ s\end{pmatrix}.
\]
In the first case the first column is \(2b(u,1,0,0)^t\).
In the second it is \(2(0,0,u,1)^t\).
The two-block pairs of the second column are
\[
 \begin{array}{c|cc}
 \text{case}&\text{first block}&\text{second block}\\ \hline
 1&-bs\,(u,1)&(z-1)(1,w)\\
 2&a(1+z)(1,w)&s\,(u,1).
 \end{array}
\]
These formulas follow by substituting \(z=q-ps\) in the two
choices of \(N\).

The second column has a unit
coefficient on \(U_1\) in each case, since \(w\) belongs to
the maximal ideal and \(2,a\) are units. Thus the first two
columns span \(U_\eta\oplus U_1\).
The third space is the perpendicular of the first for the
block alternating pairing, so it adds the other graded
line of the block containing \(U_1\). Consequently
\(N^{-1}P^{-1}\operatorname{diag}(A,B')PN\) is upper
triangular. The ordered exponents are \(k+t,k,0,-t\) in
both cases. Multiplication by \(\delta\) multiplies all four
characters by the same unit reducing to one. In particular
the integral flag has \(N_{20}=\pm1\) and
\(N_{21}\equiv0\), and every nonzero residual off-diagonal
entry of \(N\) has degree in \(Q\).

\end{proof}

The integral ordinary flag identifies the residual labels locally. The residual path also needs the global extension pattern and its permitted ramification.

\begin{proposition}[Residual support and inertia of the induced seed]\label{prop:count-two-residual-seed-support}
In the integral induced frame, let \(R_{ij}\) denote the residual off-diagonal entries. Their support satisfies 
\[
 \operatorname{supp}_{\rm off}(R)\subseteq\{(1,0),(3,2),(3,0),(1,2)\},
\]
 with the paired entries given by the two expressions below, and the prescribed inertia is killed after the scalar
correction. These conditions, together with the local flag
constructed above, put the seed at the specified residual
vertex.
\end{proposition}

\begin{proof}
Keep the same frame
\(P\) before changing the local flags. On
\(G_F\cap G_{k_0}\) let \(x,y\) be the reduced upper
right entries of the two blocks. Direct multiplication by
\(P^{-1}\) and \(P\) gives diagonal
\((\bar\eta,1,\bar\eta,1)\) and the only possible
off-diagonal entries
\[
 R_{10}=-R_{32}=(\bar\gamma x-y)/2,\qquad
 R_{30}=-R_{12}=(\bar\gamma x+y)/2 .
\]
There is an element with \(x\ne0\) by the nonsplit lattice
choice. Since
\[
                   R_{10}+R_{30}=\bar\gamma x\ne0,
\]
the element has either the \(-w_\alpha\) entry or both the
\(w_\beta\) and \(-w_\vartheta\) entries, using \(R_{30}=-R_{12}\). On the other coset multiply by the
diagonal matrix of \(\tau\), which preserves this support
and gives the specified residual labels there.

For inertia away from \(\ell\) required to be killed in
the tower representation, the fitted data make the host
character trivial there and give trivial residual
\(\bar\eta\). The latter gives \(\eta=1\) there since
its order is prime to \(\ell\). The quadratic sign and the
cyclotomic character are also trivial on this inertia.
The two rank-two restrictions have diagonal inertial type
\(1\oplus\eta\), including for the conjugate place. Both
blocks act trivially there. On this inertia \(\delta^2=1\) and \(\bar\delta=1\). For every inertia element \(h\) we have
\[
 (\delta(h)-1)(\delta(h)+1)=0,\qquad \delta(h)+1\in O^\times
\]
because \(\ell\) is odd. Thus \(\delta(h)=1\).

At
\(\ell\) consider inertia fixing the cyclotomic roots.
Here \(\omega=\phi_0=1\) since \(k_0\) splits at \(\ell\).
The fitted tower makes the diagonal residual representation
trivial on this inertia, hence its second character
\(\overline\Psi=\bar\xi\) is one there. The restriction of
the host character to the seed layer has order prime to
\(\ell\), so \(\xi=1\) there as well. The formula for
\(\delta^2\) now gives \(\delta^2=1\), and its residue
again selects \(+1\). Thus the ordered characters are
trivial on this part of inertia.
These give the seed flag, residual vertex and ramification
conditions in the same normalization.
\label{cor:ordinary-promodularity-10}

\end{proof}

\subsubsection{Marked cup cancellation and the deep layer}
\label{par:marked-cup}

The residual labels at the two endpoints now agree with the chart. Connecting the endpoints still requires removing the obstruction given by a cup product at the marked places.

\begin{proposition}[Three links from the marked cup correction]\label{prop:count-two-residual-links}
For the normalized even-order target and seed, the marked
cup term can be canceled after the indicated finite residue
extension and cyclic layer. On a common sufficiently deep layer, write \(R_{\rm aux}\) for the auxiliary datum and \(R_{\rm root}\) for its one-root projection. The three residual lines have endpoints 
\[
 R_{\rm target}\longrightarrow R_{\rm aux}\longrightarrow R_{\rm root}\longrightarrow R_{\rm seed}.
\]
 Each vertex satisfies one of the component tests of
Proposition~\ref{prop:tower-component-propagation}.
\end{proposition}

\begin{proof}
\proofstep{Normalize the residual cocycles.}

On the starting layer \(F\), use the normalized residual
datum with its prescribed local flags. The \((2,0)\)
global and frame entries give a cocycle \(y\) and local
primitives \(n_v\) of degree \(-w_\Psi\).
Thus \(y(g_0)=0\) and
\[
 y(s)=(1-\beta_-(s))n_v
\]
on the local group at \(v\mid\ell\),
where \(\beta_-\) is the residual character of degree
\(-w_\Psi\). Denote the character of degree \(-w_\alpha\)
by \(\alpha_-\), and put \(\gamma_-=\alpha_-\beta_-\).

More generally let \(y_j\), \(j\) in a finite set, be
such cocycles for which we need one common \(x\).
The residual character values belong to a finite field
\(k_{\mathrm{res}}\). In the ambient residue field,
let \(V_y\) be the \(k_{\mathrm{res}}\)-span of the images of the
\(y_j\). Each image is finite, as a continuous image of a compact group in a discrete field. In particular
\[
 \dim_{k_{\rm res}}V_y\le\sum_j|\operatorname{im}y_j|<\infty .
\]
\proofstep{Kill the marked cup term by the auxiliary extension.}

The normalized root cocycles determine the marked obstruction. Passing to the chosen auxiliary extension removes this cup term while allowing the residual links to be compared on a common layer.


Consequently \(V_y\) is a finite additive group, including
when the \(y_j\) take transcendental values over \(k_{\mathrm{res}}\).

On the finite \(\mathbf F_\ell\)-module
\(M=k_{\mathrm{res}}\oplus V_y^{\{j\}}\) use the action
\[
 g\cdot(u,(b_j)_j)
   =\bigl(\alpha_-(g)u,\,
      (\gamma_-(g)b_j+y_j(g)\alpha_-(g)u)_j\bigr).
\]
The cocycle identity for \(y_j\) is exactly the
multiplication rule for this action. The module is
unramified outside the finite set \(S_0\) of allowed
places, which includes the places above \(\ell\).
Chebotarev gives infinitely many rational primes
splitting in the normal closure of the finite field
extension cut out by this action, the residual data,
the mod-\(\ell\) cyclotomic character and \(F\). Use
places of \(F\) over these primes as possible new
ramification places. Both \(M\) and \(\mu_\ell\) are
trivial on their decomposition groups.

Here is the cohomology argument for a cocycle with
nonzero first coordinate. Fix
\[
 e=(1,(0)_j)
\]
and add
\(t\) of the split places to \(S_0\). At the original
finite places impose zero in local \(H^1\). At each
new place allow all unramified classes and require
the restriction to inertia to take values in
\(\mathbf F_\ell e\). Local class field theory at a new place gives
\[
 H^1(G_v,M)\simeq M_{\rm unr}\oplus M_{\rm tame},\qquad M_{\rm unr}\simeq M_{\rm tame}\simeq M.
\]
The inertia condition cuts the second summand down to \(\mathbf F_\ell e\). Thus at each added place \(v\) the local factor in the cardinality formula is
\[
 \frac{|M\oplus\mathbf F_\ell e|}{|H^0(G_v,M)|}
   =\frac{|M|\ell}{|M|}=\ell .
\]
The product over the \(t\) new places is \(\ell^t\). The oddness of \(\ell\)
kills \(H^1\) at the infinite places.

The Greenberg--Wiles
cardinality formula from Poitou--Tate bounds the order
of this Selmer group below by \(2^t/C_0\), with
\[
 C_0=|M^\vee(1)|\,|M|^{[F:\Q]+\#S_0}.
\]
Indeed the dual Selmer order and the global invariants
of \(M\) are at least one, the global dual invariants
have cardinality at most \(|M^\vee(1)|\), and the
remaining local invariant factors cost at most one
\(|M|\) per original or infinite place. The number of
cocycles unramified outside \(S_0\) is finite independently
of \(t\). Each such cocycle and the action define a
homomorphism into the finite affine group
\(M\rtimes\operatorname{Aut}(M)\), so this follows by
Hermite--Minkowski and finiteness of the possible maps
for each field. Taking \(t\) large gives a cocycle
ramified at a new place. On that inertia group the
nonzero value is in \(\mathbf F_\ell e\), where the
action on \(M\) is trivial.

Write the resulting cocycle as \(c=(x,(z_j)_j)\).
The action formula gives
\[
\begin{split}
 x(gh)&=x(g)+\alpha_-(g)x(h),\\
 z_j(gh)&=z_j(g)+\gamma_-(g)z_j(h)
                 +y_j(g)\alpha_-(g)x(h).
\end{split}
\]
At an original place over \(\ell\), write \(c(s)=
(1-s)(m_v,(b_{j,v})_j)\). The resulting local equation is
\[
 z_j(s)=(1-\gamma_-(s))b_{j,v}
                  -y_j(s)\alpha_-(s)m_v .
\]
Subtract a global coboundary to make
\(x(g_0)=z_j(g_0)=0\). Its coordinates can be chosen in
this order because \(\alpha_-(g_0)\) and \(\gamma_-(g_0)\)
are both different from one, and \(y_j(g_0)=0\).
This changes the local primitives by the corresponding
constants and leaves the value on the chosen
new inertia group fixed.

The symplectic realization is explicit. Put
\[
 U(x,y,z)=
 \begin{pmatrix}
 1&0&0&0\\ x&1&0&0\\ y&0&1&0\\ 2z-xy&y&-x&1
 \end{pmatrix}.
\]
Then \(U\) preserves the alternating form and
\[
 U(x,y,z)U(x',y',z')=
 U(x+x',y+y',z+z'+yx').
\]
The diagonal matrix of the four residual characters
acts on these three coordinates by
\(\alpha_-,\beta_-,\gamma_-\), respectively. Thus the
global representation is
\[
 g\longmapsto
 U(x(g),y_j(g),z_j(g))\operatorname{diag}(\bar\chi_i(g)).
\]
The local frame
\(U(m_v,n_{j,v},b_{j,v})\) conjugates it to the marked
diagonal on the local group. Both marked \(-w_\Psi\)
entries are retained. The common \(-w_\alpha\) global
and local entries come from \(x\) and \(m_v\).

At each added rational prime \(p\), the base field
is split and \(p\equiv1\bmod\ell\). The degree-\(\ell\)
subfield of \(\Q(\zeta_p)\) cuts inertia by the unique
index-\(\ell\) subgroup of its tame pro-\(\ell\)
quotient. On that inertia the cocycle is a homomorphism to
an elementary abelian \(\ell\)-group. Its tame pro-\(\ell\)
restriction factors through
\[
                        \Z_\ell/\ell\Z_\ell .
\]
Thus the new extension kills it.
The compositum of these auxiliary extensions is
abelian of odd degree and unramified at \(\ell\).

In particular, this step preserves tameness at the
starting layer. The nonzero global entry survives
restriction. If \(x(g)\ne0\), its normalization gives
\[
 x([g_0,g])=(\alpha_-(g_0)-1)x(g)\ne0
\]
for the commutator \(g_0gg_0^{-1}g^{-1}\). For an abelian extension \(L'/L\) of the starting layer, projection to \(\operatorname{Gal}(L'/L)\) kills commutators. Thus
\[
 g_0gg_0^{-1}g^{-1}\in G_{L'}.
\]
The new distinguished
element is an allowed power of \(g_0\), so it still
gives diagonal matrices. This yields the residual
datum with its original allowed ramification in the
subsequent deep layers.
\label{cor:ordinary-promodularity-11}

The three lines use the following operations on these
residual data. Write \(R_0,R_1\) for the target and auxiliary matrices on the first line. For its retained entry one has 
\[
 (1-s)(R_0)_{20}+s(R_1)_{20}=(R_0)_{20}\quad\text{in }k[s],
\]
 because the common \(-w_\Psi\) coordinate is fixed, and the same identity holds in the frames. Then retain the \(-w_\alpha\)
extension and discard the other entries. The local
flags of this second endpoint can be chosen to share
the \((2,0)\) entry of the seed, using the equality of
the two residual characters on the relevant local
subgroup. Finally interpolate between this one-root
point and the seed, with the common local \(-w_\Psi\)
flag entry and with the global \(-w_\Psi\) entry zero
on both sides. The block orders and the entries fixed at both endpoints are
\[
 \begin{array}{c|c|c|c}
 \text{line}&\text{ordered blocks}&\text{common entry}&\text{paired entry}\\\hline
 1,3&\{0,2\},\{1,3\}&(2,0)&(3,1)\\
 2&\{0,1\},\{2,3\}&(1,0)&(3,2).
 \end{array}
\]
The matrices are block triangular.

In each order the symplectic equation determines the paired entry from the common one, so both diagonal blocks are fixed, also in the frames. In either block order write \(R_1(g)=R_0(g)+D(g)\).
In these block coordinates the differences all occupy one strict off-diagonal block. For instance, in lower block order,
\[
 \begin{pmatrix}0&0\\ X&0\end{pmatrix}
 \begin{pmatrix}0&0\\ Y&0\end{pmatrix}=0 .
\]
In particular \(D(g)D(h)=0\) for all \(g,h\). Subtracting the two endpoint multiplication laws gives
\[
 D(gh)=R_0(g)D(h)+D(g)R_0(h).
\]
Set \(R_s=R_0+sD\). Square-zero multiplication eliminates
\(s^2 D(g)D(h)\), so
\[
                R_s(g)R_s(h)=R_0(gh)+sD(gh)=R_s(gh).
\]
In the symplectic and frame products the two off-block differences again have zero product (the adjoint preserves that block ideal).

For any of the resulting equations \(H=0\), substitution of the interpolated matrices gives
\[
 H(s)=(1-s)H(0)+sH(1)=0 .
\]
Thus the polynomial line respects these equations identically. The target passes the first test with
\(E(b_0)\) or the third test, by its lattice normalization.
If it has a nonzero global entry of degree \(-p_0\), the
first intermediate vertex retains it and adds \(E(-a_0)\),
so that the first test applies there. If its entire
\(-p_0\) global entry vanishes, so does that of the
first intermediate vertex. That vertex then has support
in \(V\), inherits the nonzero local \((2,0)\) entry and
passes the third test. The second intermediate vertex has
only the degree \(-a_0\) globally and takes its nonzero
\((2,0)\) frame entry from the seed.

It passes the third
test. The seed has global support in \(V\) and frame support
in \(Q\) with nonzero \((2,0)\) entry. Its ordinary local image supplies one of the following inclusions.
\[
 \{-a_0\}\subseteq\operatorname{supp}_{\rm glob}(\text{seed})
 \quad\text{or}\quad
 \{b_0,-t_0\}\subseteq\operatorname{supp}_{\rm glob}(\text{seed}).
\] In the second case the tower retains an element
at each \(\ell\)-adic place on which the original
\(\alpha\)-character remains nontrivial after the prescribed
power change. Its residual value on that element equals the
value of \(\bar\chi_0/\bar\chi_1\), giving the fourth test.
 The cocycle
bound \(C\) is taken after that extension. The large
degree and local degree are then obtained in the same
cyclotomic layer, and the seed is constructed there.
The order of choices allows the preceding component
proof to apply to that chain.
\label{cor:ordinary-promodularity-12}

\end{proof}

\subsubsection{Characteristic-\(\ell\) specialization}

Characteristic-zero propagation uses ordinary support at both levels. In residue characteristic the same ideals instead constrain the characteristic polynomials, a conclusion valid even when the point has coefficients of positive transcendence degree.

\begin{proposition}[Characteristic polynomials on the split residual layer]\label{prop:tower-characteristic-ell}
In characteristic \(\ell\), a normalized target killed by \(J\)
and with the prescribed split residual inertia has, for its matrices \(r(g)\) and the prescribed diagonal characters \(d_i\), the identities 
\[
 \det(X-r(g))=\prod_i(X-d_i(g))\qquad(g\in G_F).
\]

The assertion holds over the original valuation ring of the point.
\end{proposition}

\begin{proof}
Suppose the target valuation ring has characteristic \(\ell\).
Let \(a:O\to V\) be its coefficient map and
\(y:A\to V\) the normalized point. The kernel of \(a\)
contains \(\ell\), so it contains the uniformizer of
the coefficient DVR. The same holds after the finite
coefficient extension used to define \(B\).
For each \(s\) in the full inertia of a marked place,
normalization gives the equality in \(V\)
\[
    y(\psi_{v,i}(s))=a\bigl((r_0)_{ii}(s)\bigr).
\]
The split residual system of the fitted tower is trivial on
that inertia. Thus \((r_0)_{ii}(s)-1\) lies in the
coefficient maximal ideal, and the displayed equality makes
\(y(\psi_{v,i}(s)-1)=0\). The restriction of \(y\) to \(B\)
kills \(J\) by propagation. The maximal-ideal power in Proposition~\ref{prop:tower-closed-specialization}
therefore maps to zero.

Since \(V\) is a domain,
\(B\to V\) factors through its finite residue field. All
characteristic polynomials on \(G_F\) specialize to those
of the split residual system, including for the original
point since conjugation preserves characteristic polynomials.
This gives
the finite-index characteristic-polynomial comparison
used below, even when the original characteristic-\(\ell\)
point has a transcendental coordinate.
\label{cor:ordinary-promodularity-13}

\end{proof}

\subsubsection{Classical specializations of both Hecke quotients}
\label{par:count-two-classical}
\hypertarget{source-ordinary-count-ii-point-output}{}

For a characteristic-zero target the two ordinary evaluations yield representations. To compare their Frobenius polynomials with those of the original point, the evaluations must first be realized at classical weight.

\begin{proposition}[Classical control at the propagated target]\label{prop:two-link-classical-control}
At a characteristic-zero even-order target reached by the
two-link propagation, each evaluation factors through a classical
ordinary eigensystem in its link. Let \(r_i\) be the Galois representation of the eigensystem at level \(i\), \(r\) the representation at the chosen determinant point and \(\kappa\) the common twisting character. Then 
\[
 \det(X-r_i(g))=\det(X-\kappa(g)r(g))\qquad(g\in G_L).
\]

\end{proposition}

\begin{proof}
Use the same Hecke quotients \(R_i/J_i\) as in
Paragraph~\ref{par:count-two-ideals}. At a characteristic-zero
target \(y\) reached by Paragraph~\ref{par:count-two-propagation},
the evaluations \(x_i=y|_B f_i\) factor through these quotients.
Write \(m=(j+1,j,1,0)\) for the rational inertial exponents.
On the selected inertia groups above \(\ell\), \(\kappa\) is
an integral cyclotomic power up to a finite character. It
adds a common integer to the four exponents at each place.
Thus the coefficients have the algebraic parts given by
\(\lambda_{\tau,3-i}=-m_{w,i}-i\), with \(m_{w,i}\) the exponents
after twisting. The difference in the highest weight order is
\[
       \lambda_{\tau,2-i}-\lambda_{\tau,3-i}
                  =m_{w,i}-m_{w,i+1}-1\ge0 .
\]
The difference inequality gives dominance. After the algebraic
part is removed the remaining compact torus character is smooth. These conditions put the point in the arithmetic ordinary locus used in the cohomology calculation.

Here is the control argument at this weight. Write \(O_x\)
for a coefficient DVR containing the weight values, \(\pi_x\)
for its uniformizer, \(W_\lambda\) for the algebraic coefficient
representation, and \(p_\lambda\) for its lowest weight quotient
for the upper unipotent group. Use a lattice in \(W_\lambda\)
stable under the integral local groups and decomposed into
algebraic torus weights. Such a lattice is given by the integral
highest weight construction at the embeddings, tensored with
the required integral determinant powers. Write \(L\subset W_\lambda\) for this integral lattice and \(L_\mu\) for its weight submodules. Then
\[
 L=L_{\rm low}\oplus\ker(p_\lambda|_L),\qquad L_{\rm low}\simeq O_x,
 \qquad \ker(p_\lambda|_L)=\bigoplus_{\mu\ne\mu_{\rm low}}L_\mu .
\]
This gives the rank-one projection integrally.
Choose \(c\) so that the smooth character kills \(T_c\) and
\(\bar N_c\) acts trivially on the coefficient lattice modulo \(\pi_x\).
Use \(J_c=NT\bar N_c\) for the following comparison.

Trivial stabilizers identify the classical module with the
finite module of values on double coset representatives.
In the left-invariant realization \(F_\ell\) of
Section~\ref{sec:count-two-seed-specialization}, the dominant
translation \(z\) acts by the sum with coefficient matrices
\(W_\lambda(b)W_\lambda(z)\), where \(b\) ranges over
\(N/zNz^{-1}\). Normalize each matrix by the lowest-weight
scalar of \(z\). The normalized \(W_\lambda(z)\) is integral
and kills the kernel of \(p_\lambda\) modulo \(\pi_x\):
For any other weight \(\mu\), write
\(\mu-\mu_{\rm low}=\sum_\alpha n_\alpha\alpha\) with
\(n_\alpha\ge0\), some \(n_\alpha>0\).

Then
\[
 v_{\pi_x}\!\left(\frac{\mu(z)}{\mu_{\rm low}(z)}\right)
      =\sum_\alpha n_\alpha v_{\pi_x}(\alpha(z))>0.
\]
Modulo \(\pi_x\), the normalized diagonal operator kills all
these higher weight spaces. The integral unipotent matrices
\(W_\lambda(b)\) preserve the quotient projection, so the same
holds for the normalized translation sum. The resulting operator intertwines the
ordinary operator on the scalar functions under \(p_\lambda\).
On characteristic-zero functions the projection is into the
scalar module with lower unipotent depth allowed to increase
with precision, by continuity of the coefficient action.

Modulo \(\pi_x\), the scalar module at depth \(c\) is obtained
by projecting the algebraic values on the same cosets. Let \(P:C\to S\) be this projection on the
depth-\(c\) residue spaces. It is onto, commutes with the
normalized translation \(U\), and \(U\ker P=0\).
Since \(P\) is equivariant and onto, \(P(U^a C)=U^a S\)
for every \(a\).

At the stable images its kernel is
\[
    C^{\rm ord}\cap\ker P \ \subseteq\ C^{\rm ord}\cap\ker U=0 .
\]
The last equality uses invertibility on \(C^{\rm ord}\).
This proves the comparison on the stable images. The integral ordinary projectors reduce
to these stable-image projectors. The scalar ordinary module
is finite free with the value filtration of
Paragraph~\ref{par:ordinary-functions}; its ordinary reduction
is exactly the stable image at depth \(c\). Thus the projected
classical ordinary module maps isomorphically onto the scalar
ordinary module, by freeness and reduction modulo \(\pi_x\).
The spherical Hecke actions and their residual localizations
commute with the projection. Together with the scalar weight
base change of Section~\ref{sec:count-two-seed-specialization},
this gives the classical realization on each \(M_i\).

Any field evaluation of \(R_i/J_i\) at this weight occurs in
that specialization. Indeed, putting \(A_i=R_i/J_i\), finiteness and faithfulness imply
\[
 \operatorname{Supp}_{A_i}M_i=V(\operatorname{Ann}_{A_i}M_i)
 =\operatorname{Spec}A_i,\qquad M_i\otimes_{A_i}\kappa(\mathfrak p)\ne0 .
\]
The last assertion is Nakayama's lemma at the evaluation prime \(\mathfrak p\).
After extending coefficients, that fiber has the evaluation
as an eigensystem. In the automorphic decomposition choose a
contributing constituent with these spherical eigenvalues.
Its semisimple Galois attachment has the characteristic
polynomials of \(\kappa\,r|_{G_L}\), by the spherical
normalization on the dense set of split good Frobenius classes.
This equality holds on all of \(G_L\), including at reducible
targets. At an irreducible target it identifies the actual
Galois representations.
\hfill\mbox{\cite{newton-thorne-unitary}}\label{cor:ordinary-count-ii-point-output-1}

\end{proof}

\subsubsection{The Frobenius ratio test at a full block}
\label{par:count-two-purity}

A hypothetical full two-dimensional block imposes a determinant ratio at good primes. The polynomials just obtained from both ordinary levels give an incompatible restriction on their roots.

\begin{proposition}[Exclusion of a full block]\label{prop:count-two-full-block-excluded}
At a characteristic-zero full-block target, write \(\alpha_i(w)\) for the Frobenius roots at a good prime in the block test. The two-link specialization gives 
\[
 |\alpha_i(w)/\alpha_j(w)|<Nw\qquad\text{for every }i,j.
\]

This contradicts the determinant ratio of the two blocks.
\end{proposition}

\begin{proof}
In the transfer of a contributing constituent from link \(i\),
a discrete Speh summand of length at least two would give two
summands of its Galois representation differing by
\(\epsilon_\ell\) or its inverse. Take the chosen Frobenius
element \(g_i\) at the exceptional place of this link.
All four roots of the characteristic polynomial at \(g_i\)
reduce to the same nonzero residual eigenvalue, by the
construction of \(J_i\) and \(f_i\). The eigenvalues of both
summands are units, since they are roots of that integral
polynomial with unit constant coefficient. If \(z\) and \(z\epsilon_\ell(g_i)^{\pm1}\) are such roots,
then \(z\) is a unit. Their reductions cannot agree:
\[
              \bar z\,\bar\epsilon_\ell(g_i)^{\pm1}\ne\bar z .
\]
This contradicts the common nonzero residual eigenvalue.

Consequently every discrete summand is cuspidal. The
cuspidal summands in the unitary base change are tempered
at the good places, by the results cited after
Corollary~1.3 of Newton and Thorne. The common algebraic
normalization multiplies every root at a place by the same
power of its residue norm. Thus all roots have the same
absolute value there, and in particular each ordered ratio
of absolute values is strictly less than the residue norm.

The purity condition for that calculation is the ratio condition just proved. It passes to the rational representation:
If \(\alpha_i\) are the rational roots, restriction to a
place \(w\mid p\) of residue degree \(f\) gives \(\alpha_i^f\).
The common scalar twist cancels. The inequalities over \(L\) imply
\[
              |\alpha_i/\alpha_j|^f<Nw=p^f
            \quad\Longrightarrow\quad|\alpha_i/\alpha_j|<p.
\] The property
also holds at primes splitting completely in \(L\), which
suffices for the block-determinant test. Apply the link to
the nonsplit representative constructed in
Paragraph~\ref{par:full-block-extension}. Equality of polynomials
gives the same ratio bound at the split block point.
\hfill\mbox{\cite{newton-thorne-unitary}}\label{cor:ordinary-count-ii-point-output-2}

\end{proof}

\subsubsection{Attachment at an irreducible target}
\label{par:count-two-glue}

At an irreducible target, equality of the characteristic polynomials identifies the restricted representations themselves. Using both levels then compares their local parameters outside their respective exceptional sets.

\begin{proposition}[Cuspidal attachment from the two links]\label{prop:count-two-cuspidal-attachment}
At a rational irreducible target of the even-order construction,
the restrictions to \(G_L\) are absolutely irreducible. The
cuspidal transfers from both ordinary links give the same
representation there; if \(S_1,S_2\) are the respective sets of exceptional places, the compatible local parameters cover both choices because 
\[
 S_1\cap S_2=\varnothing.
\]

\end{proposition}

\begin{proof}
For a rational irreducible target, restriction to \(G_L\)
is absolutely irreducible. The restriction to a normal
subgroup is semisimple, by taking translates of a simple
subspace. If the restriction were reducible, the action of
the abelian quotient \(\operatorname{Gal}(L/\Q)\) on its
commutant would have a nontrivial character. An eigenvector in the commutant is a nonzero intertwiner
for a self-twist \(\chi\), hence invertible by irreducibility on
the rational group. It yields
\[
              (1-\chi(g))\operatorname{tr}r(g)=0 .
\]
On any coset where \(\chi\ne1\) the trace is zero. The reduction of the trace on such a coset
is a linear combination of the four distinct restricted
residual labels with nonzero coefficients. Linear independence
of characters rules out its identical vanishing. This proof
also applies to any target obtained by the change of frame
or scalar extension in the modular-locus argument.

Each link now gives a cuspidal transfer over \(L\): an isobaric
sum with several terms would make the attachment reducible.
At a common good prime the same twist and normalization give
\[
     P_{\Pi_1,w}(T)=\det(T-\kappa r(\mathrm{Frob}_w))
                         =P_{\Pi_2,w}(T),
\]
where \(P_{\Pi_i,w}\) denotes the algebraically normalized local
polynomial. Both links use the same specialization \(y\). Strong multiplicity
one identifies the transfers. For each selected place, use
the local comparison at a level where it is outside the
exceptional pair of places over \(L^+\). Indeed if \(S_i\) is the exceptional pair over the \(i\)-th choice of \(v_0\), their rational primes are distinct. Consequently
\[
 S_1\cap S_2=\varnothing,\qquad S_1^c\cup S_2^c
 =\{\text{all finite places}\}.
\]

At each of the exceptional places the other link therefore supplies the local comparison. At the other places outside the level the
transfer is unramified, and at the marked places above
\(\ell\) the split local comparison gives the Weil--Deligne
parameter of the same target.
The common scalar \(\kappa\) is locally algebraic at \(\ell\)
with integral labeled exponents and finite inertia away from
\(\ell\). By reciprocity it is the \(\ell\)-adic avatar of an
algebraic Hecke character. Twisting the transfer by its inverse
gives the cuspidal attachment of \(r|_{G_L}\), with the rational
similitude as its multiplier. In particular the target weights
are recovered after removing the twist.
\hfill\mbox{\cite{newton-thorne-unitary}}\label{cor:ordinary-count-ii-point-output-3}

\end{proof}

\subsubsection{Ordinary flag and rational level}
\label{par:count-two-host-level}

The attachment is obtained after restriction to the fitted field. To count at a prescribed rational level, the original representation must have the required ordered inertia exponents and conductor bounds.

\begin{proposition}[The rational ordinary flag and its level]\label{prop:count-two-rational-level}
For an irreducible even-order target with \(j\ge2\) in
Section~\ref{sec:odd-prime-cases} and the attachment of
Proposition~\ref{prop:count-two-cuspidal-attachment}, the rational
ordinary flag has the prescribed exponents and zero monodromy.
The rational inertia conditions give the level bound for the ensuing quadratic restriction, with 
\[
 M=\ell\,2^{n_{\rm II}}\mathcal N.
\]

\end{proposition}

\begin{proof}
\proofstep{Recover the rational ordinary characters.}

At \(\ell\) write the ordered characters of \(r\) as
\(\kappa_i=\nu_i\epsilon_\ell^{m_i}\) with
\(m=(j+1,j,1,0)\) and \(\nu_i\) finite on inertia. The
character \(\kappa_i/\kappa_k\) can equal the cyclotomic
character on a full decomposition group only if
\(m_i-m_k=1\). Restrict to an open inertia subgroup
killing \(\nu_i/\nu_k\) and use the infinite cyclotomic image
there to see this. As \(j\ge2\), the consecutive gaps are
\((1,j-1,1)\), each at least one.

If \(m_i-m_k=1\) with
\(i<k\), then \(k=i+1\). At the retained Frobenius lift of
Paragraph~\ref{par:count-ii-tower}, the cyclotomic
character and the residual \(\bar\Psi\) are one, and
\(\bar\eta\) has order greater than four. In both permitted frames the labels at this lift are
\[
         (\bar\eta,1,\bar\eta,1)\qquad(\bar\eta\ne1).
\]
Thus all adjacent labels differ at a point where
\(\bar\epsilon_\ell=1\). This rules out the cyclotomic equality at the
marked place of \(L\) as well as over \(\Q_\ell\).

The target is de Rham at \(\ell\), since it is so over the
fitted CM field and de Rhamness descends along a finite
extension. Its local Weil--Deligne representation is
therefore defined. By exactness on the ordinary flag its
graded Weil lines, on the rational Weil group, are
\[
       \mu_i(s)=\iota(\nu_i(s))\,\ell^{m_i\deg(s)}
\]
in the arithmetic Frobenius convention, where \(\deg(s)\)
is the exponent of arithmetic Frobenius in the unramified
quotient.

At the marked place of \(L\) use the restriction
of these lines. A norm ratio between two lines would give
\(\kappa_i/\kappa_k=\epsilon_\ell\) in one order: taking
\(\ell\)-adic valuations at Frobenius forces the exponent
difference to be one, because the \(\nu_i\) take unit
values. Equality then forces the two \(\nu_i\) to agree
by density of the Weil group. The Weil--Deligne relation on the \((i,k)\) entry of \(N\)
forces a norm ratio \(\mu_i/\mu_k=|\cdot|^{\pm1}\) if
that entry survives. Indeed the ratio constraint follows entrywise from the Weil--Deligne identity
\[
 (\mu_i(s)/\mu_k(s)-|s|)N_{ik}=0
\]
in the convention \(r(s)Nr(s)^{-1}=|s|N\). Switching conventions inverts the norm. Since both orders have been excluded, \(N=0\) over \(\Q_\ell\) and after restriction.

At the marked place of \(L\) all \(|\mu_i/\mu_k|\) on
Weil elements are one, by temperedness of the glued
cuspidal transfer and its local compatibility. Removing
the scalar twist preserves this equality. It holds on the
rational Weil group by taking positive powers into this
local subgroup. Choose \(s\) in the rational Weil group
mapping to the uniformizer \(\ell\) under reciprocity. Then
\(\epsilon_\ell(s)=1\) and \(\deg(s)\ne0\). Were \(\kappa_i(s)/\kappa_k(s)\) torsion, its complex modulus
would be one.

Since \(\epsilon_\ell(s)=1\), the displayed
Weil lines would then give
\[
       |\mu_i(s)/\mu_k(s)|=\ell^{(m_i-m_k)\deg(s)}\ne1 .
\]
Temperedness on the restricted Weil group forces absolute value one. If \(s^a\) lies there, then
\[
 |(\mu_i/\mu_k)(s^a)|=|(\mu_i/\mu_k)(s)|^a .
\]
Thus the violation persists on restriction, a contradiction. The character ratios in question have infinite order. A power of \(s\) fixes \(L\), and
further positive powers give the same conclusions over
finite local extensions. This element will be used for
adjoint rigidity.

Away from \(\ell\) the rational chart kills an open subgroup
of inertia. At primes of \(S_b\) it imposes the inertia
quotient from \(\xi\). Elsewhere inertia factors through
the residual character quotient. The given finite characters
have conductors dividing
\[
 M=\ell\,2^{n_{\rm II}}\mathcal N.
\]

The resulting inertia image is finite and abelian, and
Frobenius acts trivially on it because all these characters
come from rational characters. The Frobenius-semisimplified
local parameter is thus a sum of Weil lines with zero
monodromy, and each smooth unit character is trivial on
\(1+M\Z_p\). At \(\ell\) the graded Weil lines just
computed have smooth unit characters from the finite
inertia factors of the rational flag, also of conductor
dividing \(M\). These conductor and local conditions enter the rank-four packet comparison. \hfill\mbox{\cite{newton-thorne-unitary}}\label{cor:ordinary-count-ii-point-output-4}

\end{proof}

\subsubsection{Descent and admissible quadratic restrictions}
\label{par:count-two-descent}

The rational local calculation fixes the level. To use the count of cuspidal representations, one must still verify descent and cuspidality after each imaginary quadratic restriction allowed in that count.

\begin{proposition}[Quadratic packet restrictions of the even-order target]\label{prop:count-two-quadratic-packets}
After the stated descent to \(K\), every admissible imaginary
quadratic restriction on the fixed packet list has the regular
algebraic cuspidal attachment at the level \(M\) from the rational local calculation. Its ordered weights are 
\[
 (j+1,j,1,0).
\]

\end{proposition}

\begin{proof}
\proofstep{Descend to the fitted quadratic field.}

Let \(K\subset L\) be the imaginary quadratic field of the
fitted tower. We descend the untwisted transfer to \(K\)
along a chain of cyclic prime degree extensions inside the
abelian extension \(L/K\). At each step the upper cuspidal
transfer is invariant under the cyclic Galois group:
on the Galois side conjugation carries the restriction to an isomorphic representation, since
\[
 r(ghg^{-1})=r(g)r(h)r(g)^{-1}.
\]
Thus it preserves the good Frobenius parameters. Strong multiplicity one proves invariance.
Cyclic cuspidal descent then gives a cuspidal representation
over the lower field. Its weight is regular algebraic,
as can be checked after base change at infinity. Apply
unramified Galois attachment for regular algebraic cuspidal
representations over CM fields to this descent. Its
semisimple representation has the same restriction to
the upper field as \(r\), by unramified compatibility and
Chebotarev. This restriction is absolutely irreducible
since the restriction to \(G_L\) is.

After identifying the restrictions, let \(r_1,r_2\)
be the two extensions to the lower group \(G\), and let
\(H\triangleleft G\) be the upper group. For \(g\in G\) put
\[
 C_g=r_1(g)r_2(g)^{-1}.
\]
Write \(r_H\) for the common restriction. Normality of \(H\) gives, for \(h\in H\),
\[
 C_g r_H(h)C_g^{-1}=r_1(g)r_H(g^{-1}hg)r_1(g)^{-1}=r_H(h).
\]
By irreducibility and Schur's lemma this ratio is a
scalar \(\zeta(g)\). Writing \(r_1(g)=\zeta(g)r_2(g)\) gives
\[
 \zeta(gh)r_2(gh)=r_1(g)r_1(h)=\zeta(g)\zeta(h)r_2(gh).
\]
Hence \(\zeta\) is a character, and \(\zeta|_H=1\)
makes it a finite character of the cyclic quotient.
Twist the descended cuspidal representation by the inverse
finite character.
Its attachment is now \(r\) restricted to the lower field.
The polarization of \(r\) and strong multiplicity one give
the corresponding essential conjugate duality on the
descended cuspidal representation. This proves the assertion
at every stage with the same normalization.
\hfill\mbox{\cite{clozel-rajan,harris-lan-taylor-thorne}}\label{cor:ordinary-count-ii-point-output-5}

\proofstep{Pass to each admissible quadratic field.}

The count requires cuspidality and the level bound over each admissible field.


Fix its eighteen imaginary quadratic fields of pairwise coprime discriminants.
In choosing the ramified prime of \(K\) in
Paragraph~\ref{par:count-ii-tower}, avoid these fixed
discriminants in addition to the rational ramification set.
Dirichlet's theorem in that paragraph allows any such finite
exclusion. For \(K'\) on the list with \(r|_{G_{K'}}\) irreducible,
write \(\eta_{KK'/K'}\) for the quadratic character of
\(KK'/K'\). We use
\[
        r|_{G_{K'}}\not\simeq
               r|_{G_{K'}}\otimes\eta_{KK'/K'} .
\]
This suffices for irreducibility on the index-two subgroup. Indeed, reducibility on that index-two group
would give the quadratic self-twist of \(G_{K'}\)
corresponding to \(KK'/K'\).

At the chosen ramified prime
of \(K\), inertia lies in \(G_{K'}\), acts trivially on \(r\)
and takes the value \(-1\) on the quadratic character.
For such an inertia element \(s\), if \(\nu\) is the quadratic character, then
\[
 \operatorname{tr}r(s)=\dim r,\qquad
 \operatorname{tr}((r\otimes\nu)(s))=-\dim r .
\]
These characteristic-zero traces distinguish the two representations. Cyclic base change
from \(K\) to \(KK'\) is cuspidal by irreducibility there,
and the same descent to \(K'\) now gives the required
attachment. The argument applies to any specified finite
list by the same finite exclusion. No degree or discriminant
of the intermediate tower enters the packet bound.

At a place \(w\mid p\) of \(K'\), restricting the rational
Weil lines gives the inertia characters on the descended
side. Their smooth unit characters factor through the
local norm on units. If \(x\in1+M\mathcal O_{K'_w}\), multiplication by \(x\)
on \(\mathcal O_{K'_w}\) has matrix \(I+MA\) over \(\Z_p\).
Thus
\[
           N_{K'_w/\Q_p}(x)=\det(I+MA)\in1+M\Z_p.
\] Away from \(\ell\)
local--global compatibility for the polarized cuspidal
representation identifies the full Frobenius-semisimplified
Weil--Deligne parameters. Above \(\ell\), compatibility up
to semisimplification suffices. All ratios of the Weil
lines have absolute value one on Frobenius, so the
semisimplification admits only zero monodromy. Hence the
local parameters on the descended side are sums of Weil
lines and their inertia characters have the asserted bounds.
Under local Langlands the local representation is a
principal series subquotient.

It is generic by cuspidality,
and a generic subquotient of a reducible induction of
linked rank-one characters has a segment of length greater
than one and nonzero monodromy. So it is the irreducible
principal series of those lines. Let \(K(M)\) be the principal congruence subgroup. If \(\chi_B\) denotes the full inducing character of the Borel \(B\), including its modulus factor, then \(\chi_B|_{B\cap K(M)}=1\). Hence
\[
 f(bk)=\chi_B(b)\quad(b\in B,\ k\in K(M)),\qquad
 f|_{G\setminus BK(M)}=0
\]
is well-defined and right \(K(M)\)-invariant, with \(f(1)=1\). Where
\(p\nmid M\) the inducing characters are unramified and
there is a spherical vector. The weights after descent
are \(j+1,j,1,0\), from the local de Rham flag. This gives
the quadratic restriction at level \(M\) used in the
packet count.
\hfill\mbox{\cite{blggt,blggt-lg,zelevinsky}}\label{cor:ordinary-count-ii-point-output-6}
\end{proof}

\subsection{The fitted tower for the rank-four case}
\label{sec:count-three-tower}

The even-order construction used an induced Hilbert system. For the remaining nonprincipal odd-order character in the indicated congruence range, the starting representation will instead have two rank-two tensor factors.

\begin{established}[Inertia in the odd-order case]\label{hyp:count-three-host}
In the odd-order case of Lemma~\ref{lem:odd-prime-positive-ranges},
take \(\xi\ne1\) of odd order \(o_\xi\), with
\(\ell\nmid o_\xi\), \(\ell\ge5\),
\(\ell\nmid\mathcal Np_1\) and \(\ell-1\mid2j\).
Use \(b=p_1\), \(n=n_{\rm III}\) and the character
\(\eta\) of Paragraph~\ref{par:dyadic-residual}.
Write \(I_\ell\) for the inertia group and \(S_b\) for the places
with nontrivial \(\ell\)-primary inertia in \(\xi\). Then
\[
      \xi(-1)=1,\qquad \eta|_{I_\ell}=\xi|_{I_\ell}=1,\qquad S_b=\varnothing .
\]
\end{established}
\begin{proof}
The order of \(\xi(-1)\) divides both two and \(o_\xi\).
The conductor of \(\xi\) divides \(\mathcal N\); the conductor
of \(\eta\) divides \(2^n p_1\) by its construction. Both are
prime to \(\ell\). Finally every inertia image under \(\xi\)
has order dividing \(o_\xi\).
\end{proof}
Write \(\sigma=\pm1\) for the sign of the second coordinate in
the chosen rational ordinary flag.

\subsubsection{Three auxiliary quadratic characters}

The two tensor labels must separate on the starting field, while the character used for descent has a different prescribed restriction. The auxiliary quadratic fields impose these conditions simultaneously.

\begin{proposition}[Fitted tame fields for the odd-order case]\label{prop:count-three-starting-fields}
For the data of Lemma~\ref{hyp:count-three-host}, the fitted rank-four tower can start with the
three quadratic characters and a tame field \(F_0\) constructed
below. Write \(\lambda\) for the descent character and \(\lambda_0\) for the seed character, of prime conductor \(q_0\). Their two quadratic fields 
\[
 G_K=\ker\lambda,\qquad K'=\mathbb Q(\sqrt{-q_0})
\]
 have the indicated splitting at the marked primes.
\end{proposition}

\begin{proof}
Choose an odd quadratic character \(\lambda_0\) with prime
conductor \(q_0\) outside \(S\), split at \(S\), and
write \(K'=\Q(\sqrt{-q_0})\).
Choose two more odd quadratic characters \(\lambda,\lambda'\)
with distinct prime conductors \(q,q'\), splitting all
the primes in \(S\cup\{q_0\}\).
We may avoid any fixed finite set when choosing these primes.
For each choice Dirichlet's theorem supplies a prime
congruent to \(-1\) modulo eight times the product of the
odd primes required to split. Quadratic reciprocity and
the supplementary law at two then give the splitting
conditions. Put
\[
 \begin{split}
 H&=\{g\in G_{\Q}:\eta(g)=\lambda_0(g),\
                        \bar\Psi(g)=\bar\xi(g)\},\\
 J_{\mathrm{quad}}&=\ker\lambda\cap\ker\lambda',\qquad
 G_{F_0}=H\cap\ker(\lambda/\lambda'),\qquad G_K=\ker\lambda .
 \end{split}
\]
These are open normal subgroups defined by finite abelian
characters. In the second condition in \(H\), the quotient
\(\bar\Psi/\bar\xi=\bar\omega^j\) has order dividing two.
Complex conjugation lies in \(H\), since \(\eta\) and
\(\lambda_0\) are odd and \(\Psi\) and \(\xi\) are even.

It kills \(\lambda/\lambda'\), so \(F_0\) is totally real.
The latter quotient has order two, as seen on inertia at
\(q\), so \([F_0:\Q]\) is even.
Intersect \(G_{F_0}\) with \(\ker\lambda\) to get
\[
              G_{F_0K}=H\cap\ker\lambda\cap\ker\lambda'.
\]
At \(q\), the unramified character \(\lambda'\) kills inertia.
The defining relation \(\lambda=\lambda'\) on \(G_{F_0}\)
makes \(\lambda\) trivial there too. Since \(\lambda\) is already
unramified elsewhere, \(F_0K/F_0\) is unramified at all finite places.
The field \(K\) is imaginary and split at every prime
in \(S\). Above each place of \(F_0\) over \(\ell\),
choose a decomposition embedding on the specified side of
\(K\). This is possible because the two places of \(K\)
above \(\ell\) are split and the chosen side specifies one
of the two extensions to \(F_0K\).

The inertia image in \(\operatorname{Gal}(F_0/\Q)\) at
\(\ell\) has exponent dividing two, since here the only
possibly nontrivial defining quotient is \(\bar\omega^j\).
In particular \(F_0\) is tame there.
On \(G_{F_0}\), the reduced characters satisfy
\[
        \bar\alpha=\lambda_0,\qquad\bar\Psi=\bar\xi.
\]
Both are trivial on inertia at \(\ell\), so their products
give the identity full diagonal there.
At any finite prime away from \(\ell\), if the
residual diagonal is the identity on an inertia element
of \(F_0\), then \(\bar\Psi=1\), and hence \(\xi=1\)
there by injectivity of reduction on its roots of unity.
The order condition for the restriction of \(\xi\) already
holds on all of \(G_{F_0}\) and at each subsequent layer.

\end{proof}

\subsubsection{The normalizer and an inertia element at \(\ell\)}

Restricting to the chosen field could identify the residual labels. An explicit normalizing element distinguishes the tensor factors there, and an inertia element tests the marked order at the residual prime.

\begin{proposition}[Marked elements on the starting field]\label{prop:count-three-marked-elements}
The starting field of Proposition~\ref{prop:count-three-starting-fields}
retains the normalizer \(g_0\) distinguishing the two tensor labels
and an inertia element \(u_0\) above \(\ell\) for the marked chart. The finite characters take the values 
\[
 \lambda_0(g_0)=\eta(g_0)=-1,\qquad \omega(g_0)=1.
\]
 Suitable powers of \(g_0\) and \(u_0\) lie in the deeper layer.
\end{proposition}

\begin{proof}
\proofstep{Choose \(g_0\) and \(u_0\) by their residual character values.}

Take \(a\) in inertia at two with cyclotomic value \(-1\)
on the \(2^n\)-th roots of unity, using total ramification.
Its \(\eta\)-value is \(-1\), and its \(\omega\)-value is
one. Raise it to \(o_\xi\) to kill \(\xi\) while preserving both
values. Choose \(b\) in inertia at \(q_0\) with
\(\lambda_0(b)=-1\). The characters \(\xi,\eta,\omega\)
are all unramified there.

Choose \(c\in G_{\Q}\) with
\(\omega(c)=1\), \(\eta(c)=1\) and \(\xi(c)\ne1\)
as follows. The odd order character \(\xi\) has no
dyadic part and its conductor is prime to \(p_1\ell\).
By the cyclotomic Chinese remainder theorem choose an
element with nontrivial \(\xi\)-value which acts
trivially on the \(\ell\)-th and \(p_1\)-th roots of unity.
Raise the chosen element \(c_0\) to \(2^n\). Its values
satisfy
\[
 \eta(c_0^{2^n})=1,\qquad \omega(c_0^{2^n})=1,\qquad
                        \xi(c_0^{2^n})\ne1 .
\]
The last assertion uses the odd order of the nontrivial
\(\xi(c_0)\).

Now \(g_0=abc^2\) lies in \(H\cap J_{\mathrm{quad}}\).
Indeed, the two factors \(a,b\) lie in the decomposition
groups at split primes of \(K\) and the field of
\(\lambda'\), and those two characters kill squares.
The \(\lambda_0\)-value of \(g_0\) and its \(\eta\)-value
are both \(-1\), and \(\omega(g_0)=1\).
Set \(z=\bar\xi(c)^{2\sigma}\), a nontrivial element of
odd order. The residual eigenvalues at \(g_0\), after
division by the first one and inversion, are
\[
                              1,\quad -1,\quad z,\quad -z.
\]
For a root \((u,v)\) the value is \((-1)^u z^v\),
where \((u,v)\ne(0,0)\) and each coordinate is \(0,\pm1\).

If \(v=0\), it is \(-1\). If \(v=\pm1\), it is nontrivial
because \(z^{\pm1}\) has odd order greater than one and so
is neither \(1\) nor \(-1\). In particular
\[
            \operatorname{ord}(-z)
                   =2\operatorname{ord}(z)\ge6
\]
for the root of bidegree \((1,1)\). A residual quadratic
self-twist of the entire list which is trivial on the
kernel of \(\omega\) fixes the eigenvalues at \(g_0\).
The four distinct labels force the permutation to be
the identity, and cancellation of any character in the
list shows that the self-twist is trivial. Moreover,
on \(G_{F_0}\) the root \((1,0)\) is precisely
\(\lambda_0\), including on the two cosets of
\(G_{F_0K'}\).

Take a rational inertia element \(u\) above \(\ell\)
whose \(\bar\omega\)-value generates \(\mathbf F_\ell^\times\).
Its square \(u_0\) belongs to \(H\cap J_{\mathrm{quad}}\).
At this element all diagonal residual characters are one, but
\[
                \operatorname{ord}\bar\omega(u_0)=(\ell-1)/2>1.
\]
Conjugating \(u_0\) gives this condition at
every marked place. It proves the comparison of each
residual root with the cyclotomic character and the
local condition at \(\ell\).
Chebotarev applied to \(F_0KK'\), the residual characters
and \(\bar\omega\) now gives arbitrarily large rational
primes split in these three fields, with trivial residual
diagonal and nontrivial cyclotomic value. We may also
apply Chebotarev to the identity class to get primes
split in \(F_0K\) with both residual diagonal and
cyclotomic value trivial. The two choices have distinct
uses in the ordinary comparisons and the residual chain.

\end{proof}

\subsubsection{Passage to deeper layers}

The component argument requires large degrees after tame base change. The marked elements constructed at the starting field must continue to distinguish the labels upon passage to these deeper layers.

\begin{proposition}[Fitted fields and primes for the four-link path]\label{prop:count-three-fitted-layers}
After adjoining the degree-\(\ell\) cyclotomic fields of the
added tame primes and passing to a cyclotomic layer, the degree
and local degree may be taken as large as required in
Proposition~\ref{prop:tower-component-propagation}. If \(C\) is the cohomological constant of the tame intermediate field, the eventual layer can in particular satisfy 
\[
 [F:\mathbb Q]\ge30(C+1)^2.
\]
 The elements \(g_0,u_0\) of Proposition~\ref{prop:count-three-marked-elements} retain their properties after taking suitable powers. The two ordinary
auxiliary primes at this stage use the cyclotomic value at a power of \(u_0\);
the extra conditions at the splitting prime \(w\) for the seed
are imposed separately.
\end{proposition}

\begin{proof}
\proofstep{Restrict powers of \(g_0,u_0\) to the deeper layers.}

The properties just proved persist in any abelian
\(\ell\)-power enlargement of \(F_0\) with the same \(K\).
Take a common \(\ell\)-power of \(g_0\) and \(u_0\)
which lies in the new group. If \(z\) is a residual character value of order \(e\) prime to \(\ell\), then
\[
 \operatorname{ord}(z^{\ell^a})=e/\gcd(e,\ell^a)=e .
\] This retains the marked tests. Splitting and
inertia restrictions persist by restriction. For
\(g_0\) the first root is still \(-1\), so the
quadratic restriction concerning \(K'\) continues to
give both cosets. At a finite set of extra split primes
\(w\equiv1\pmod\ell\), the compositum with the degree
\(\ell\) subfields of \(\Q(\zeta_w)\) absorbs every
\(\ell\)-torsion tame inertia quotient at those primes.
These cyclic fields have total ramification of degree
\(\ell\) at their own conductor and are unramified
at \(\ell\). The resulting intermediate field \(F_1\)
is still tame at \(\ell\).

The depth bound is chosen on \(F_1\) after the finite set
of extra primes has been specified. Use the cohomological
constant \(C\) for the cyclotomic \(\mathbf Z_\ell\)-layers
over this tame field. Then take a layer \(F\) with
\[
 d=[F:\Q]\geq30(C+1)^2.
\]
If \(\Q_m\) is the degree-\(\ell^m\) cyclotomic layer, then
\[
          [F_1\Q_m:F_1]=[\Q_m:\Q_m\cap F_1]\longrightarrow\infty .
\]
The denominator in this degree comparison stays bounded by
\([F_1:\Q]\).
This is the same order of choices as in the rank-four
component argument. In the rank-four case the prime-to-\(\ell\)
order of \(\xi\) already gives the deep character
condition at every layer. The two ordinary exceptional
primes can be chosen on \(FK\) using this power of \(u_0\). The splitting prime \(w\) used for the
seed instead needs the identity Frobenius class modulo
the residual and cyclotomic characters. Its additional
local conditions enter the tensor construction.
\label{cor:ordinary-count-iii-tower-1}
\end{proof}

\subsection{The two tensor factors and their integral flags}
\label{sec:count-three-tensor}

With the base field fixed, the four residual labels prescribe ratios for the two tensor factors. The integral structures on those factors will determine which extensions occur on reduction.

Use \(F_0,K,K',g_0\) from Section~\ref{sec:count-three-tower}.
Write \(\sigma=1\) for the original order and \(\sigma=-1\)
for the exchanged order. Let \(D_i\) denote the ordered
characteristic-zero characters of the diagonal rational datum,
and let \(d_i\) be their reductions. Their quotient
\(d_0/d_1\) on \(G_{F_0}\) is \(\lambda_0\).
This is the character of \(K'/\Q\), and it is nontrivial
at \(g_0\). In particular \(K'\) is contained in neither
\(F_0\) nor \(F_0K\), and both exclusions persist on the
\(\ell\)-power layers by taking an appropriate \(\ell\)-power of \(g_0\).

\subsubsection{The quadratic factor}

The first factor must distinguish the two orders of the tensor labels. Inducing an algebraic Hecke character from the auxiliary quadratic field prescribes both its inertia weights and its action at the chosen normalizer.

\begin{proposition}[The quadratic integral factor]\label{prop:count-three-quadratic-factor}
The algebraic Hecke character of \(K'\) given by the
class-number construction yields the integral rank-two
quadratic factor, with the specified diagonal action at \(g_0\). Its two inertia exponents at the chosen place above \(\ell\) are 
\[
 (m,-m).
\]

\end{proposition}

\begin{proof}
Let \(h'\) be the class number of \(K'\). For each fractional
ideal \(\mathfrak a\), choose a generator \(a\) for
\(\mathfrak a^{h'}\). For a positive multiple \(r\) of
the number of roots of unity of \(K'\), the character
\[
                 \mathfrak a\longmapsto (a/\bar a)^r
\]
is independent of the generator because changing \(a\) by a
root of unity multiplies this value by an \(r\)-th power equal
to one. Multiplication of the ideals multiplies the generators
up to such a unit. On a principal ideal \((x)\) its value is
\[
                       (x/\bar x)^{rh'}.
\]
These exponents are its infinity type.

It is unramified at every finite prime as a Hecke character
and conjugation inverts it. Its \(\ell\)-adic avatar
\(\kappa'\) is unramified away from \(\ell\).
Above \(\ell\), the two places are split, and on their
inertia groups the exponents are \(m_0,-m_0\), where
\(m_0=rh'\), up to choosing the first place.
The avatar takes values in the finite local coefficient
field given by the chosen embedding of \(K'\).
Its values are units by continuity and compactness.
Raise this character to a positive power to make its
reduction one and keep \(m_0\) for the resulting
positive exponent. The inertia exponents are then
still opposite, and conjugation still inverts the
character.

For any positive integer \(a\), take
\[
 \tau=\operatorname{Ind}_{G_{K'}}^{G_{\Q}} (\kappa')^a
\]

and write \(m=am_0\). Since \(g_0\notin G_{K'}\), conjugation by it inverts
\(\kappa'\) and fixes \(g_0^2\). Thus
\[
 \kappa'(g_0^2)=\kappa'(g_0^2)^{-1}.
\]
This value reduces
to one, so it is one because \(\ell\) is odd.
Consequently in the two induced basis vectors the action
of \(g_0\) interchanges the vectors. The change of basis to their sum and difference has matrix
\[
 S=\begin{pmatrix}1&1\\1&-1\end{pmatrix},\qquad
 \det S=-2\in O^\times,\qquad
 S^{-1}\begin{pmatrix}0&1\\1&0\end{pmatrix}S=\operatorname{diag}(1,-1).
\]
The new eigenbasis is therefore integral.
In this basis, on \(G_{K'}\) the matrices have the form
\[
                    \begin{pmatrix}u&v\\v&u\end{pmatrix},
       \qquad u^2-v^2=1.
\]
They preserve \(\operatorname{diag}(1,-1)\), and the
other coset acts by multiplying a matrix of this form by
\(\operatorname{diag}(1,-1)\).

In the residue field
the matrices are \(\operatorname{diag}(1,\lambda_0)\)
on the whole rational group. At a place above \(\ell\),
the two eigenvectors are \((1,s)\) and \((1,-s)\)
with \(s=1\) or \(-1\), chosen so that the first inertia
exponent is \(m\). Under conjugation of the place,
the sign is preserved on \(G_{K'}\) and is exchanged
on its other coset. These choices can all be made
over one finite coefficient DVR.

\end{proof}

\subsubsection{Determinant twist for the other factor}

The other rank-two factor will come from an ordinary form. Its determinant is constrained by the four target characters, which fixes the finite twist before the weight is chosen.

Put
\[
             \phi=\eta\xi^\sigma,\qquad
             \Phi=\epsilon_\ell^{k-1}\phi .
\]
The finite character \(\phi\) is odd and is unramified
at \(\ell\). First take a rational prime \(w\) split in \(F_0K\), with
\(w\equiv1\pmod{\ell\operatorname{cond}(\phi)}\) and trivial
residual diagonal. Such primes follow by Chebotarev at the
identity of the compositum. A character \(\theta\) of order
\(\ell\) modulo \(w\) is available and reduces to one.
It is killed on inertia by adjoining the degree-\(\ell\)
subfield of \(\Q(\zeta_w)\).
After fixing this level, take \(k\) odd with
\(k-1\equiv j\pmod{\ell-1}\).
This congruence has odd solutions because \(\xi(-1)=1\)
forces \(j\) even. The Chinese remainder theorem
also lets us impose \(\ell\nmid k\) and any lower
bound determined by the chosen level and \(m_0\).

\subsubsection{Divisibility of the Eisenstein constant}

A cuspidal eigenform with the desired residual system will be obtained by congruence with an Eisenstein series. The constant term must therefore vanish in the residue field of the chosen coefficient ring.

\begin{proposition}[Small Eisenstein constant at the auxiliary conductor]\label{prop:count-three-small-constant}
For \(\chi=\phi\theta\) at the selected prime \(w\), let \(a_0\) denote the constant of the weight-\(k\) Eisenstein series. In the coefficient DVR \(O\), Proposition~\ref{prop:count-three-cusp-congruence} uses the divisibility 
\[
 a_0=\tfrac12 L(1-k,\chi)\in\mathfrak m_O.
\]
 The calculation
uses the original fixed weight, before the final choice of a deep
tower layer.
\end{proposition}

\begin{proof}
Let \(M=\operatorname{cond}(\phi)\) and
\(\chi=\phi\theta\), primitive of conductor \(N=Mw\).
For the Eisenstein series with positive Fourier
coefficients \(\sum_{d\mid n}\chi(d)d^{k-1}\), the
constant term is \(L(1-k,\chi)/2\).
Denote by \(\chi_0\) the character on the same
modulus induced by \(\phi\), so the two characters
have the same reduction. We check that their
special values are integral and congruent.

For any of the characters \(\gamma\) of modulus \(N\)
or \(M\), grouping the Hurwitz zeta series gives
\[
 -k L(1-k,\gamma)=
  (\ell N_\gamma)^{k-1}
       \sum_{a=1}^{\ell N_\gamma}
          \gamma(a) B_k\bigl(a/(\ell N_\gamma)\bigr),
\]
where \(N_\gamma=N\) or \(M\) is the modulus used
in the grouping and \(B_k(X)\) is the Bernoulli
polynomial. Here \(\gamma(a)\) is taken modulo
\(N_\gamma\), even when \(\ell\mid a\).
Expand the polynomial in Bernoulli numbers and
group by \(a\bmod N_\gamma\).
For \(i\ge1\), the coefficient
\((\ell N_\gamma)^{i-1}B_i\) is \(\ell\)-integral by
the von Staudt--Clausen denominator bound and
the integrality of \(B_1=-1/2\) at odd \(\ell\).

For \(i=0\), the coefficient of \(\gamma(b)\)
is the sum of the \(k\)-th powers of all
\(\ell\) representatives of the class of \(b\),
divided by \(\ell N_\gamma\).
The representatives give each residue modulo \(\ell\) once.
Because \(k\) is odd,
\[
          \sum_{r\in\mathbf F_\ell}r^k=0
\]
by pairing \(r\) with \(-r\). This cancels the remaining
factor \(\ell\) in the denominator.
Dividing by \(k\) is allowed as \(\ell\nmid k\),
and at the common modulus \(N\) each coefficient
in this linear expression is independent of
the character. This proves the congruence
of the two special values and integrality of
the special value for \(\phi\).
Finally,
\[
 L(1-k,\chi_0)
     =(1-\phi(w)w^{k-1})L(1-k,\phi)
\]
has positive \(\ell\)-adic valuation. The first
factor is zero modulo \(\ell\) by the choice of
\(w\). Hence the Eisenstein constant belongs to
the maximal ideal of the coefficient DVR.
The weight congruence and this additional choice
of \(k\) leave the eventual lower bound on \(k\) free.
\hfill\mbox{\cite{washington-cyclo,bernoulli-values}}\label{cor:ordinary-count-iii-seed-1}

\end{proof}

\subsubsection{An integral form with unit constant}
\label{par:count-three-cusp-correction}

Divisibility of one Eisenstein constant is only the first step toward a cusp form. An integral correcting form must change that constant while prescribing its values independently at the other cusps.

\begin{proposition}[Cusp correction with unit constant term]\label{prop:unit-constant-cusp-correction}
Let \(N\ge3\), \(\ell\) be odd and prime to \(N\), and
\(\chi(-1)=(-1)^k\). With the separating weight \(m\) of
\eqref{eq:count-three-correction-degree}, assume \(k\ge m+10\).
For the correcting form \(F\), let \(a_n(F;c)\) denote the expansion coefficients at a cusp \(c\), in its chosen integral coordinates. Then 
\[
 a_n(F;c)\in O\quad\text{at every cusp},\qquad a_0(F;\infty)=1,
\]
 with \(O\) the coefficient DVR. This is
the correction used for the primitive Eisenstein series.
\end{proposition}

\begin{proof}
Let \(N\geq3\) with \(\ell\nmid N\), and let
\(\zeta\) be a primitive \(N\)-th root of unity. We use its complex
value for Fourier expansions and its \(\ell\)-adic value for
integrality. For \(a\in(\Z/N)^\times\) and \(r\geq3\), take the
normalized weight-\(r\) Eisenstein sum over \((c,d)\equiv(0,a)\bmod N\)
with summand \((cz+d)^{-r}\). Denote it by \(E_{r,a}\). The normalization is
\[
 E_{r,a}(z)=\frac{N^r(r-1)!}{(-2\pi i)^r}
          \sum_{(c,d)\equiv(0,a)\bmod N}(cz+d)^{-r}.
\]
Absolute convergence gives holomorphy and the right slash action of
\(\gamma=\bigl(\begin{smallmatrix}*&*\\c_\gamma&d_\gamma\end{smallmatrix}\bigr)\)
transports the row \((0,a)\) to \((ac_\gamma,ad_\gamma)\).

In particular the form is modular on \(\Gamma_1(N)\). For
\(\gamma\in\Gamma_0(N)\), the slash is \(E_{r,ad_\gamma}\).
If \(c_\gamma\not\equiv0\bmod N\), the transported sum contains
only terms with \(c\ne0\). Summing the absolute values over \(d\)
gives a bound \(O_r((|c|\operatorname{Im} z)^{1-r})\) for
\(\operatorname{Im} z\geq1\). Writing \(y=\operatorname{Im}z\ge1\) gives the explicit tail bound
\[
 \sum_{c\ne0}(|c|y)^{1-r}
 =2y^{1-r}\sum_{c\ge1}c^{1-r}\longrightarrow0
 \qquad(y\to\infty),
\]
since \(r>2\). The same bound on the
nonconstant part and the convergent \(c=0\) sum give holomorphy
at every cusp.

Set \(L_r(x)=(x\,d/dx)^{r-1}(x/(1-x))\). Differentiating the
partial fraction expansion of \(\pi\cot(\pi t)\) gives
\[
  \frac{(r-1)!}{(-2\pi i)^r}
      \sum_{b\in\Z}(t+b)^{-r}
     =L_r(e^{2\pi i t}).
\]
For \(\operatorname{Im}t>0\) the right side has expansion
\(\sum_{e\geq1}e^{r-1}e^{2\pi i et}\). For the constant
term take \(t=a/N\). For \(c>0\) use \(t=(cz+a)/N\).
The terms with \(c<0\) give the contribution with \(-a\)
and the factor \((-1)^r\). Thus at infinity
\[
 a_0(E_{r,a})=L_r(\zeta^a),\qquad
 a_n(E_{r,a})=\sum_{e\mid n}e^{r-1}
              \bigl(\zeta^{ae}+(-1)^r\zeta^{-ae}\bigr)
       \quad(n\geq1).
\]
The numerator of \(L_r\) is integral and its denominator is
\((1-x)^r\). Reduction is injective on the \(N\)-th roots of
unity, so all these coefficients are integral at \(\ell\).

We next separate the constant at \(a=1\) from that at any
\(a=d\ne\pm1\), with the same weight in the two evaluations.
In the residue field put \(t=x/(1-x)^2\) and
\(u=(1+x)/(1-x)\) for \(x=\bar\zeta^a\). Direct differentiation gives
\[
 L_3(x)=tu,\qquad L_4(x)=t(1+6t),\qquad
 L_5(x)=tu(1+12t),\qquad u^2=1+4t.
\]
At any unit label \(L_3(x)\ne0\), because \(x\ne1,-1\).

In residue characteristic other than two or three, equality of
both \(L_3^2\) and \(L_3L_5\) at \(x,y\) forces
\(t(x)=t(y)\) by taking their quotient.
In characteristic three \(L_4=t\) suffices. Clearing denominators
in \(t(x)=t(y)\) gives \((x-y)(1-xy)=0\).
For \(x=\bar\zeta\), \(y=\bar\zeta^d\), both alternatives
are excluded. At least one of the three families
\[
      E_{4,a},\qquad E_{3,a}^2,\qquad E_{3,a}E_{5,a}
\]
therefore has constants at \(a=1,d\) differing by a unit.

For each such \(d\) choose one family \(H_{d,a}\) of weight \(r_d\).
For a level-one form \(L_{r_d}\) of weight \(r_d\)
with integral coefficients and constant one, set
\[
        H'_{d,a}=H_{d,a}-a_0(H_{d,d})L_{r_d}.
\] Such
forms exist for all even nonnegative weights except two,
by products of \(E_4\) and \(E_6\), including the empty product.
The resulting family has zero constant at \(a=d\) and unit
constant at \(a=1\). Multiplication of the correcting factors gives the even weight
\begin{equation}\label{eq:count-three-correction-degree}
                 m=\sum_{d\ne\pm1}r_d\le8\varphi(N).
\end{equation}
It has zero constants at all \(a\ne\pm1\) and retains a
unit at \(a=1\). The same right slash rule holds for all the
families, since the coefficients in these operations are fixed
as \(a\) varies.

For \(k\geq m+10\), multiply this family by \(E_{3,a}\) if
\(k\) is odd and by \(E_{3,a}^2\) if \(k\) is even. Complete the weight with a product of \(E_4\) and \(E_6\) having constant term one.
The forms \(F_a\) of weight \(k\) so obtained vanish at all
cusps outside the \(\Gamma_0(N)\)-orbit of infinity by the
weight-three factor. At infinity \(c_1=a_0(F_1)\) is a unit. The constant terms
are
\[
       a_0(F_a)=
       \begin{cases}
       c_1,&a=1,\\
       (-1)^kc_1,&a=-1,\\
       0,&a\ne\pm1 .
       \end{cases}
\]
The second case is obtained by slashing by \(-I\).

For \(\chi(-1)=(-1)^k\), set
\[
    G=\frac{1}{2c_1}\sum_{a\in(\Z/N)^\times}\chi(a)^{-1}F_a .
\]
The right slash rule and reindexing \(a\mapsto ad_\gamma\)
give \(G|_k\gamma=\chi(d_\gamma)G\) on \(\Gamma_0(N)\).
Its coefficients are integral because \(2c_1\) is a unit.
Its constant is one, and it vanishes away from the orbit of
infinity. In particular the construction works at the level
\(N=Mw\) with \(w\equiv1\bmod\ell\), even though
\(\ell\mid\varphi(N)\).
\label{cor:ordinary-count-iii-eigenform-1}

\end{proof}

\subsubsection{The cuspidal congruence}

The integral correction can now be multiplied by the small Eisenstein constant. Subtracting it enforces cuspidality while preserving the positive Fourier coefficients modulo the maximal ideal.

\begin{proposition}[A cuspidal congruence at the chosen conductor]\label{prop:count-three-cusp-congruence}
For the primitive Eisenstein character \(\chi\) of conductor \(N=Mw\)
and the small constant calculated for the seed, the cusp correction gives an integral cusp form \(f\) satisfying 
\[
 a_n(f)\equiv\sum_{e\mid n}\chi(e)e^{k-1}\pmod{\mathfrak m}\qquad(n\ge1).
\]
 The
congruence includes the ordinary Hecke eigenvalue at \(\ell\)
used in the lift.
\end{proposition}

\begin{proof}
For the primitive character \(\chi=\phi\theta\) of conductor
\(N=Mw\), the weight-\(k\) Eisenstein series \(E\) of character
\(\chi\) is the sum of \(E_{k,a}\) with weights
\(\chi(a)^{-1}/(2\mathcal G(\chi^{-1}))\) over the unit labels.
Here \(\mathcal G(\chi^{-1})=\sum_a\chi(a)^{-1}\zeta^a\)
is the nonzero primitive Gauss sum, with characters extended
by zero on nonunits. Primitivity gives
\(\sum_a\chi(a)^{-1}\zeta^{ae}=\chi(e)\mathcal G(\chi^{-1})\)
also when \(e\) is a nonunit.

For a nonunit, multiplying
\(a\) by an element fixing \(e\) and on which \(\chi\) is
nontrivial gives zero by primitivity. Parity then gives
\[
    a_n(E)=\sigma(n):=\sum_{e\mid n}\chi(e)e^{k-1}
    \quad(n\geq1),\qquad a_0(E)=L(1-k,\chi)/2 .
\]
For the constant, the same Gauss-sum identity applied
to \(\sum_{e\geq1} e^{k-1}\zeta^{ae}x^e\) for \(|x|<1\)
gives the finite sum of \(L_k(\zeta^a x)\).

At \(x=e^t\), the generalized Bernoulli generating
function gives the continued value
\[
 \left.\left(\frac{d}{dt}\right)^{k-1}
            \frac{\sum_{e=1}^N\chi(e)e^{et}}{1-e^{Nt}}
       \right|_{t=0}=L(1-k,\chi).
\]
The primitive Gauss factor and parity give the normalization of
the constant term above. The cusp vanishing and the character follow
from the corresponding transformations of \(E_{k,a}\).
Set \(f=E-a_0(E)G\). Its constant term at infinity is zero,
and the other positive coefficients satisfy
\[
      a_n(f)\equiv a_n(E)=\sigma(n)\pmod{\mathfrak m}\qquad(n\ge1)
\]
because \(a_0(E)\in\mathfrak m\) and \(G\) has integral coefficients. At cusps in the orbit of infinity,
the vanishing follows by applying the common character law
to the zero constant of \(f\) at infinity. Here the smallness
of \(a_0(E)\) is the earlier Bernoulli and Euler-factor
calculation for the chosen \(k\).

For a prime \(p\), the Hecke operator in this character space has
coefficient \(a_{pn}+\chi(p)p^{k-1}a_{n/p}\), where the second
term is zero if \(p\nmid n\), and \(\chi(p)=0\) if \(p\mid N\).
The divisor sums satisfy the same formula with eigenvalue
\(1+\chi(p)p^{k-1}\). In fact, if \(n=p^b n'\) with
\(p\nmid n'\) and \(v=\chi(p)p^{k-1}\), this is the polynomial
identity
\[
  (1+v+\cdots+v^{b+1})+v(1+v+\cdots+v^{b-1})
      =(1+v)(1+v+\cdots+v^b),
\]
with the second sum empty at \(b=0\).
Thus the reduction of \(f\) is a simultaneous Hecke
eigenvector, nonzero since \(a_1(f)\equiv1\).

\end{proof}

\subsubsection{Simultaneous lifting and the ordinary line}

A nonzero mod-\(\ell\) cusp form gives the desired residual eigenvalues. The second tensor factor requires one characteristic-zero eigenform realizing them simultaneously, including the ordinary eigenvalue at \(\ell\).

\begin{proposition}[Eigensystem of the corrected cusp form]\label{prop:count-three-cuspidal-eigensystem}
The cuspidal congruence of Proposition~\ref{prop:count-three-cusp-congruence}
gives the simultaneous cuspidal eigenform \(g\) and the ordinary local line used for the rank-two factor. After a finite coefficient extension the normalized eigenform satisfies 
\[
 a_p(g)\equiv1+\chi(p)p^{k-1}\pmod{\varpi}\qquad(p\text{ prime}).
\]
 The weight can be
fixed before the choice of the final deep layer.
\end{proposition}

\begin{proof}
\proofstep{An ordinary eigensystem with the selected residual line.}

Choose a finite local coefficient field \(E_\ell\) containing the
coefficient fields of the forms in the construction, and let \(O\)
be its valuation ring with uniformizer \(\varpi\). The
character-\(\chi\) cusp space of level \(N\) has a model over
this field. Indeed normalized newforms have algebraic
coefficients in number fields, and the newform decomposition
spans the cusp space. Since the character has conductor \(N\),
every constituent is new at level \(N\). Both the character
and the coefficients of \(f\) are in the chosen coefficient
field, so \(f\) belongs to this model. One can check this by
taking finitely many Fourier coefficients that identify the
coordinates in a newform basis.

Let \(V_O\) consist of the forms with all coefficients in \(O\).
Take a finite list \(n_1,\ldots,n_s\) separating the forms. Its coefficient map gives
\[
 V_O\hookrightarrow O^s,\qquad h\longmapsto
 (a_{n_1}(h),\ldots,a_{n_s}(h)).
\]
A submodule of a finite free module over a DVR is finite free.
It also contains a basis after integral scaling, by
integrality of the normalized newforms. It is finite free
over \(O\). The coefficient formula makes it stable under
all prime Hecke operators, including those at \(p\mid N\). The coefficient condition gives the saturated reduction
\[
       \ker(V_O\longrightarrow \text{residue expansions})
                                      =\varpi V_O.
\] The mod-\(\varpi\) eigenline of \(f\) therefore
defines a character of the commutative \(O\)-algebra of
Hecke operators on \(V_O\). Write \(\mathfrak m\) for its
kernel. This algebra is finite and torsion-free over \(O\).
Write \(\mathbb T\) for the Hecke algebra. A minimal
prime \(\mathfrak p\subset\mathfrak m\) avoids the uniformizer
since it is a nonzerodivisor.

Normalize \(\mathbb T/\mathfrak p\)
in its fraction field and choose a place over \(\mathfrak m\);
this gives one characteristic-zero evaluation
\[
 \mathbb T\longrightarrow
      \operatorname{Frac}(\mathbb T/\mathfrak p)
\]
with all the prescribed residual Hecke values.
In the generic fibre the faithful module has a nonzero
factor at \(\mathfrak p\). After extending scalars, the
maximal ideal of that Artinian local factor acts nilpotently.
If \(M_\mathfrak p\) is this local generic factor and
\(\mathfrak n\) its maximal ideal, choose the last nonzero
\(\mathfrak n^a M_\mathfrak p\). It is killed by \(\mathfrak n\),
so a nonzero vector there is a common eigenvector.
The eigenvector has nonzero first coefficient, since the
Hecke recursion from a zero first coefficient would give
zero coefficients at every positive index. Normalizing it
gives a cuspidal eigenform \(g\) of level \(N\) and character
\(\chi\) with
\[
    a_p(g)\equiv1+\chi(p)p^{k-1}\pmod{\varpi}
       \qquad(p\text{ prime}),
\]
where \(\varpi\) now denotes a uniformizer after coefficient
extension. This is the simultaneous lifting argument of
Deligne and Serre. \cite[Lemma 6.11]{deligne-serre}

Use the Galois representation of the weight-\(k\) newform
\(g\), with arithmetic Frobenius polynomial
\(X^2-a_p(g)X+\chi(p)p^{k-1}\) for \(p\nmid\ell N\).
It is absolutely irreducible, with determinant
\(D=\Phi\theta\). Since \(\theta\) has order \(\ell\), its residual character is trivial. At every good prime the trace congruence and determinant now give
\[
 \bar P_p(X)=X^2-(1+\bar\Phi(\operatorname{Frob}_p))X+
 \bar\Phi(\operatorname{Frob}_p)
 =(X-1)(X-\bar\Phi(\operatorname{Frob}_p)).
\]
Chebotarev and semisimplification yield \(\bar\Phi\oplus1\). For \(p\mid N\) the local component has conductor
exponent equal to that of the central character, by primitivity
of \(\chi\).

It is principal series with one unramified factor.
Indeed, for a supercuspidal parameter the two-dimensional
irreducible Weil representation has invariants of dimension
zero or two on each normal ramification group. Its determinant
has conductor at most half as large by the Artin conductor
formula. For a special representation the conductor is one
with unramified central character if the twist is unramified.
For a ramified twist of conductor exponent \(a>0\) the
conductor is \(2a\), while that of its central character is
at most \(a\). For two principal series characters with conductors
\(a_1,a_2\), the conductor and determinant bounds read
\[
                a(W)=a_1+a_2,\qquad a(\det W)\le\max(a_1,a_2).
\]
At a ramified prime equality thus forces precisely one
\(a_i\) to be zero. Local--global
compatibility gives the two inertia characters \(1,\phi\theta\).
Outside \(\ell N\) the representation is unramified.
\cite{deligne-galois,carayol-local}

Since \(k\geq3\) and \(\ell\nmid N\), the congruence at
\(p=\ell\) gives \(a_\ell(g)\equiv1\). The unit root
\(\alpha_\ell\) of the Hecke polynomial at \(\ell\) is
congruent to one. The ordinary local theorem for this
newform, or for its unit-root \(\ell\)-stabilization, gives
an unramified quotient \(\nu\) with
\(\nu(\operatorname{Frob}_\ell)=\alpha_\ell\) and invariant
line of character \(D\nu^{-1}\). Its inertia character is
\(\epsilon_\ell^{k-1}\). The two generators of its local data satisfy
\[
         \bar\nu|_{I_\ell}=1,\qquad
         \bar\nu(\operatorname{Frob}_\ell)=1 .
\]
Hence \(\bar\nu\) is trivial on the decomposition group. The high-weight line thus reduces to \(\bar\Phi\)
on that entire group, including when
\(k-1\equiv0\bmod\ell-1\).
For the dual of the attached representation the
zero-exponent line has residual character one and determinant
\(D^{-1}\); the twist by \(D\) gives exactly the preceding
order of the characters. \cite{hida-ordinary-galois}\label{cor:ordinary-count-iii-eigenform-2}

The weight was chosen after \(N\) and the initial positive
exponent \(m_0\) of the quadratic factor. In addition to
the cusp-correction cutoff \(m+10\), we may impose
\(k\geq18m_0+3\) and the earlier congruence and coprimality
conditions. No bound on the eventual deep layer is used in
this choice.
\end{proof}


\subsubsection{The integral lattice and its ordinary line}

The lift and twist fix the rational determinant. The residual endpoint depends also on the lattice in this representation, in particular on which ordinary line survives in its reduction.

\begin{proposition}[The ordinary lattice of the rank-two factor]\label{prop:count-three-ordinary-lattice}
The cuspidal lift and determinant twist give an integral rank-two representation \(r\) with 
\[
 \det r(g)=\Phi(g)\qquad(g\in\ker\theta).
\]

One can choose its lattice and the marked eigenbasis so that
the ordinary local line has the reduction and ordered
character required for the tensor.
\end{proposition}

\begin{proof}
\proofstep{Choose the ordinary integral lattice.}

The integral rank-two datum we now construct is
\(\rho_1\) with the following properties. Its determinant
is \(\Phi\) on \(\ker\theta\). At \(g_0\) it is diagonal,
and its reduction on \(G_{\Q}\) is
\[
                  \begin{pmatrix}\bar\Phi&0\\b&1\end{pmatrix}
\]
with \(b\ne0\). Away from \(\ell\) it is the identity on
inertia where both \(\phi\) and \(\theta\) are trivial.
At each chosen place \(A\) above \(\ell\), it has an
integral invariant line generated by \((1,x_A)\),
with ordered inertia characters \(\epsilon_\ell^{k-1},1\).
Start with the cuspidal attachment just constructed. Its
ordinary invariant line is the character \(D\nu^{-1}\), with
residue \(\bar\Phi\) on the entire decomposition group.

Choose a stable lattice and use the two eigenprojectors
of \(\rho(g_0)\). If the two eigenvalues are \(\lambda_1,\lambda_2\), the projectors are
\[
 e_i=\frac{\rho(g_0)-\lambda_j I}{\lambda_i-\lambda_j}
 \qquad(\{i,j\}=\{1,2\}).
\]
The denominator is a unit by residual separation, so both preserve the lattice. Order the basis
with the \(\bar\Phi(g_0)\)-eigenvalue first. The traces
of \(\rho(g)\) and of \(\rho(g_0g)\) determine its
two reduced diagonal entries as \(\bar\Phi(g),1\).
The first diagonal entry in a product is
\(a_ga_h+b_gc_h\) if \(b_g,c_h\) are the off-diagonal
entries. Comparing with the reduced diagonal character gives
\[
                  \overline{b_gc_h}=0\qquad(g,h\in G_\Q).
\]

The two sets of off-diagonal entries have finite
minimum valuations \(v_+,v_-\), by irreducibility,
and \(v_++v_->0\). Multiplying the second basis vector by \(\varpi^{v_-}\)
changes the minima to
\[
           v'_-=v_--v_-=0,\qquad v'_+=v_++v_->0.
\]
Thus a lower-left entry survives and the upper-right ones all
vanish after reduction.
At each place over \(\ell\), a rational Frobenius
lift with trivial cyclotomic value separates \(\bar\Phi\)
from one, by Paragraph~\ref{par:dyadic-residual}.
The reduced ordinary line is the first, \(\bar\Phi\)-eigenline there. A generator \((u,v)\) of the lifted line has \(u\) a unit and admits the normalization
\[
 u^{-1}(u,v)=(1,x_A),\qquad x_A=v/u.
\]


The cuspidal rank-two representation is irreducible
on every open subgroup. Indeed, a reducible
restriction to an open normal subgroup that contains
a power of an inertia element with distinct eigenvalues
has at most two invariant lines. Such an element
exists here since the local cyclotomic exponents are
distinct. By irreducibility on the rational group
there must be two lines, exchanged by the other
coset of their stabilizer. On the stabilizer,
the two reduced line characters are \(1,\bar\Phi\).
Conjugation exchanges them and fixes \(\bar\Phi\),
so \(\bar\Phi\) is one on the stabilizer. Writing \(H\) for the index-two stabilizer, each square belongs to \(H\). Thus
\[
 \bar\Phi(g)^2=\bar\Phi(g^2)=1\qquad(g\in G_\Q).
\] At the mark this contradicts
\(\bar\Phi(g_0)^2=z^2\ne1\). The latter follows from
\(\bar\Phi(g_0)=-z\), with \(z\) of nontrivial odd
order. The argument for arbitrary open subgroups
follows by passing to the normal core.

Put \(u=(\epsilon_\ell/\omega)^{1/2}\), the root with
values in \(1+\ell\Z_\ell\), and define
\[
 \delta =
 \begin{cases}
 u^{j+2-k} & \sigma=1,\\
 \omega^j\xi\,u^{j+2-k} & \sigma=-1 .
 \end{cases}
\]
The identities \(\omega^{k-1}=\omega^j\) and
\(\omega^{2j}=1\) give
\(\delta^2\Phi=\vartheta\).
For the original order, \(\bar\delta=1=d_3\) and
\(\bar\delta\bar\Phi=d_0=\bar\vartheta\).

For the exchanged order, \(\bar\delta=\bar\Psi=d_3\)
and \(\bar\delta\bar\Phi=d_0=\bar\alpha\).
Together with \(d_0/d_1=\lambda_0\) and
\(d_0d_3=d_1d_2\), this identifies all four reductions.
On any inertia group away from \(\ell\) where the
residual diagonal is trivial, \(\omega=1\) and the
values of both finite characters \(\eta,\xi\) are
one by prime-to-\(\ell\) injectivity. Thus
\(\delta=\phi=1\) there.
The same conclusions hold on the part of inertia
above \(\ell\) in the cyclotomic kernel, since
\(\eta,\xi\) are unramified above \(\ell\).

\end{proof}

\subsubsection{Tensor representation and residual endpoint}

We can now tensor the two integral bases in the marked order. This realizes the endpoint as an actual residual representation with a flag, data that a list of four characters alone would not specify.

\begin{proposition}[Integral tensor and its marked endpoint]\label{prop:count-three-tensor-endpoint}
With the ordinary bases of the two factors, their twisted tensor
is symplectic and has the required ordered flag on the fitted
field. Its reduction gives the designated endpoint and the local diagonal characters of the rank-four chain. Write \(b_i\) for the inertia exponent on the \(i\)th ordered line after the scalar twist. Then 
\[
 2(b_1,b_2,b_3,b_4)=2(k-1+m,k-1-m,m,-m)+(j+2-k)(1,1,1,1).
\]

\end{proposition}

\begin{proof}
Let \(E\) be the free rank-two module of \(\rho_1\), with
the standard alternating form, and \(H\) that of \(\tau\)
with symmetric form \(\operatorname{diag}(1,-1)\).
For bases \((e_0,e_1)\), \((f_0,f_1)\), take the tensor
basis
\[
               (e_0 f_0,e_0 f_1,-e_1 f_1,e_1 f_0).
\]
The tensor form is alternating and has the prescribed
pairings in indices \(0,3\) and \(1,2\). In this basis,
\(\delta(\rho_1\otimes\tau)\) has similitude
\(\delta^2\det\rho_1\). Thus after killing \(\theta\)
on the deep layer its similitude is exactly
\(\vartheta\). On the tame layer the same equality
holds after reduction, since \(\theta\) reduces to one.
Both factors are diagonal at the normalizer and all
its powers.

\proofstep{Compute the tensor flag in the ordered basis.}

For the local flag let \(L_x=\begin{pmatrix}1&0\\x&1\end{pmatrix}\).
Conjugate the tensor first by \(L_{x_A}\otimes1\) in
the displayed basis. Then conjugate by
\[
             1+s_A E_{10}-s_A E_{32},
\]
where \(s_A\) is the sign of the chosen eigenline of
\(\tau\) at \(A\). Both are lower unipotent symplectic
matrices. Their product is the flag matrix, with
\((1,0)\) entry \(s_A\).
Over characteristic zero the local conjugate is upper
triangular, since the first step triangularizes the
rank-two factor and the second orders the two lines
of the symmetric factor in each block. Its inertia
characters, in order, are
\[
       \delta\epsilon_\ell^{k-1+m},\quad
       \delta\epsilon_\ell^{k-1-m},\quad
       \delta\epsilon_\ell^m,\quad
       \delta\epsilon_\ell^{-m}.
\]

On the cyclotomic kernel the characters are all one.
They reduce to the required characters on the full
marked decomposition groups. To see the last point,
the marked place splits in \(K'\), so \(d_0=d_1\)
and \(d_2=d_3\) there. Conjugation by \(L_{x_A}\) preserves the residual lower
triangular order. Invariance of the first line then forces its
lower-left entry to vanish, so its reduction remains
\[
                          \operatorname{diag}(\bar\Phi,1).
\]
This proves the local diagonal conditions.

Outside \(\ell,w\), the residual tensor kills the
required inertia. Indeed where the original
residual diagonal is trivial, \(\lambda_0=1\) on
the restricted group. There \(\tau\) is unramified,
and \(\phi=\delta=\theta=1\) on the indicated inertia.
On the deep layer the corresponding integral claim
follows because the extra inertia has been killed.
The construction gives the tame scout \(S\) and
the compatible integral matrices on every such
layer, since they all come by restriction from the
same rational tensor. The chosen frames over
different layers transport as the lower factors
for this same ordinary flag. In the quadratic
factor the place sign changes by \(\lambda_0\)
on conjugation, and the rank-two line transports
by \(\rho_1\). These two actions give the flag
transport simultaneously on both factors.

In the residue field, the matrix is
\[
       \bar\delta\,
       \begin{pmatrix}
        \bar\Phi&0&0&0\\
        0&\lambda_0\bar\Phi&0&0\\
        0&-\lambda_0 b&\lambda_0&0\\
        b&0&0&1
       \end{pmatrix}.
\]
Thus only \(-w_\beta,-w_\vartheta\) can occur
globally off the diagonal. The lower frame entry
\(s_A\) stays nonzero. If \(b\) became zero on
\(G_F\), then conjugating the powered normalizer
by any rational Galois element would give
\[
           0=b(g)\bigl(\bar\Phi(g_0^{\ell^r})-1\bigr).
\]
This uses the product rule
\(b(gh)=b(g)\bar\Phi(h)+b(h)\), and the powered
normalizer still lies in \(G_F\) with
\(\bar\Phi\ne1\). Indeed \(\bar\Phi\) at the normalizer
is the ratio \(d_0/d_3\). Hence \(b\) is nonzero
on \(G_F\), and both indicated global roots occur.
The tensor supplies exactly the sign and support
conditions of Section~\ref{sec:count-three-chain}.

Finally, the doubled inertia exponents of this
tensor are the four integers
\[
 (2(k-1+m),2(k-1-m),2m,-2m)+(j+2-k)(1,1,1,1).
\]
Their adjacent differences are
\(4m,2(k-1-2m),4m\).
They are positive even integers at least four if
\(k\ge2m+3\). Removing the scalar by twisting
by \(\delta^{-1}\) changes the doubled inertial
exponent by \(-(j+2-k)\). Its sum with the first
doubled weight is even, as needed for the
ordinary seed comparison.
\hfill\mbox{\cite{ray-class}}\label{cor:ordinary-count-iii-seed-2}

\end{proof}

\subsubsection{Irreducibility after cyclotomic restriction}

To obtain an automorphic starting point we still need irreducibility after cyclotomic restriction and regularity of the weights after opposite shifts of the quadratic factor.

\begin{proposition}[Irreducibility and separation of infinity exponents]\label{prop:count-three-tensor-irreducible}
If the rank-two cuspidal factor is absolutely irreducible on all
open subgroups, the tensor restricted to the fitted CM layer
remains absolutely irreducible after killing the cyclotomic
character. For \(k\ge18m_0+3\) there is \(1\le a\le9\) for which the
opposite shifts \(\pm am_0\) give four distinct infinity parameters
and 
\[
 k-1-2am_0\ge2.
\]

\end{proposition}

\begin{proof}
Suppose the rank-two cuspidal input is absolutely
irreducible on all open subgroups of \(G_{\Q}\).
Put \(L=FK\) and \(M=LK'\).
Let \(H=G_M\cap\ker\epsilon_\ell\). If \(H\) contains a
nonscalar matrix, an invariant line for \(H\) belongs to the
at most two eigendirections of that matrix. For an invariant line \(l\), normality of \(H\) makes every \(G_M\)-translate an \(H\)-invariant line. Hence
\[
 [G_M:\operatorname{Stab}(l)]=|G_M l|\le2 .
\]
Continuity makes the stabilizer closed, so this finite-index subgroup is open and preserves \(l\).
If the subgroup acts by scalars, the projective
action factors through the procyclic compact
cyclotomic image of \(G_M\). A lift of a
topological generator preserves an eigenline,
and so does the closure of its projective cyclic
group. In both cases this contradicts the assumed
strong irreducibility.

The character of \(M/L\) is nontrivial on
\(G_L\cap\ker\epsilon_\ell\). Otherwise this quadratic
character would factor through \(\epsilon_\ell(G_L)\)
and kill its pro-\(\ell\) part. The powered normalizer
lies in \(G_L\) with \(\omega=1\) and
\[
 \lambda_0=-1,
\]
a contradiction.
Over \(G_M\cap\ker\epsilon_\ell\), the tensor without
its scalar is the sum of the two irreducibles
\(\rho\otimes(\kappa')^a\) and
\(\rho\otimes(\kappa')^{-a}\).
They are inequivalent. Choose an inertia element
\(i\in G_M\) above \(\ell\) with infinite cyclotomic
order, and let \(g\in G_L\) be the powered
normalizer outside \(G_M\). The commutator belongs to \(G_M\) by normality and to \(\ker\epsilon_\ell\) by abelianity. Since conjugation by \(g\) inverts \(\kappa'\), its value is
\[
 \kappa'(igi^{-1}g^{-1})
 =\kappa'(i)\kappa'(gi^{-1}g^{-1})
 =\kappa'(i)^2 ,
\]
of infinite order. On this commutator the ratio of the determinants of
the two rank-two summands is
\[
 \frac{\det(\rho\otimes(\kappa')^a)}
           {\det(\rho\otimes(\kappa')^{-a})}
       =(\kappa')^{4a},
\]
whose value there is
\[
 \kappa'(i)^{8a}\ne1.
\]

Thus the summands are inequivalent. The restriction to
\(G_M\cap\ker\epsilon_\ell\) is a direct sum with two
inequivalent simple summands, so a nonzero proper
invariant subspace for that subgroup must be one summand.
The other coset in \(G_L\cap\ker\epsilon_\ell\) exchanges
them, so neither is invariant for the whole group.
This also proves irreducibility on \(G_L\).
Scalar twisting has no effect on these conclusions.

For the infinity exponents in the induction
argument, the two conjugate places of \(K'\)
have two rank-two parameters each. Twisting
by \((\kappa')^a\) shifts these two parameter sets
by opposite integers \(am_0\) and \(-am_0\).
Only differences of the original parameters
in the two sets can equal \(2am_0\).
Writing the unshifted entries as \(b_1,b_2\) and \(c_1,c_2\), an overlap in either orientation requires
\[
  2am_0\in\{\pm(c_j-b_i):1\le i,j\le2\}.
\]
This set has at most eight elements, whereas \(2am_0\) takes nine distinct values in the indicated range. One of \(a=1,\ldots,9\)
therefore separates the four parameters. For the adjacent gaps,
\[
                       k-1-2am_0\ge2
\]
follows from the earlier choice \(k\ge18m_0+3\).
\label{cor:ordinary-count-iii-seed-3}

\end{proof}

\subsubsection{Cuspidal induction on the fitted CM layer}
\label{par:count-three-cuspidal-induction}

The separation of the infinity exponents is used here to obtain a regular cuspidal representation. Its restriction to the fitted CM field supplies the automorphic realization of the integral tensor.

\begin{proposition}[Cuspidal induction of the separated tensor]\label{prop:count-three-cuspidal-induction}
For the chosen infinity separation and the tensor of
Propositions~\ref{prop:count-three-tensor-endpoint}
and \ref{prop:count-three-tensor-irreducible}, the restriction to the fitted CM layer \(L\) has the regular algebraic polarized cuspidal attachment used for the seed. For the separated tensor before the common scalar twist, the ordered infinity labels and their consecutive gaps are 
\[
 (k-1+m,k-1-m,m,-m),\qquad(2m,k-1-2m,2m).
\]

\end{proposition}

\begin{proof}
With \(a\) chosen by the infinity separation, put
\(L=FK\) and \(M'=LK'\). The latter extension has degree
two over \(L\) because the powered normalizer belongs
to \(G_L\) and acts nontrivially on \(K'\). The extension
\(M'/K'\) is solvable, since both \(L\) and \(K'\) are
abelian over \(\Q\). Let \(\pi\) be the cuspidal representation
of \(\mathrm{GL}_2(\mathbb A_\Q)\) of the newform \(g\).
Base change to \(K'\), twist by the Hecke character of
\((\kappa')^a\), and base change to \(M'\). These are
cuspidal base changes.

At each cyclic step, failure of
cuspidality would give a nontrivial self-twist by a character
of the cyclic extension. Comparison at unramified places would give a Galois self-twist by a nontrivial quotient character \(\nu\). An intertwiner \(T\) restricts on the upper subgroup \(H\) to an element of its scalar commutant. Hence
\[
 T=cI,\qquad T r(g)T^{-1}=\nu(g)r(g)
 \quad\Longrightarrow\quad \nu(g)=1\quad\text{for all }g,
\]
which is impossible. This holds by the rank-two strong irreducibility
and is preserved by the character twist. At unramified places, base change takes the powers
dictated by the residue degrees and the twist multiplies
by the rank-one value. Thus the attached representation is
\[
                 \rho|_{G_{M'}}\otimes(\kappa')^a .
\]
\cite{clozel-rajan}\label{cor:ordinary-count-iii-admission-1}

Write \(R=\rho\otimes\operatorname{Ind}_{G_{K'}}^{G_\Q}
(\kappa')^a\), with the tensor basis already specified.
The tensor identity for induction gives
\[
  R|_{G_L}
   \simeq\operatorname{Ind}_{G_{M'}}^{G_L}
                 \bigl(\rho|_{G_{M'}}(\kappa')^a\bigr).
\]
The two constituents on \(G_{M'}\) are inequivalent by
the determinant test above. Hence the cuspidal representation
over \(M'\) is inequivalent to its conjugate over \(L\),
by unramified Galois comparison. Quadratic automorphic
induction is cuspidal under this hypothesis. At an embedding \(\tau:L\hookrightarrow\C\), if \(\tau_+,\tau_-\)
are the two extensions to \(M'\), the infinity multiset is
\(\Lambda_{\tau_+}\sqcup\Lambda_{\tau_-}\). If \(B_+,B_-\) denote the unshifted rank-two multisets at the two embeddings, the character of \(K'\) gives
\[
 \Lambda_{\tau_+}\sqcup\Lambda_{\tau_-}
 =(B_++am_0)\sqcup(B_--am_0)
\]
after choosing the positive embedding first.

Their
restrictions to \(K'\) are the two distinct embeddings,
so the parameters are the four separated by the choice
of \(a\). They are algebraic in rank four: the rank-two
algebraic shifts have the same half-integral parity as
the rank-four shifts, and the twists add integers.
Thus the induced representation is regular algebraic,
with \(R|_{G_L}\) as its Galois representation by comparison
at the unramified places. For the polarization, write \(W=\rho|_{G_{M'}}\otimes(\kappa')^a\) and let \(\iota\) denote complex conjugation. The rank-two factor extends over \(\Q\), its determinant is \(D\), and \((\kappa')^\iota=(\kappa')^{-1}\). The exterior pairing in rank two therefore gives
\[
 W^\iota\simeq \rho|_{G_{M'}}\otimes(\kappa')^{-a}
       \simeq W^\vee\otimes D|_{G_{M'}} .
\]
Induction and comparison at unramified places give the corresponding essential conjugate duality of the automorphic representation by strong multiplicity one.
\cite{clozel-rajan}\label{cor:ordinary-count-iii-admission-2}

The finite-place conditions can be checked before induction.
At primes in \(S\) and at \(w\), the extension \(M'/L\)
splits, by the splitting in \(K'\) at \(S\) and the
choice at \(w\). At primes away from \(S\cup\{w\}\), the
quadratic extension \(M'/L\) is unramified. In fact the
quadratic character there is the restriction of \(\bar\alpha\)
on \(G_L\), and
\[
 \alpha=\eta\epsilon_\ell/\omega
\]
is
unramified away from \(S\). The representation
\(\rho|_{G_L}\) is unramified at these places as well, since
the conductors of \(\phi\) and \(\theta\) are supported on
\(S\cup\{w\}\). The Hecke character of \(K'\) is
unramified away from \(\ell\), so the twisted base change
is spherical at the nonsplit places of \(M'/L\).
Here one uses the local compatibility of the rank-two
attachment away from \(\ell\), under restriction and
character twist.

At a split place of \(M'/L\) the
parameter of the induction is the direct sum from its
two places. At a nonsplit unramified place the inertia group \(I\) is the same for the quadratic extension and the base. On restricting the induced parameter to inertia, its two coset summands both restrict trivially:
\[
 \left.(\operatorname{Ind}\tau)\right|_I
 =\left.\tau\right|_I\oplus\left.\tau^c\right|_I
 =\mathbf{1}^{\oplus4}.
\]
The summands have zero monodromy, hence the induction is spherical there. These are
the split and unramified local comparisons for the induced representation.

In particular the induced rank-four representation matches
the tensor at all places away from \(\ell\), including
monodromy. At \(\ell\) the quadratic extension is split. The ordered
exponents of the tensor before the scalar shift and their gaps are
\[
      (k-1+m,k-1-m,m,-m),\qquad
      (2m,k-1-2m,2m).
\]
Every gap is at least two. The local Weil lines are unlinked
by the valuation calculation in
Section~\ref{sec:count-two-seed-specialization}.
The local--global compatibility theorem at \(\ell\)
applies to the regular algebraic polarized cuspidal
representation, since those gaps give Shin regularity.
\cite{blggt,blggt-lg}\label{cor:ordinary-count-iii-admission-3}

\end{proof}

\subsubsection{The eigenpoint at both ordinary levels}
\label{par:count-three-actual-seed}

The cuspidal tensor has to give an ordinary specialization at each of two levels. With the chosen local flags one can impose both sets of Hecke relations at the same integral point.

\begin{proposition}[Support of the integral tensor seed]\label{prop:count-three-ordinary-attachment}
The normalized integral tensor gives a seed evaluation on both
ordinary Hecke quotients. In the framed ring over \(F\), write \(f_i\) for their level maps. For these maps \(f_i\) and ideals \(J_i\), every prime \(P\) below the seed kernel in \(B\) satisfies 
\[
 f_1(J_1)B+f_2(J_2)B\subseteq P\subseteq\ker y_{\rm seed}.
\]
 This is the initial support condition for the four
residual links.
\end{proposition}

\begin{proof}
On the framed integral tensor the scalar is \(\delta\);
the twist over \(L\) by \(\delta^{-1}\) recovers
\(\delta^{-1}(\delta R|_{G_L})=R|_{G_L}\), the attached
cuspidal representation. Write \(e_0=j+2-k\)
for the exponent of \(\delta^2\) on an open subgroup.
There is a positive integer \(h\) with
\[
 (\delta^{-2}\epsilon_\ell^{e_0})^h=1,
\]
since this
character has finite image. Along with the even doubled
gaps and the parity of the doubled first exponent plus
the doubled correction, this is the scalar correction data
in the ordinary seed interface. The required irreducibility
on the cyclotomic kernel was proved above. The choice of
\(\delta\) gives determinant multiplier \(\vartheta\) on \(G_F\)
because \(\theta\) is trivial there.

For each link use the parallel weight shift of
Paragraph~\ref{par:count-two-seed-support}, which is
needed here as the seed's doubled weights can be odd.
Write \(\beta_i\) for the link companion and \(h_0\) for
the first doubled seed exponent. Take the shift
\(n=|h_0|\). All adjacent doubled gaps are even. Thus for every \(r\)
\[
 h_r\equiv h_0\equiv |h_0|=n\pmod2,\qquad (h_r+n)/2\in\Z .
\]
This gives integral seed exponents for \(\kappa\beta_i^n\) at the chosen embeddings. Thus on \(G_L\) the shifted
evaluation has the cuspidal attachment just constructed
twisted by the algebraic Hecke character of
\(\kappa\beta_i^n\delta\).

At a selected embedding, if \(m^{(R)}_r\) denotes the
\(r\)-th integral untwisted tensor exponent, the scalar exponent is
\[
             \frac{h_r+n}{2}-m^{(R)}_r\in\Z .
\]
It is independent of \(r\), since the companion acts by one
scalar on all four lines.
At the conjugate embeddings integrality follows from
the link pairing with its integral cyclotomic multiplier
and the rank-two determinant \(D\). The anticyclotomic
companion preserves the link pairing and its unitary
normalization.

For the tame types, after killing \(\theta\) on \(G_F\)
the inertia of \(R\) away from \(\ell\) has finite image of
order prime to \(\ell\). Where \(K'\) splits this follows
from the rank-two inertia calculation and unramifiedness
of \(\kappa'\). At the other places the quadratic extension
is unramified and the whole seed is unramified. The
finite inertia of \(\kappa\beta_i^n\delta\) away from \(\ell\) has
order prime to \(\ell\), by the common twist and the
formula for \(\delta\). If \(e\) is the order of this finite inertia image, each inertial root satisfies
\[
 z^e=1,\qquad e\in\mathbb Z_\ell^\times.
\]
Thus the polynomial \(X^e-1\) has simple roots modulo \(\ell\). Each inertial root \(z\) is consequently determined by its reduction:
\[
 z^e=1,\qquad (X^e-1)'(z)=e z^{e-1}\in O^\times
 \quad\Longrightarrow\quad z=[\bar z].
\]
Here the brackets denote the Teichmüller lift. This is the type at each prescribed
tame place, with zero monodromy.

At \(w\) it is unramified
after killing \(\theta\). At both exceptional places the
seed is unramified as required: these places are outside
\(S\cup\{w\}\), and their inertia is killed by the common
link twist. The finite places where it can ramify are
split over \(F\), including those over \(\ell\).
Above \(\ell\), the same exponent gaps and the ordered
local characters from the tensor flag give the ordinary
vector using the ordered translation in
Section~\ref{sec:count-two-seed-specialization}.
At inert places over \(F\) the attachment is spherical.
The compact unitary descent and the projection to
ordinary scalar functions therefore apply with exactly
the prescribed types and one exceptional place on each
of the two levels.
\cite{thorne-unitary-descent,newton-thorne-unitary}\label{cor:ordinary-count-iii-admission-4}

Let
\[
 f_{i,n}:R_i\to B
\]
be the shifted link map to the
common framed ring and \(y_{\rm seed}\) the tensor point.
The characteristic coefficients of \(y_{\rm seed}\circ f_{i,n}\)
agree with those of the ordinary eigenpoint at split good
primes, hence everywhere by density. The local-character
coordinates agree by the matched weights and uniformizer
translations. The spectral and local matches give
\[
 J_i\subseteq\ker(y_{\rm seed}\circ f_{i,n}).
\]

By the finite-twist invariance in
Paragraph~\ref{par:count-two-seed-support},
\[
                  f_{i,n}(J_i)B=f_i(J_i)B
                             \subseteq\ker y_{\rm seed}.
\]
This transfers support to the original link. The hypotheses of the adjoint comparison in the same
paragraph hold on this tensor: it is cuspidal, irreducible
on the cyclotomic kernel, has the separated ordinary
exponents, and is unramified at the exceptional place.
The comparison gives vanishing at every prime below
the seed evaluation kernel. In the normalized framed
ring over \(F\), the sum of the two extended ideals therefore
satisfies the seed condition for the four links.
Propositions~\ref{prop:tower-residual-propagation},
\ref{prop:count-three-residual-path} and
\ref{prop:count-three-extra-vertices} propagate this same ideal
to the normalized target in the rank-four construction.
\end{proof}


\subsection{The four residual links for the rank-four case}
\label{sec:count-three-chain}

The tensor now supplies a point in ordinary support. To reach the target we construct residual intermediate representations over a common finite field, so successive comparisons refer to the same local characters.

\begin{proposition}[The residual path over one finite field]\label{prop:count-three-residual-path}
Use the tower of Proposition~\ref{prop:count-three-fitted-layers},
the normalized tensor of Proposition~\ref{prop:count-three-tensor-endpoint}
and the support conditions of Lemma~\ref{hyp:count-three-path-endpoints}.
In the notation for the two bridge points \(P_3,P_4\), the common intermediate datum \(P_2\), target \(P_0\) and seed \(S_F\), the four residual lines carrying the ordered flags form the path 
\[
 P_0\longrightarrow P_3\longrightarrow P_4\longrightarrow P_2\longrightarrow S_F.
\]
 All lines and the good endpoints can be
defined over a single finite residue field after the indicated
coefficient extension.
\end{proposition}

\begin{proof}
Use the ordered frame on the tame field \(F_0\) from
Section~\ref{sec:count-three-tower}. Write \(d_i\) for
the four diagonal residual characters, and put
\(A=d_1/d_0\), \(B=d_2/d_0\), so \(d_3/d_0=AB\).
The roots of these three ratios are
\(-w_\alpha,-w_\Psi,-w_\vartheta\). Each of the
ratios is nontrivial at the normalizer. A residual
datum comprises its global matrices and one lower
unipotent symplectic frame \(N_v\) at each marked
place. The conjugated local matrices are upper
triangular with prescribed diagonal throughout
the chain.

\begin{established}[The support of the normalized target and the tensor seed]\label{hyp:count-three-path-endpoints}
On the fitted tower of Proposition~\ref{prop:count-three-fitted-layers}, write for off-diagonal entries:
\[
 \begin{split}
 S_+&=\{w_\beta,-w_\alpha,-w_\Psi,-w_\vartheta\},\\
 S_-&=\{-w_\beta,-w_\alpha,-w_\Psi,-w_\vartheta\},\\
 S_0&=S_+\cap S_- .
 \end{split}
\]
Write \(\mathcal A\) for the residual data with
global support in
\(\{-w_\beta,-w_\Psi,-w_\vartheta\}\) and
\[
 (N_v)_{10}\ne0
\]
at every marked place.
Here support refers to entries that are allowed to
be nonzero. The normalized rational target restricts
to a datum \(P_0\) with global support in \(S_+\) and frame support in
\(S_0\). Its two possible initial conditions are
\begin{enumerate}
\item the global roots \(-w_\Psi\) and \(w_\beta\) occur,
\item the global support lies in
\[
 S_+\setminus\{-w_\Psi\},
\]
all the frame
entries \((N_v)_{20}\) are nonzero and
the global root \(-w_\alpha\) occurs.
\end{enumerate}


The input from the tensor construction is one more
datum \(S\) on \(F_0\), with global support in \(S_-\).
It can be ramified at one extra split rational prime
\(w\) with \(w\equiv1\pmod\ell\). On every deeper
layer where that tame inertia has been killed, the
construction gives a characteristic-zero integral
seed with residual datum \(S_F\). This is the
restriction of \(S\) with its transported frames.
The transported seed satisfies
\[
      S_F\in\mathcal A,\qquad
        -w_\beta,\,-w_\vartheta\in\operatorname{supp}(S_F).
\]
\end{established}
\begin{proof}
Proposition~\ref{prop:rational-rank-four-torus} gives the two
alternatives for the marked rational target in its normalized
basis. Lemma~\ref{lem:count-two-target-restriction} gives the restriction
calculation for these shapes: a nonzero normalized cocycle of
character \(\theta\) remains nonzero on the normal subgroup
containing a marked element \(\tau\) with \(\theta(\tau)\ne1\).
Indeed its vanishing there would give
\((1-\theta(\tau))f(g)=0\) for every \(g\) by conjugation.
Proposition~\ref{prop:count-three-fitted-layers} retains the
distinguishing elements here, and the same lemma treats the
transported local frame in the second shape.
For the other endpoint use Proposition~\ref{prop:count-three-tensor-endpoint}.
In the tensor basis its two marked factor entries supply
\(-w_\beta,-w_\vartheta\) and the prescribed nonzero local frame
entry. Killing the added tame inertia gives precisely the deep-layer
seed there, with the same marked reductions. These two calculations
apply to the two endpoints before any interpolating family is chosen.
\end{proof}

\proofstep{Simultaneous cup lifts.}
For support \(S_-\), row one of a product gives
\[
 (r(g)r(h))_{10}=r(g)_{10}d_0(h)+d_1(g)r(h)_{10}.
\]
Dividing by \(d_0(gh)\) yields the \(A\)-cocycle \(x\).
For \(S_+\), row two uses just columns zero and two,
and the same computation gives the \(B\)-cocycle \(y\). Normalization makes
these cocycles zero at the normalizer.
The first-column equations on the marked
decomposition groups give the local coboundaries
\[
 x(s)=(1-A(s))(N_v)_{10},\qquad
 y(s)=(1-B(s))(N_v)_{20}
\]
in the respective supports. In the first identity
row one of \(r(s)N_v\) uses only the first two
entries of its first column. In the second identity
row two uses only entries zero and two.

Apply the locally split cup lift of
\cite{count-two-cohomology}, with \(x\)
from \(S\) held fixed and a new \(B\)-cocycle \(y\)
chosen nonzero. Its finite coefficient module has
coordinates \((y,z')\), with action
\[
   (u,v)\longmapsto (B(g)u,(AB)(g)v+x(g)B(g)u).
\]
The result is a cochain \(z'\) with product
cross term \(x(g)B(g)y(h)\).
The cochain \(z=xy-z'\) has cross term
\(y(g)A(g)x(h)\), by the product rule for \(xy\).
Write \(\mathcal D=\operatorname{diag}(d_i)\) and
\[
 U(x,y,t)=
 \begin{pmatrix}
 1&0&0&0\\ x&1&0&0\\ y&0&1&0\\ t&y&-x&1
 \end{pmatrix}.
\]
At a marked place write \(x=(1-A)m_v\) and
\(y=(1-B)n_v\). For this local primitive the change \(z=xy-z'\) gives
\[
  z'=(1-AB)o'_v-xBn_v
     \quad\Longrightarrow\quad
  z=(1-AB)(m_vn_v-o'_v)-yAm_v .
\]
Here \(x=(1-A)m_v\) and \(y=(1-B)n_v\) are used on the same
decomposition group.

Consequently the matrices \(U(x,y,2z-xy)\mathcal D\)
with frames \(U(m_v,n_v,m_vn_v-2o'_v)\)
give a datum \(P_2\) of support \(S_0\).
In particular \(P_2\) has the
same \((1,0)\) coordinate as \(S\) in both global
matrices and local frames, and its root
\(-w_\Psi\) occurs. The new ramification is only
at finitely many extra split rational primes as
in the cup lift.

Next retain the \(B\)-cocycles of \(P_0\) and
\(P_2\) simultaneously. Use the direct sum of
the two cup coordinates with a common \(A\)-coordinate.
The finite action is
\[
 (u,v_0,v_2)\longmapsto
  (A(g)u,(AB)(g)v_0+y_0(g)A(g)u,
           (AB)(g)v_2+y_2(g)A(g)u).
\]
A cocycle \((x,v_0,v_2)\) for this action has the usual
\(A\)-cocycle law in its first coordinate and satisfies
\[
 v_i(gh)-v_i(g)-(AB)(g)v_i(h)=y_i(g)A(g)x(h),
                                  \qquad i=0,2 .
\]
These are the two cup equations with one common \(x\).
Prescribe the ramification direction \((1,0,0)\).
The duality argument gives a solution with \(x\ne0\)
and both lifts locally split at the old places. Hence it constructs two data \(P_3,P_4\)
with global and frame support in \(S_0\), using the
retained \((2,0)\) entries of \(P_0,P_2\),
respectively. They have the same \((1,0)\) entry in
each matrix and frame. Their common root
\(-w_\alpha\) occurs. The same local constants
are used for the common \(x\) at each marked place.

For every new prime, including \(w\), let \(C_w\) be
the degree-\(\ell\) cyclotomic subfield at that conductor
and set \(F_1=F_0\prod_w C_w\). The diagonal on each added
inertia group is trivial. For that inertia group \(I_w\), the cocycle factors as
\[
 I_w\longrightarrow I_w^{\rm tame,(\ell)}
 \simeq\Z_\ell(1)\longrightarrow\mathbf F_\ell(1)
 \longrightarrow \text{the residual coefficient module}.
\]
The degree-\(\ell\) cyclotomic field at \(w\) is totally ramified there, so its inertia subgroup lies in the kernel of the middle reduction.

On inertia at
\(w\) the tensor input will use that quotient as
well. Thus restriction removes the extra ramification
of all data at once. Over this tame \(F_1\), fix the
cohomological bound \(C\) for the growing cyclotomic
layers, and choose the layer \(F\) of the indicated
depth. The good target remains the restriction of
the same rational point. The identities of retained
entries and the nonvanishing needed for the links
persist on restriction with the transported flags.
Here the identities of global entries are
restrictions, and the nonzero root cocycles survive
by the \(\ell\)-power calculation of
\cite{count-two-cohomology}.

For the local frames one uses the identities
\[
 (N_v')_{i0} =
 \frac{(r(g)N_{v_0})_{i0}}{(r(g)N_{v_0})_{00}}
 \qquad (i=1,2)
\]
whenever the indicated conjugating matrix has unit
first pivot. This follows from the first column in
the lower--upper factorization defining the
transported frame. Each support is triangular in
a suitable coordinate order, so the principal pivots
are units and elimination keeps the same support. In both supports the denominator
equals \(d_0(g)\). For \(i=2\) with support \(S_+\),
row two uses only columns zero and two, and for
\(i=1\) with support \(S_-\), row one uses only
columns zero and one. This proves that the equalities
of retained frame coordinates survive transport.
The formula for \(i=1\) also transports the
nonzero seed frame coordinate to \(P_2\).
The rational lattice shape already supplies the
marked frame on the restricted target, so the
comparison uses its actual local flag.

\proofstep{The links and their endpoints.}
The path over \(F\) is
\[
                     P_0\longrightarrow P_3
                  \longrightarrow P_4
                  \longrightarrow P_2
                  \longrightarrow S_F .
\]
For the first and third links the two endpoints
have global and frame support in \(S_+\) and
the same \((2,0)\) data. Order the blocks as
\(\{0,2\},\{1,3\}\). The matrices are block lower
triangular in this order. Their first diagonal
block is fixed because it consists of the two
diagonal entries and the common \((2,0)\)
entry. In bases dual for the pairing, write \(A,D\) for the two diagonal blocks. The similitude equation in the paired blocks is
\[
 A^{\mathsf T}D=\mu I,\qquad D=\mu A^{-\mathsf T}.
\]
Thus fixing \(A\) and the similitude fixes \(D\). For the second
and fourth links the global support is \(S_-\)
and the frames are lower triangular.
Use blocks \(\{0,1\},\{2,3\}\), whose diagonal
blocks are fixed by the common \((1,0)\) data and
the similitude.

Thus at each link, for both the
representation and the frames, the difference
lies in the lower-left block ideal of a block
triangular algebra. This ideal is square-zero
and is preserved by the symplectic adjoint in
each order because the two blocks are paired.
In that algebra let \(A_s=(1-s)A_0+sA_1\) and
\(B_s=(1-s)B_0+sB_1\) be matrices built from two endpoint
pairs. Their differences have zero product, giving
\[
 A_s B_s
      =(1-s)A_0B_0+sA_1B_1,
 \qquad
 (A_s)^\dagger=(1-s)A_0^\dagger+sA_1^\dagger .
\]
For each link write its global matrices \(r_s\) and local
frames \(N_{v,s}\) by affine interpolation. The two product
identities yield
\[
 r_s(g)r_s(h)=r_s(gh),\qquad
 r_s(g)^\dagger r_s(g)=\nu(g)I,\qquad
 r_s(h)N_{v,s}=N_{v,s}B_{v,s}.
\]
In the last identity \(h\) is local and \(B_{v,s}\) is the
interpolation of the conjugated upper triangular matrices.
It retains the same ordered diagonal. In particular the normalization, prescribed
unramified quotients and residual diagonal hold
on all four polynomial lines.

The first bridge point \(P_3\) has the same
\((2,0)\) data as \(P_0\) and a nonzero
\(-w_\alpha\) root. Thus it has the first type
of ordinary good residual datum if the target
has global \(-w_\Psi\), and the frame type if
the target has only the local \((2,0)\) data.
For \(P_4\) the two nonzero roots are explicit:
\[
 \{-w_\alpha,-w_\Psi\}\subseteq\operatorname{supp}_{\rm glob}(P_4).
\]
This verifies the first ordinary test using the global matrices alone.

Write \(r_2\) and \(N_{2,v}\) for the global matrices and marked frames of \(P_2\).
The \((1,0)\) entry agrees with the seed. In support \(S_0\) the symplectic adjoint pairs it with \((3,2)\), while the local frames retain their nonzero seed entry. Thus
\[
 (r_2(g))_{10}=(r_2(g))_{32}=0\quad\text{for every }g,
 \qquad (N_{2,v})_{10}\ne0\quad\text{at every marked }v.
\]
It follows that \(P_2\in\mathcal A\); its root \(-w_\Psi\) occurs. The endpoint
\(S_F\) is in \(\mathcal A\) with both
\(-w_\beta,-w_\vartheta\) occurring. The last
two good-point conditions are the extra cases
required for the rank-four case.
\label{cor:ordinary-count-iii-chain-1}

\begin{construction}
\(\displaystyle
P_0 \xrightarrow[\ (2,0)\ ]{S_+} P_3
    \xrightarrow[\ (1,0)\ ]{S_-} P_4
    \xrightarrow[\ (2,0)\ ]{S_+} P_2
    \xrightarrow[\ (1,0)\ ]{S_-} S_F
\)
\caption{The four affine residual lines. Above each arrow is the allowed global support on that line, below is the entry shared by the two endpoints in both the global and frame data. Fixing the paired diagonal blocks makes the interpolation multiplicative.}\label{fig:four-residual-line-supports}
\end{construction}

\proofstep{Descent of the residual path.}

The links have been verified after making coefficient choices for their cocycles. They enter the component argument only once all four descend to the same finite residue field.

The matrices in a polynomial line are locally
constant over the discrete ring \(\overline{\mathbf F}_\ell[X]\).
Each continuous polynomial-valued matrix has finite image by
compactness and discreteness. Let \(S\) be the union of all coefficients of those images and frames across the four lines. Both the set of matrices and their polynomial degrees are finite, so
\[
 k'=\mathbf F_\ell(S),\qquad [k':\mathbf F_\ell]<\infty .
\]
All four matrix families have entries in \(k'[X]\).
The algebraic equations, inertia equations and
nonvanishing conditions can be read there
because scalar extension is injective. The four
lines and their good endpoints consequently
descend to a single finite residue field. This
extension can be made together with the
coefficient extensions for the actual
characteristic-zero seed, the small target model
and the two ordinary levels.
\label{cor:ordinary-count-iii-chain-2}
\end{proof}

\subsection{Component propagation at the extra vertices in the rank-four construction}
\label{sec:count-three-components}

Having the required residual links describes a path of points. At its extra vertices one still needs to pass between components without losing the condition that both ordinary ideals vanish.

\begin{proposition}[Propagation through the extra vertices]\label{prop:count-three-extra-vertices}
At the extra vertices of the residual path in
Proposition~\ref{prop:count-three-residual-path}, use the
depth and inertia-coincidence dimension bounds of
Proposition~\ref{prop:tower-component-propagation}. If \(P,P'\) are any two component primes through such a vertex and \(J\) denotes the ordinary ideal there, 
\[
 J\subseteq P\quad\Longrightarrow\quad J\subseteq P'.
\]
 Along the
four links this gives the backward propagation from the integral
seed of Proposition~\ref{prop:count-three-ordinary-attachment}.
\end{proposition}

\begin{proof}
For the first three vertices of
Section~\ref{sec:count-three-chain}, the ordinary
component argument of Paragraph~\ref{par:count-two-propagation}
already applies. We prove the additional statement
for a residual datum in \(\mathcal A\) having global root
\(-w_\Psi\), or both \(-w_\beta,-w_\vartheta\).
The ideal in the argument is the same sum of the two
ordinary Hecke ideals.

\proofstep{The reducible loci at an extra vertex.}
Use the nine supports in the component argument.
The ninth is
\[
                   Y=\{\pm w_\beta,\pm w_\vartheta\}.
\]
For every other support meeting the extra vertex, the
tangent dimension is at most \(d+20C\).
Indeed, use the symplectic cocharacter with exponents
\((0,0,1,1)\). It sends the global matrices to the
diagonal and the frame at \(v\) to
\[
                     N_v=1+b_v(E_{10}-E_{32}),
                     \qquad b_v\ne0.
\]
The tangent dimension can only increase at this
limit. The first-column flag equation at a datum in
\(\mathcal A\) gives \(d_0=d_1\) on each marked
decomposition group, since global row one has no
off-diagonal entries and \((N_v)_{10}\ne0\).
The same equality holds at the limit.

For a global tangent cocycle \(U\) and a frame
variation \(V_v\), put
\(Z(s)=U(s)+\mathcal D(s)V_v\mathcal D(s)^{-1}-V_v\)
at the local group. Its conjugate by \(N_v\) is
upper triangular. Direct multiplication gives
\[
 Z_{10}=b_v^2 Z_{01}+b_v(Z_{00}-Z_{11}),\qquad
 Z_{21}=0,\qquad Z_{30}=0 .
\]
For the last equality use also entries \((2,0),(3,1)\)
of the upper triangular condition.
First fix the two independent global diagonal
cocycles. This costs at most \(2C\).
After fixing the two diagonal cocycles, the local equations give
\[
       U_{10}=b_v^2U_{01},\qquad [U_{21}]=[U_{30}]=0
         \text{ in the local root cohomology}.
\]
The root character for \(U_{10},U_{01}\) is trivial here,
so their frame terms vanish.

For the eight supports preceding \(Y\), the ones
that can meet these vertices are rows \(1,4,5,7,8\)
in the order of the displayed list.
Fix the global cocycle indicated for each surviving row.
\[
 \begin{array}{c|ccccc}
 \text{row}&1&4&5&7&8\\\hline
 \text{entry}&(1,2)&(0,3)&(0,1)&(1,2)&(0,3)
 \end{array}
\]
The image dimension of this coordinate map costs at most \(d+C\).
In the resulting kernel the two upper long-root
entries \((0,3),(1,2)\) vanish, either by this choice
or by the support.

In each row either \((0,1)\)
or \((1,0)\) also vanishes globally. The local
relation and \(b_v\ne0\) force both to vanish
locally. Thus every odd root entry is locally
a coboundary. The remaining global entries
have dimension at most \(10C\) by the strict
and even cohomology bounds. At most \(6g\)
frame coordinates remain, where \(g\le C\)
is the number of places over \(\ell\). Including
the three fixed cocycles, the tangent bound is
\[
                   2C+(d+C)+10C+6g\le d+20C .
\]
The completed support quotient therefore has
dimension at most \(d+20C+1\), with the extra
parameter from the coefficient DVR.

\proofstep{The support \(Y\).}
A representation supported in \(Y\) preserves
the two planes with indices \(\{0,3\}\) and
\(\{1,2\}\). Its two determinants are both the
similitude \(\nu\).
At any local point where the first lower flag
entry \((N_v)_{10}\) is nonzero, the ordinary first
vector has nonzero projection on each plane.

Let \(u_1,u_2\) be the two nonzero projections of
the first flag vector. Both satisfy
\(r(h)u_i=\psi_{v,0}(h)u_i\) on the decomposition group,
since the two plane projections commute with \(r(h)\).
The determinant \(\nu\) of each plane determines its other
graded character. Thus the full composition multiset is
\[
       \{\psi_{v,0},\psi_{v,0},\nu/\psi_{v,0},\nu/\psi_{v,0}\}.
\]
Since the pairing gives \(\psi_{v,3}=\nu/\psi_{v,0}\),
the multiset forces on the same decomposition group the coincidence
\[
               \psi_{v,1}=\psi_{v,0}\quad\text{or}\quad
               \psi_{v,1}=\psi_{v,3}.
\] At every prime of the completion at a
vertex in \(\mathcal A\) the frame entry stays
nonzero, since it is a unit at the vertex.

The dimension comparison with the characteristic
polynomial and ordinary character ring also holds
on the part of \(Y\) outside the first eight
supports. On this part each plane has both of
its off-diagonal entries nonzero somewhere.
Over the fraction field, the four eigenvalues
of the normalizer are algebraic over the base.
If \(\lambda_i\) are the marked eigenvalues, the four traces are
\[
    \operatorname{tr}(r(g_0)^j r(g))
            =\sum_i\lambda_i^j r(g)_{ii}\qquad(0\le j<4).
\]
Invert the Vandermonde matrix to read the diagonal entries. The same holds for
the diagonal entries of \(r(gh)\).
Subtracting products of the known diagonals
therefore determines
\[
                  r(g)_{ij}r(h)_{ji}
\]
for each opposite pair of allowed off-diagonal
positions and all \(g,h\).

Choose one nonzero
value in one direction for each plane. All other
entries in both directions in that plane are
then obtained from these products by division.
Thus two field generators, together with the
four entries determining each local lower
symplectic frame, suffice over an algebraic
extension of the base field. The finite type
dimension comparison increases the dimension
by at most \(2+4g\), just as at the absolutely
irreducible primes.

\proofstep{The selected arc.}

An arc avoiding the excluded loci lets us apply the one-point support result.

At an intersection of adjacent components through
an extra vertex, start on the modular component
and cut by the coefficient uniformizer. Choose
a prime in the intersection with quotient
dimension at least \(N=2d+1-2g\).
The depth inequalities give
\[
     d+20C+1<N,\qquad 2d+1-d_v+2+4g<N .
\]
If \(Z\) is the chosen integral closed subset, its
dimension is at least \(N\). The first inequality gives
\(Z\not\subseteq\) any of the eight small reducible supports.
The second gives the same exclusion for inertia coincidence loci,
using the dimension comparison on \(Y\).
In particular the prime cannot be supported
in \(Y\): its unit lower frame entry would
force a coincidence on the full decomposition
group and hence on the open inertia.
The prime therefore has absolutely irreducible
generic representation and separated inertia
characters.

In the local coordinate domain of \(Z\), choose \(0\ne f_i\) from each forbidden-locus ideal and \(0\ne t\) from the maximal ideal. The element
\[
 f=t\prod_i f_i\ne0
\]
belongs to the maximal ideal and vanishes on each forbidden locus. Lemma~\ref{lem:local-curve-avoidance} gives a one-dimensional quotient retaining that product. Normalizing as in the component argument gives the centered formal arc with generic point outside the exclusions. Local vanishing at this
arc puts the shared Hecke ideal in the adjacent
component. This proves the required implication
from one modular component to all modular
components at both new kinds of vertices.
It applies to every vertex in the residual
chain, so the backwards argument of
Paragraph~\ref{par:count-two-propagation} applies
once the integral seed lies in the ordinary
support with its cuspidal attachment.
\label{cor:ordinary-count-iii-components-1}
\end{proof}


\subsection{The target in the rank-four construction at the rational level}
\label{sec:count-three-target}

Ordinary support has reached the target at the two chosen levels. Its Frobenius identities now allow the full-block exclusion and the rational conductor bounds needed to count the remaining representations.

\begin{proposition}[Rank-four bound at the original level]\label{prop:count-three-host-output}
For the data of Lemma~\ref{hyp:count-three-host}, with \(j\ge2\), and
using the four-link support of Proposition~\ref{prop:count-three-extra-vertices},
the characteristic-zero full blocks are excluded by the determinant test
of Proposition~\ref{prop:count-two-full-block-excluded}. Irreducible
targets have the rational flag, quadratic packet attachment and
level \(M_{\mathrm{host}}=\ell\,2^{n_{\rm III}}\mathcal Np_1\), with
\(n_{\rm III}\) defined in \eqref{eq:ordinary-dyadic-exponents}. Their normalized cohomology length satisfies

\[
 \log\ell\,s_\ell(\xi,j)\le(\ell\mathcal N j)^C
\]
 with the
constant independent of the fitted fields, exceptional places
and coefficient extension.
\end{proposition}

\begin{proof}
\proofstep{Classical control and full-block exclusion.}

After the four-link propagation the actual target kills
the sum of the two ordinary ideals. On the chosen CM
layer it has the target exponents \(j+1,j,1,0\), plus
the common link exponent. The common twist cancels in successive differences, giving
\[
 (j+1)-j=1,\qquad j-1\ge1,\qquad 1-0=1
\]
for \(j\ge2\). Subtracting the index from
the oppositely signed Hodge numbers gives a dominant
weight, exactly as in Proposition~\ref{prop:two-link-classical-control}.
The ordinary control argument there puts each evaluation
in the classical part of the same quotient. At the
exceptional prime of each link the residual matrix is
scalar, whereas the residual cyclotomic character has
nonidentity value. The rank-four tower supplies Frobenius elements with these two
residual properties at primes chosen separately from the seed prime. The Speh
argument of Proposition~\ref{prop:count-two-full-block-excluded} therefore
applies. The Frobenius modulus ratios at good primes are
strictly less than the residue norm.

For a full block on the
rational chart, first apply Proposition~\ref{prop:full-block-extension}
and run the propagation on the resulting nonsplit point.
The same ratio bound holds for the split point since the
characteristic polynomials agree. Adding the inertial exponents within each block gives its determinant exponent,
\[
 (j+1)+1=j+2,\qquad j+0=j.
\]
At a split good prime the ratio of determinant moduli is
\[
 p^{j+2}/p^j=p^2.
\]
For roots \(\lambda_i,\mu_i\) of the
respective blocks the two-link bound instead gives
\[
              \frac{|\lambda_1\lambda_2|}{|\mu_1\mu_2|}
                   =|\lambda_1/\mu_1|\,|\lambda_2/\mu_2|<p^2 .
\]
This is the required contradiction. The degree-zero and length comparisons of
Section~\ref{sec:rank-four-deformation} use this exclusion
on the chart for the rank-four construction.

\proofstep{Glue the two cuspidal attachments.}

At an irreducible target the four restricted residual
labels are distinct at the powered normalizer, so the
trace self-twist argument of
Proposition~\ref{prop:count-two-cuspidal-attachment} gives irreducibility
on \(G_L\). The two cuspidal transfers have equal polynomials
on the common good primes. Strong multiplicity one identifies
them. At each exceptional pair use the other level's local
comparison, and at \(\ell\) retain the comparison on the
marked side. Remove the common scalar twist to obtain the
cuspidal attachment of \(r|_{G_L}\).
\cite[Theorem 1.2 and Corollary 1.3]{newton-thorne-unitary}

\proofstep{The rational local ratio test.}
De Rhamness at \(\ell\) follows by descent from \(L\).
On the ordinary flag over \(\Q_\ell\), the four characters and their pairwise ratios are
\[
 \kappa_i=\nu_i\epsilon_\ell^{m_i},\qquad
 \frac{\kappa_i}{\kappa_h}=\frac{\nu_i}{\nu_h}\epsilon_\ell^{m_i-m_h},
 \qquad m=(j+1,j,1,0),
\]
with finite inertia for the \(\nu_i\). At an arithmetic Frobenius lift \(s\) with
\(\epsilon_\ell(s)=1\), the ordered residual labels
of the rational flag are
\[
                 (\bar\eta\bar\Psi,\bar\Psi,\bar\eta,1)(s).
\]
The adjacent characteristic ratios at \(s\) are
\[
 \bar\kappa_0/\bar\kappa_1=\bar\eta,\qquad
 \bar\kappa_1/\bar\kappa_2=\bar\Psi/\bar\eta,\qquad
 \bar\kappa_2/\bar\kappa_3=\bar\eta .
\]
The rational Frobenius test of Paragraph~\ref{par:dyadic-residual}
makes each of them nontrivial there.
This includes the middle pair when \(j=2\).
The Weil lines on the flag have parameters
\(\iota(\nu_i(g))\,\ell^{m_i\deg(g)}\) on the rational
Weil group.

Write \(\mu_i(g)=\iota(\nu_i(g))\ell^{m_i\deg(g)}\).
If two lines were related by the norm character, take the
order in which \(\mu_i(g)/\mu_h(g)=\ell^{\deg(g)}\).
Applying the coefficient-field \(\ell\)-valuation at a
Frobenius lift gives
\[
       m_i-m_h=1,\qquad
       \nu_i(g)/\nu_h(g)=1\quad\text{on the Weil group}.
\]
In the first computation the \(\nu_i\) are units. Weil-group
density then identifies the two finite-inertia factors. The residual values at \(s\) rule this out
for every pair with exponent difference one.

Indeed, on two such Weil lines of characters \(\mu_i,\mu_h\) the monodromy equation gives
\[
 (\mu_i(g)/\mu_h(g)-|g|)N_{ih}=0\qquad(g\text{ in the Weil group}),
\]
with the inverse norm in the opposite convention. The preceding exclusion in both orders forces every \(N_{ih}=0\).
Thus \(N=0\) rationally, and \(N|_{G_{L_v}}=0\) after restriction.
This use of the rational flag precedes restriction:
the quadratic root \(\bar\alpha\) is trivial at the
selected decomposition groups above \(\ell\) in \(L\),
where the rational separation need not hold.
\cite{colmez-fontaine}\label{cor:ordinary-count-iii-point-output-1}

At the selected places above \(\ell\) in \(L\), the
cuspidal transfer is tempered after undoing the common
algebraic normalization. As its parameter has zero
monodromy, the absolute values of all the ratios of
the Weil lines there are one. Scalar twisting cancels
in a ratio. Write \(f_{ih}\) for a ratio of two rational Weil lines.
For a rational Weil element \(g\), choose \(a>0\) with
\(g^a\) in the finite-index local restriction from \(L\).
Then
\[
                         |f_{ih}(g)|^a=|f_{ih}(g^a)|=1 .
\]
This proves the modulus equality on the rational group.
At the rational lift \(s\), \(\epsilon_\ell(s)=1\).
If the ratio of the finite-inertia factors were torsion at that
element, its complex absolute value would be one. This would give
\[
       |\mu_i(s)/\mu_h(s)|=\ell^{(m_i-m_h)\deg(s)}\ne1,
\]
contrary to the modulus equality.
Thus all these ratios have infinite order, and the same
is true on a positive power in any prescribed finite
extension's decomposition group.

\proofstep{Rigidity, descent and the valuation.}

Classicality supplies Frobenius data on the fitted field. Bounding the original character component requires counting the descended representations and their integral congruences.

The normalizer in \(G_L\) has trivial mod-\(\ell\)
cyclotomic character and retains four distinct labels.
Suppose a sign \(\chi\) through cyclotomic reduction
twists \(\bar\kappa_i\) to \(\bar\kappa_{\sigma(i)}\).
Writing \(g_0\) for this normalizer gives
\[
 \bar\kappa_{\sigma(i)}(g_0)=\chi(g_0)\bar\kappa_i(g_0)
 =\bar\kappa_i(g_0)\quad\Longrightarrow\quad\sigma(i)=i .
\]
The identity \(\chi\bar\kappa_i=\bar\kappa_i\) now holds as characters, so cancellation gives \(\chi=1\). The simple-spectrum
element just constructed and the argument of
Proposition~\ref{prop:marked-invariant-rigidity} give irreducibility on
\(G_{L(\mu_{\ell^\infty})}\).
The same adjoint vanishing theorem and local tangent
argument apply to the ordinary cuspidal attachment.
In particular the infinite-order ratios apply on the
open local subgroups needed for the weight-gap-one
extensions. The rational tangent class then vanishes
in the normalized chart.

\proofstep{Return to the original level.}

For the original field level take
\[
 M_{\mathrm{host}}=\ell\,2^{n_{\rm III}}\mathcal Np_1.
\]

Away from \(\ell\) the rational chart factors inertia
through the residual character quotient because
\(S_b=\varnothing\). It gives finite abelian inertia
with conductors dividing \(M_{\mathrm{host}}\).
At \(\ell\) the finite parts of the ordered rational
inertia characters have conductors dividing the same
level, and the rational ratio test supplies the zero
monodromy condition for this order of weights.
For quadratic descent choose the ramified prime of
\(K\) in the fitted tower outside the fixed packet
discriminants as in Proposition~\ref{prop:count-two-quadratic-packets}.
This is allowed in the simultaneous splitting
conditions of Section~\ref{sec:count-three-tower}.

At inertia at this ramified prime the original representation acts trivially and the quadratic sign acts nontrivially. The incompatible characteristic polynomials
\[
 (X-1)^4\quad\text{and}\quad (X+1)^4
\]
give the self-twist exclusion for each admissible quadratic restriction as in Proposition~\ref{prop:count-two-quadratic-packets}. Cyclic descent
and unramified comparison give the rationally
normalized cuspidal attachment on each member of the
packet list for which the restriction is irreducible.
At \(w\mid p\) of a member \(K'\) of the list, a
restricted unit character is the rational one composed
with \(N_{K'_w/\Q_p}\). For \(u\in1+M_{\mathrm{host}}\mathcal O_{K'_w}\), use the
matrix of multiplication in an integral local basis:
\[
 N_{K'_w/\Q_p}(u)=\det(I+M_{\mathrm{host}}A)
                         \in1+M_{\mathrm{host}}\Z_p .
\]
This gives triviality of the inducing unit characters on
the required congruence group. The weights on that side
are \(j+1,j,1,0\).
\cite{clozel-rajan,blggt,blggt-lg,thorne-adjoint}\label{cor:ordinary-count-iii-point-output-2}

\proofstep{Apply the pointwise valuation bound.}

The indexed trace product of the pointwise bound
therefore applies to the rank-four case with this \(M_{\mathrm{host}}\).
Choose \(p\geq5\) prime to \(2\ell\mathcal Np_1\),
with \(p\leq6\ell\mathcal Np_1+1\).
The reducible target and all other points have the
same trace comparison as in the even-order case, with the
\(2^{n_{\rm III}}\)-th roots included.
Each irreducible point maps to a distinct record in the fixed
packet product \(D_p\), with the two possible rational descents
retained as separate labels.

The cardinality bound of
Proposition~\ref{prop:marked-rational-packets} is polynomial in
\(M_{\mathrm{host}}(j+1)\).
The norm, integrality and length arguments then give
\[
      \log\ell\,s_\ell(\xi,j)\leq(\ell\mathcal N j)^C .
\]
Here \(p_1\) is fixed by the primitive target, and
\(2^{n_{\rm III}}\) and the order of the root group
in the trace are bounded by fixed powers of
\(\ell\mathcal N\). Thus \(C\) is independent of the
fitted fields, the weight of the Eisenstein seed, its
prime \(w\), both exceptional places and the degree of
the coefficient extension. This gives the rank-four positive-component valuation estimate with the same host and
parameters as the common-window calculation.
\end{proof}

